\SourcePage{240}
\chapter{量子上同调}\label{chap:quantum-cohomology}
\index{量子上同调|hyperpage}
本章将看到，量子上同调实际上有两种，即所谓\emph{大量子上同调环}与\emph{小量子上同调环}。
我们还将研究 \emph{Gromov--Witten 势}，它可视为树级 Gromov--Witten 不变量的指数型生成函数。
Gromov--Witten 势满足的 WDVV 方程，自然导出大量子上同调环的结合律，
并对\hyperref[chap:gromov-witten]{第~\ref*{chap:gromov-witten} 章}研究的例子产生一些有趣的推论。
另一个重要主题是 Dubrovin 形式体系，它说明量子上同调如何在上同调环上给出 Hodge 结构变分。

量子上同调是理解镜像对称的数学版本所需的最后一项要素。
如本章将看到，它给出三点函数的清晰表述。
此外，它导出上一段提到的 Hodge 结构变分，这是 Kähler 模上的变分。
镜像对称的一个关键思想是：经由镜像映射，这个 Hodge 结构变分对应于
镜像族的复模空间上的几何 Hodge 结构变分。
本章末尾定义数学上的镜像对并讨论镜像猜想时，将详细解释这一点。

从历史上看，量子上同调的思想最初在物理文献中以 Calabi--Yau 三维簇的\emph{手征环}的形式出现。
手征环被理解为通常上同调环的形变，见 \SourceCite{LVW}，
在大半径极限点与通常上同调环一致。
量子乘积也可在\IndexLocation{indexloc:1336:240}拓扑量子场论中作更一般的理解，见 \SourceCite{Witten1,Witten2}；\hyperref[sec:topological-qft]{附录 B.6} 将讨论这一点。
量子上同调最初的数学构造适用于半正辛流形，见 \SourceCite{RT1,MS}。

本章采用的量子上同调方法几乎完全是代数几何的。
介绍性的论述见 \SourceCite{Aluffi,FP,MS,Morrison7,Voisin3}。

\section{小量子上同调}\label{sec:small-quantum-cohomology}
\index{量子上同调!小量子上同调|hyperpage}
光滑射影簇 $X$ 的小量子上同调环，用\hyperref[chap:gromov-witten]{第~\ref*{chap:gromov-witten} 章}的 Gromov--Witten 不变量 $\langle I_{0,3,\beta}\rangle$ 定义。

\subsection{定义与性质}\label{subsec:small-quantum-definition}
取 $H^*(X,\mathbb Q)$ 的齐次基 $T_0=1,T_1,\ldots,T_m$，令
$g_{ij}=\int_XT_i\cup T_j$，则 $(g^{ij})=(g_{ij})^{-1}$ 为其逆矩阵。
回忆\hyperref[chap:moduli]{第~\ref*{chap:moduli} 章}，若 $\omega\in H^2(X,\mathbb C)$ 的虚部为 Kähler 类，则 $\omega$ 是一个复化 Kähler 类。
\SourcePage{241}
\begin{definition}\label{def:small-quantum-product}
设 $\omega$ 是光滑射影簇 $X$ 上的复化 Kähler 类。
对 $a,b\in H^*(X,\mathbb C)$，定义
\[
a*b=\sum_{i,j}\sum_{\beta\in H_2(X,\mathbb Z)}
\langle I_{0,3,\beta}\rangle(a,b,T_i)g^{ij}q^\beta T_j,
\qquad q^\beta=e^{2\pi i\int_\beta\omega}.
\]
称 $a*b$ 为 $a,b$ 的\emph{小量子乘积}。
\end{definition}
\index{量子乘积!小量子乘积|hyperpage}
按此写法，这一定义确实不够精确，因为对 $\beta$ 的求和可能无限，因而需要收敛性结果。
\hyperref[subsec:formal-small-quantum]{第~\ref*{subsec:formal-small-quantum} 小节}将说明几种赋予
定义~\ref{def:small-quantum-product} 严格意义的方法：以适当的形式表达式代替 $q^\beta$。
事实上，这才是通常的做法。不过，我们仍愿先给出上面这个不完整的定义，
因为它明确显示小量子乘积依赖 Kähler 模，而这是量子上同调的核心思想之一。
请读者暂且“悬置疑虑”，直到该小节；届时会看到，本节与下一节的所有断言，
经适当解释后，都是完全严格的。

采用物理学中“升指标”的记号，定义
\[
T^i=\sum_jg^{ij}T_j,
\]
于是 $T^0,\ldots,T^m$ 构成 $T_0,\ldots,T_m$ 的对偶基，满足
$\int_XT^i\cup T_j=\delta^i_j$。这样，小量子乘积可更简洁地写为
\begin{equation}\label{eq:small-quantum-product}
a*b=\sum_i\sum_{\beta\in H_2(X,\mathbb Z)}\langle I_{0,3,\beta}\rangle(a,b,T_i)q^\beta T^i.
\end{equation}
本章与下一章都使用此记号。

\hyperref[chap:gromov-witten]{第~\ref*{chap:gromov-witten} 章}为有理上同调类定义了 Gromov--Witten 不变量。
在上面的 $*$ 公式中，我们通过线性性将定义扩张到复系数，
使 $\langle I_{0,3,\beta}\rangle(a,b,T_i)$ 对 $a,b\in H^*(X,\mathbb C)$ 有意义。
还要注意，对 $\beta\in H_2(X,\mathbb Z)$ 与齐次 $a,b$，
次数条件~\eqref{eq:gw-top-degree-condition} 表明定义~\ref{def:small-quantum-product} 中的
Gromov--Witten 不变量 $\langle I_{0,3,\beta}\rangle(a,b,T_i)$ 为零，除非
\begin{equation}\label{eq:quantum-three-point-degree}
\deg a+\deg b+\deg T_i=2\dim X-2\int_\beta\omega_X,
\end{equation}
其中 $\omega_X$ 是 $X$ 的典范类。因此，上述定义只对这些 $\beta$ 求和。
也可检验，二元运算 $*$ 不依赖所选的基 $T_i$。

对于 Calabi--Yau 流形，有如下猜想。
\begin{conjecture}\label{conj:small-quantum-convergence}
对于 Calabi--Yau 流形，只要 $\omega$ 的虚部足够大，定义~\ref{def:small-quantum-product} 中的和便收敛。
\end{conjecture}
暂且假设此猜想成立。如前所述，\hyperref[subsec:formal-small-quantum]{第~\ref*{subsec:formal-small-quantum} 小节}
将通过重新解释 $q^\beta$，说明如何避免收敛问题。
幸运的是，对一些射影流形，已知上述求和性质良好，
因而能在 Kähler 模的意义下严格定义小量子乘积。
\begin{proposition}\label{prop:quantum-finite-fano}
设 $X$ 是光滑射影簇，满足以下条件之一：
\begin{enumerate}[label=(\roman*)]
\item $X$ 是 Fano 簇；
\item $X$ 上有半单 Lie 群的可迁作用。
\end{enumerate}
则对任意 $a,b\in H^*(X,\mathbb C)$，定义~\ref{def:small-quantum-product} 中的和为有限和。
\end{proposition}
\SourcePage{242}
\begin{proof}
需证明：$X$ 上满足式~\eqref{eq:quantum-three-point-degree} 且可由亏格零稳定映射表示的类 $\beta$ 只有有限多个。
先设 $X$ 为 Fano 簇。忽略有限的挠部分后，这些 $\beta$ 是 $H_2(X,\mathbb R)$ 中的格点。
Fano 条件意味着 $-\omega_X$ 位于 Kähler 锥的内部，
所以线性泛函 $\beta\mapsto-\int_\beta\omega_X$ 在有效 $\beta$ 生成的锥的闭包上严格为正（原点除外）。
于是，对任意常数 $C$，满足 $-\int_\beta\omega_X<C$ 的 $\beta$ 只有有限多个。
由式~\eqref{eq:quantum-three-point-degree}，得到 Fano 情形的结论。

再设 $X$ 为齐次空间 $G/P$。此时 $X$ 的上同调已知；
由亏格零稳定映射表示的类 $\beta$，是有限多个类 $\beta_1,\ldots,\beta_r$ 的非负整系数线性组合，
每个 $\beta_i$ 都是非恒定映射 $\mu_i:\mathbb P^1\to X$ 的像类。
由式~\eqref{eq:quantum-three-point-degree}，只需证明
$\int_{\beta_i}\omega_X<0$，等价地，$\int_{\beta_i}c_1(T_X)>0$。

可迁群作用表明，$T_X$ 由全局截面生成，因此 $\mu_i^*(T_X)$ 也由全局截面生成。
它是 $\mathbb P^1$ 上的向量丛，所以
$\mu_i^*(T_X)\simeq\bigoplus_j\mathcal O_{\mathbb P^1}(a_j)$，其中所有 $a_j\geq0$。
若能证明某个 $a_j>0$，便有 $\int_{\beta_i}c_1(T_X)>0$。
否则，$\mu_i^*(T_X)$ 平凡，而由于 $T_{\mathbb P^1}\simeq\mathcal O_{\mathbb P^1}(2)$，
自然映射 $T_{\mathbb P^1}\to\mu_i^*(T_X)$ 必为零。
这与 $\mu_i$ 非恒定矛盾，命题得证。
\end{proof}
命题的第二部分说明，对任意射影空间、Grassmann 簇、二次超曲面、旗流形，
以及这些簇的任意乘积，定义~\ref{def:small-quantum-product} 中的和都是有限的。

接着证明 $H^*(X,\mathbb C)$ 在小量子乘积下成为环。
\begin{theorem}\label{thm:small-quantum-ring}
$H^*(X,\mathbb C)$ 在小量子乘积下是有单位元的超交换环。
\end{theorem}
\begin{proof}
所谓超交换，是指
\[
a*b=(-1)^{\deg a\deg b}b*a,
\]
与杯积类似。这是因为\hyperref[sec:gw-properties]{第~\ref*{sec:gw-properties} 节}的等变性公理给出
$\langle I_{0,3,\beta}\rangle(a,b,T_i)=(-1)^{\deg a\deg b}\langle I_{0,3,\beta}\rangle(b,a,T_i)$。
至于单位元，取 $T_0=1=[X]\in H^0(X,\mathbb C)$ 为 $X$ 的基本类。
它是杯积的单位元，下面证明也是 $*$ 的单位元。
先注意，当 $\beta\ne0$ 时，$\langle I_{0,3,\beta}\rangle(a,T_0,T_i)=0$，
这是等变性公理与\IndexLocation{indexloc:0502:242}基本类公理的直接推论。因此
\[
\begin{aligned}
a*T_0&=\sum_i\langle I_{0,3,0}\rangle(a,T_0,T_i)T^i\\
&=\sum_i\left(\int_Xa\cup T_0\cup T_i\right)T^i=a\cup T_0=a.
\end{aligned}
\]
第二个等式由\IndexLocation{indexloc:0505:242}点映射公理得到，第三个等式由 $T^i$ 的定义得到。
所以 $T_0$ 是 $*$ 的单位元。

最后一步是证明 $*$ 满足结合律。
\hyperref[sec:gw-properties]{第~\ref*{sec:gw-properties} 节}的\IndexLocation{indexloc:0507:242}分裂公理在证明中起关键作用。
给定 $H^*(X,\mathbb C)$ 中的齐次类 $a,b,c$，先考虑 $(a*b)*c$。
用定义~\ref{def:small-quantum-product} 展开，并利用
\SourcePage{243}
$T^i=\sum_jg^{ij}T_j$，得到
\[
\sum_\beta\sum_{t,u}\left(\sum_{\beta=\beta_1+\beta_2}\sum_{r,s}
g^{rs}\langle I_{0,3,\beta_1}\rangle(a,b,T_r)
\langle I_{0,3,\beta_2}\rangle(T_s,c,T_t)\right)g^{tu}q^\beta T_u.
\]
现考虑 Gromov--Witten 类 $I_{0,4,\beta}(a,b,c,T_t)$。
分裂公理描述它沿式~\eqref{eq:gw-splitting-gluing-map} 定义的映射
\[
\varphi:\overline M_{0,3}\times\overline M_{0,3}\longrightarrow\overline M_{0,4}\simeq\mathbb P^1
\]
的拉回。注意，除非 $I_{0,4,\beta}(a,b,c,T_t)\in H^0(\overline M_{0,4})$，此拉回为零；
下面假设它确实位于此空间。由于 $\varphi$ 是 $\mathbb P^1$ 中一点的嵌入，可将
\begin{equation}\label{eq:small-associativity-left}
\varphi^*I_{0,4,\beta}(a,b,c,T_t)
\end{equation}
视为一个数。分裂公理于是说，式~\eqref{eq:small-associativity-left}
恰好是上面显示式的大圆括号内的量。

转而考虑 $a*(b*c)$，得到
\[
\sum_\beta\sum_{t,u}\left(\sum_{\beta=\beta_1+\beta_2}\sum_{r,s}
g^{rs}\langle I_{0,3,\beta_1}\rangle(b,c,T_r)
\langle I_{0,3,\beta_2}\rangle(a,T_s,T_t)\right)g^{tu}q^\beta T_u.
\]
由等变性公理，
$\langle I_{0,3,\beta_2}\rangle(a,T_s,T_t)=
(-1)^{\deg a\deg T_s}\langle I_{0,3,\beta_2}\rangle(T_s,a,T_t)$；
而由式~\eqref{eq:quantum-three-point-degree}，有
$\deg a+\deg T_s+\deg T_t\equiv0\pmod2$。因此
\[
\langle I_{0,3,\beta_2}\rangle(a,T_s,T_t)
=(-1)^{\deg a(\deg a+\deg T_t)}\langle I_{0,3,\beta_2}\rangle(T_s,a,T_t).
\]
所以 $a*(b*c)$ 表达式中括号内的量为
\begin{equation}\label{eq:small-associativity-right}
(-1)^{\deg a(\deg a+\deg T_t)}\widetilde\varphi^*I_{0,4,\beta}(b,c,a,T_t),
\end{equation}
其中 $\widetilde\varphi$ 与 $\varphi$ 类似。
注意 $\widetilde\varphi$ 是 $\mathbb P^1$ 中另一个点的嵌入，故与 $\varphi$ 线性等价。

比较 $(a*b)*c$ 与 $a*(b*c)$ 的表达式，可见，只要证明
式~\eqref{eq:small-associativity-left} 与式~\eqref{eq:small-associativity-right} 相等，定理立即成立。
所以只需证明
\[
I_{0,4,\beta}(a,b,c,T_t)
=(-1)^{\deg a(\deg a+\deg T_t)}I_{0,4,\beta}(b,c,a,T_t).
\]
这是等变性公理与由式~\eqref{eq:gw-top-degree-condition} 得到的
$\deg a+\deg b+\deg c+\deg T_t\equiv0\pmod2$ 的直接推论。
\end{proof}
注意，上述证明与\hyperref[subsec:plane-genus-zero]{第~\ref*{subsec:plane-genus-zero} 小节}
证明 $\mathbb P^2$ 的 Gromov--Witten 不变量递推关系时所用的方法，有许多共同之处。
这里出现了奇数次上同调类，因此还需处理符号。
关于结合律的直观解释，见 \SourceCite{Morrison7}。

除赋予环结构外，小量子乘积还有其他性质。
先描述 $a*b$ 的次数。下一小节的例子将说明，即使 $a,b$ 齐次，$a*b$ 也未必齐次。
不过，$a*b$ 中出现的次数并非任意。
\SourcePage{244}
\begin{proposition}\label{prop:small-quantum-degree}
若 $a,b\in H^*(X,\mathbb C)$ 齐次，则其小量子乘积为齐次类之和 $a*b=\sum_lc_l$，且
\[
\deg c_l\equiv\deg a+\deg b\pmod2.
\]
此外，若 $\omega_X\in H^2(X,\mathbb Z)$ 的类可被整数 $r$ 整除，则
$\deg c_l\equiv\deg a+\deg b\pmod{2r}$。
\end{proposition}
\begin{proof}
由于只有 $\deg T_i+\deg T_j=2\dim X$ 时 $g^{ij}$ 才可能非零，
\IndexLocation{indexloc:1133:244}次数条件~\eqref{eq:quantum-three-point-degree} 表明，在定义~\ref{def:small-quantum-product} 中，
$T_j$ 的系数非零只可能在
\[
\deg a+\deg b=\deg T_j-2\int_\beta\omega_X
\]
时发生。命题立即成立。
\end{proof}
还可描述小量子乘积与 $X$ 上交配对的关系。
记此配对为 $g(a,b)=\int_Xa\cup b$。
\begin{proposition}\label{prop:quantum-frobenius-pairing}
对所有 $a,b,c\in H^*(X,\mathbb C)$：
\begin{enumerate}[label=(\roman*)]
\item $g(a*b,c)=g(a,b*c)=\sum_{\beta\in H_2(X,\mathbb Z)}\langle I_{0,3,\beta}\rangle(a,b,c)q^\beta$。
\item $\int_Xa*b=g(a,b)$。
\end{enumerate}
\end{proposition}
\begin{proof}
容易证明 $g(T_i*T_j,T_k)=\sum_\beta\langle I_{0,3,\beta}\rangle(T_i,T_j,T_k)q^\beta$。
转而考虑 $g(T_i,T_j*T_k)$，杯积的超交换性给出
\[
g(T_i,T_j*T_k)=(-1)^{\deg T_i(\deg T_j+\deg T_k)}g(T_j*T_k,T_i),
\]
而由上述等式，$g(T_j*T_k,T_i)=\sum_\beta\langle I_{0,3,\beta}\rangle(T_j,T_k,T_i)q^\beta$。
第一部分随即由\IndexLocation{indexloc:0501:244}等变性公理得到。为证明第二部分，注意
\[
\int_Xa*b=\int_XT_0\cup(a*b)=g(T_0,a*b)=g(T_0*a,b)=g(a,b),
\]
因为 $T_0$ 是杯积与小量子乘积共同的单位元。
\end{proof}
此命题表明，$H^*(X,\mathbb C)$ 在 $*$ 与 $g$ 下是一个 \emph{Frobenius 代数}，见定义~\ref{def:commutative-frobenius-algebra}。
\index{Frobenius代数@Frobenius 代数|hyperpage}
命题~\ref{prop:quantum-frobenius-pairing} 也与\IndexLocation{indexloc:1129:244}\IndexLocation{indexloc:1144:244}小量子上同调的另一种略有不同的方法有关。
该方法从\emph{三点函数} $\langle a,b,c\rangle$ 出发，定义为
\begin{equation}\label{eq:a-three-point-function}
\index{三点函数!量子上同调|hyperpage}
\langle a,b,c\rangle=\sum_{\beta\in H_2(X,\mathbb Z)}\langle I_{0,3,\beta}\rangle(a,b,c)q^\beta.
\end{equation}
也称 $\langle a,b,c\rangle$ 为 \emph{A \IndexLocation{indexloc:0185:244}模型\IndexLocation{indexloc:0195:244}关联函数}。
容易看出，$a,b$ 的小量子乘积是满足
$g(a*b,c)=\langle a,b,c\rangle$（对所有 $c\in H^*(X,\mathbb C)$）的唯一元素 $a*b\in H^*(X,\mathbb C)$。
文献中常采用这种方式定义量子上同调，见 \SourceCite{MS,Morrison7,Voisin3}。
还要注意，式~\eqref{eq:small-quantum-product} 与式~\eqref{eq:a-three-point-function} 给出
\begin{equation}\label{eq:small-product-from-correlations}
a*b=\sum_i\langle a,b,T_i\rangle T^i,
\end{equation}
\SourcePage{245}
这再次证明，量子上同调由三点函数确定。
最后还应提到，命题~\ref{prop:quantum-frobenius-pairing} 给出
\begin{equation}\label{eq:three-point-quantum-integral}
\langle a,b,c\rangle=\int_Xa*b*c,
\end{equation}
所以三点函数能完全用小量子乘积表示。

对于 Calabi--Yau 流形，小量子乘积尤其好。
\begin{corollary}\label{cor:cy-quantum-hodge}
若 $V$ 为 Calabi--Yau 流形，则小量子乘积满足 $\deg(a*b)=\deg a+\deg b$。
此外，若 $a\in H^{p,q}(V)$、$b\in H^{r,s}(V)$，则 $a*b\in H^{p+r,q+s}(V)$。
\end{corollary}
\begin{proof}
第一个断言立即由命题~\ref{prop:small-quantum-degree} 得到，因为 $\omega_V$ 平凡。
设 $V$ 的维数为 $d$。对于第二个断言，由命题~\ref{prop:quantum-frobenius-pairing} 的第一部分，只需证明
\[
\langle I_{0,3,\beta}\rangle(a,b,c)=0,
\qquad c\in H^{u,v}(V),\quad(u,v)\ne(d-p-r,d-q-s).
\]
但由于 $\overline M_{0,3}$ 是一点，定义~\ref{def:gw-algebraic-invariants} 给出
\[
\langle I_{0,3,\beta}\rangle(a,b,c)=\int_{\pi_{1*}(\xi)}a\otimes b\otimes c,
\]
其中 $\pi_1:\overline M_{0,3}(V,\beta)\to V^3$ 是\IndexLocation{indexloc:0276:245}赋值映射，$\xi$ 为\IndexLocation{indexloc:1415:245}虚拟基本类。
由于 $\xi$ 是预期维数的代数类，由式~\eqref{eq:stable-map-expected-dimension}，
$\pi_{1*}(\xi)\in H_{2d}(V^3)$ 是代数圈的同调类。
上述积分只有在 $a\otimes b\otimes c$ 的型为 $(d,d)$ 时才可能非零，结论随之成立。
\end{proof}
推论~\ref{cor:cy-quantum-hodge} 表明：当 $V$ 为 Calabi--Yau 流形时，
$\bigoplus_{p=0}^dH^{p,p}(V)$ 是赋予小量子乘积后的 $H^*(V,\mathbb C)$ 的子环。
本章稍后研究 A 侧 Hodge \IndexLocation{indexloc:0019:245}结构变分时，将用到这一点。

\subsection{例题}\label{subsec:small-quantum-examples}
\label{subsec:a-model-correlation}
现在计算\IndexLocation{indexloc:1129:245}\IndexLocation{indexloc:1133:245}\IndexLocation{indexloc:1137:245}\IndexLocation{indexloc:1138:245}小量子上同调环的几个例子。
\SourceExampleNumber{8.1.2}{1}
\begin{example}\label{ex:projective-quantum-ring}
对 $\mathbb P^r$，使用上同调基 $H^i$，$0\leq i\leq r$，
其中 $H$ 为超平面类，指数表示杯积。我们断言
\begin{equation}\label{eq:projective-quantum-product}
H^i*H^j=\begin{cases}
H^{i+j},&i+j\leq r,\\
qH^{i+j-r-1},&i+j\geq r+1,
\end{cases}
\end{equation}
其中 $q=e^{2\pi i\int_\ell\omega}$，$\ell$ 是 $\mathbb P^r$ 中的直线。
为证明此式，取 $\beta=d[\ell]$，$d\geq0$。
由式~\eqref{eq:quantum-three-point-degree}，不变量
$\langle I_{0,3,d}\rangle(H^i,H^j,H^k)$（这里用 $d$ 代替 $\beta$）为零，除非
\begin{equation}\label{eq:projective-quantum-degree}
i+j+k=r+d(r+1).
\end{equation}
由于 $i,j,k\leq r$，$d$ 只能为 $0$ 或 $1$；进一步，若 $i+j\leq r$，则 $d=0$。
因此容易得到，当 $i+j\leq r$ 时，$H^i*H^j=H^{i+j}$。
特别地，对于 $0\leq i\leq r$，$H^i$ 对杯积与小量子乘积具有相同含义。

为完成证明，只需证明 $H*H^r=qH^0$。
此时由式~\eqref{eq:projective-quantum-degree}，$d=1$、$k=r$。
于是容易得到
\[
H*H^r=\langle I_{0,3,1}\rangle(H,H^r,H^r)qH^0.
\]
\SourcePage{246}
与\hyperref[subsec:plane-genus-zero]{第~\ref*{subsec:plane-genus-zero} 小节}一样，
$\langle I_{0,3,1}\rangle(H,H^r,H^r)=\langle I_{0,3,1}\rangle(H,[pt],[pt])=1$，
因为 $\mathbb P^r$ 中两个点确定唯一一条直线，它在第三个点与 $H$ 相交。
这证明了式~\eqref{eq:projective-quantum-product}。
还要注意，由于 $\omega_{\mathbb P^r}=\mathcal O_{\mathbb P^r}(-r-1)$ 可被 $r+1$ 整除，
此式是命题~\ref{prop:small-quantum-degree} 的一个很好的例子。

由于 $H^0=[\mathbb P^r]$ 是单位元，且在小量子乘积下 $H^{r+1}=qH^0$，
$\mathbb P^r$ 的小量子环同构于
$\mathbb C[H]/\langle H^{r+1}-q\rangle$。
另一方面，通常的上同调环同构于 $\mathbb C[H]/\langle H^{r+1}\rangle$。
当 $\omega$ 的虚部增大时，$q\to0$，可见小量子乘积在极限中成为杯积。
\end{example}

\SourceExampleNumber{8.1.2}{2}
\begin{example}\label{ex:batyrev-quantum-ring}
\SourceCite{Batyrev2} 为光滑射影环面簇 $X$ 定义了一个“量子上同调环”。
结果表明，若 $X$ 为 Fano 簇，这个环与小量子上同调环一致；否则可能不同。
为避免混淆，记 Batyrev 的环为 $H_\omega^*(X)$。
\SourceCite{QR} 证明，对于射影空间上的某些射影丛，$H_\omega^*(X)$ 与小量子上同调环一致。
例题~\ref{ex:toric-fano-and-hirzebruch-quantum-ring} 将解释为何它们在 Fano 情形一致（推广 \SourceCite{QR} 的结果），
并给出非 Fano 情形的一个例子，表明一般情况下 $H_\omega^*(X)$ 确实与小量子上同调环不同。
该例题还将说明，无论如何，$H_\omega^*(X)$ 与小量子上同调环密切相关。
现在讨论 $H_\omega^*(X)$ 的定义和性质。
我们会看到，它的性质与预期的量子上同调性质相似（至少在 Fano 情形下应当如此）。

置 $X=X_\Sigma$，其中 $\Sigma$ 是 $N_{\mathbb R}$ 中给出 $X$ 的扇。
设 $v_1,\ldots,v_s\in N$ 是 $\Sigma$ 的一维锥的本原整生成元；
每个 $v_i$ 对应环面不变除子 $D_i\subset X$。
先回忆 $X$ 的复系数上同调环的通常描述。
给每个 $v_i,D_i$ 配一个变量 $x_i$，得到多项式环 $\mathbb C[x_1,\ldots,x_s]$。
考虑其中两个理想：
\[
\begin{aligned}
P(\Sigma)&=\left\langle\sum_{i=1}^s\langle m,v_i\rangle x_i:m\in M\right\rangle,\\
SR(\Sigma)&=\langle x_{i_1}\cdots x_{i_k}:\{v_{i_1},\ldots,v_{i_k}\}\not\subset\sigma
\text{，对所有 }\sigma\in\Sigma\rangle.
\end{aligned}
\]
$SR(\Sigma)$ 是 $\Sigma$ 的 \emph{Stanley--Reisner 理想}。
众所周知，见 \SourceCite{Fulton3,Oda}，将 $[D_i]\in H^2(X,\mathbb C)$ 映到 $x_i$，诱导同构
\begin{equation}\label{eq:toric-cohomology-sr}
\index{Stanley--Reisner理想@Stanley--Reisner 理想|hyperpage}
H^*(X,\mathbb C)\simeq\mathbb C[x_1,\ldots,x_s]/(P(\Sigma)+SR(\Sigma)).
\end{equation}
还需要 Batyrev 对此作的一个修改，见 \SourceCite{Batyrev2}。
\hyperref[subsec:homogeneous-ring]{第~\ref*{subsec:homogeneous-ring} 小节}定义：
若 $\mathcal P=\{v_{i_1},\ldots,v_{i_k}\}$ 不包含在任何 $\sigma\in\Sigma$ 中，
但 $\mathcal P$ 的每个真子集都包含在某个锥中，则称 $\mathcal P$ 为\emph{本原集合}。
于是容易看出
\begin{equation}\label{eq:sr-primitive-collections}
SR(\Sigma)=\langle x_{i_1}\cdots x_{i_k}:\{v_{i_1},\ldots,v_{i_k}\}\text{ 为本原集合}\rangle.
\end{equation}
现在考虑量子上同调。对于 $\mathbb P^r$，我们从杯积的 $\mathbb C[h]/\langle h^{r+1}\rangle$
过渡到量子乘积的 $\mathbb C[h]/\langle h^{r+1}-q\rangle$。
为将此推广到环面簇 $X$，先引入一些记号。
设 $\mathcal P=\{v_{i_1},\ldots,v_{i_k}\}$ 为本原集合。
则和 $v_{\mathcal P}=v_{i_1}+\cdots+v_{i_k}$ 位于某个锥 $\sigma\in\Sigma$ 中，
\SourcePage{247}
所以可写成 $v_{\mathcal P}=c_1v_{j_1}+\cdots+c_lv_{j_l}$，
其中 $v_{j_1},\ldots,v_{j_l}$ 为 $\sigma$ 的生成元，$c_1,\ldots,c_l\geq0$。
这给出 $v_1,\ldots,v_s$ 之间的非平凡关系
\begin{equation}\label{eq:quantum-primitive-relation}
v_{i_1}+\cdots+v_{i_k}-c_1v_{j_1}-\cdots-c_lv_{j_l}=0.
\end{equation}
这些关系按如下方式对应于同调类。
由于 $X$ 光滑，$\operatorname{Pic}(X)=A_{n-1}(X)=H^2(X,\mathbb Z)$，
且光滑情形中的 $A_{n-1}(X)$ 无挠，因此\hyperref[chap:toric]{第~\ref*{chap:toric} 章}正合列~\eqref{eq:toric-chow-sequence} 的对偶为
\[
0\longrightarrow H_2(X,\mathbb Z)\longrightarrow\mathbb Z^s\longrightarrow N\longrightarrow0.
\]
这里 $(\lambda_1,\ldots,\lambda_s)\in\mathbb Z^s$ 映到 $\sum_{i=1}^s\lambda_iv_i\in N$。
于是有自然同构
\begin{equation}\label{eq:toric-relation-homology}
\left\{(\lambda_1,\ldots,\lambda_s)\in\mathbb Z^s:\sum_{i=1}^s\lambda_iv_i=0\right\}
\simeq H_2(X,\mathbb Z).
\end{equation}
特别地，关系~\eqref{eq:quantum-primitive-relation} 对应于一个同调类，
记为 $\beta(\mathcal P)\in H_2(X,\mathbb Z)$，参见引理~\ref{lem:mori-wall-curve}。
\SourceCite[定理 2.15]{Batyrev1} 中，Batyrev 证明 $\beta(\mathcal P)$ 是有效的一维圈类。

现给定复化 Kähler 类 $\omega\in H^{1,1}(X)=H^2(X,\mathbb C)$，定义\emph{量子 Stanley--Reisner 理想}
\begin{equation}\label{eq:quantum-sr-primitive}
SR_\omega(\Sigma)=\langle x_{i_1}\cdots x_{i_k}
-q^{\beta(\mathcal P)}x_{j_1}^{c_1}\cdots x_{j_l}^{c_l}:\mathcal P\text{ 为本原集合}\rangle,
\end{equation}
其中，如通常一样，$q^{\beta(\mathcal P)}=e^{2\pi i\int_{\beta(\mathcal P)}\omega}$。
Batyrev 随后定义环
\begin{equation}\label{eq:batyrev-quantum-ring}
H_\omega^*(X)=\mathbb C[x_1,\ldots,x_s]/(P(\Sigma)+SR_\omega(\Sigma)),
\end{equation}
见 \SourceCite[定理 5.3]{Batyrev2}。
称 $H_\omega^*(X)$ 为 \emph{Batyrev 量子环}。
\index{Batyrev量子环@Batyrev 量子环|hyperpage}\index{Stanley--Reisner理想@Stanley--Reisner 理想!量子理想|hyperpage}
为了解环~\eqref{eq:batyrev-quantum-ring} 与通常上同调环~\eqref{eq:toric-cohomology-sr} 的关系，
让 $\omega$ 的虚部趋于无穷大。
由于 $\beta(\mathcal P)$ 是有效圈，$\int_{\beta(\mathcal P)}\omega$ 的虚部也增大，
从而 $q^{\beta(\mathcal P)}\to0$。
因此，$\omega$ 的虚部趋于无穷大时，量子 Stanley--Reisner 理想~\eqref{eq:quantum-sr-primitive}
成为通常的 Stanley--Reisner 理想~\eqref{eq:sr-primitive-collections}。
这说明 $X$ 的通常上同调环是 $H_\omega^*(X)$ 的一个极限。

检验 $X=\mathbb P^r$ 的情形。
锥的生成元为 $e_0,\ldots,e_r$，其中 $e_1,\ldots,e_r$ 生成格，$e_0=-\sum_{i=1}^re_i$。
所以 $H^*(\mathbb P^r,\mathbb C)$ 的描述中的理想 $P(\Sigma)$ 为
$\langle x_1-x_0,\ldots,x_r-x_0\rangle$。
此外，$\mathcal P=\{e_0,\ldots,e_r\}$ 是唯一的本原集合。
由于 $e_0+\cdots+e_r=0$，式~\eqref{eq:quantum-primitive-relation} 中的 $c_j$ 都为 $0$。
因此，量子 Stanley--Reisner 理想~\eqref{eq:quantum-sr-primitive} 为
$SR_\omega(\Sigma)=\langle x_0\cdots x_r-q^{\beta(\mathcal P)}\rangle$。
结合 $P(\Sigma)$ 与 $SR_\omega(\Sigma)$，环~\eqref{eq:batyrev-quantum-ring} 化为
$\mathbb C[x_0]/\langle x_0^{r+1}-q^{\beta(\mathcal P)}\rangle$。
还需确定 $\beta(\mathcal P)\in H_2(\mathbb P^r,\mathbb Z)$。
因为它有效，而关系 $e_0+\cdots+e_r=0$ 生成 $e_i$ 之间全部关系构成的 $\mathbb Z$ 模，
必有 $\beta(\mathcal P)=[\ell]$，另见引理~\ref{lem:mori-wall-curve} 及其证明之后的一句话。
因此，重新得到前面对 $\mathbb P^r$ \IndexLocation{indexloc:1129:247}小量子上同调的描述。

环 $H_\omega^*(X)$ 还有其他描述。若
\begin{equation}\label{eq:toric-quantum-general-relation}
\sum_{i=1}^sa_iv_i=\sum_{i=1}^sb_iv_i,\qquad a_i,b_i\in\mathbb Z_{\geq0},
\end{equation}
\SourcePage{248}
则由式~\eqref{eq:toric-relation-homology}，
$a-b=(a_1-b_1,\ldots,a_s-b_s)$ 给出元素 $\beta(a,b)\in H_2(X,\mathbb Z)$。
\SourceCite{Batyrev2} 中，Batyrev 证明量子 Stanley--Reisner 理想也可写为
\begin{equation}\label{eq:quantum-sr-all-relations}
SR_\omega(\Sigma)=\langle x_1^{a_1}\cdots x_s^{a_s}
-q^{\beta(a,b)}x_1^{b_1}\cdots x_s^{b_s}:a_i,b_i\text{ 满足式~\eqref{eq:toric-quantum-general-relation}}\rangle.
\end{equation}
这一描述说明，$H_\omega^*(X)$ 只依赖 $\Sigma$ 的一维锥与复化 Kähler 类 $\omega$。
所以，具有相同一维骨架的不同扇，量子上同调密切相关。
例如，设 $X_1,X_2$ 是光滑环面簇，它们的扇有相同一维骨架，
并假设它们的 Kähler 锥在\hyperref[sec:gkz-decomposition]{第~\ref*{sec:gkz-decomposition} 节}
描述的 GKZ 分解中有共同的墙。
那么 $X_1,X_2$ 通过小双有理变换联系；若该变换保典范，则是一条翻越。
因此，它们的通常上同调环可能不同构。
不过，利用 Batyrev 量子环，可以按以下方式从一个过渡到另一个：
\[
\begin{array}{ccccc}
\text{通常}&&\text{量子}\quad\text{量子}&&\text{通常}\\[-2pt]
H^*(X_1,\mathbb C)&\longleftarrow&H_\omega^*(X_1)\longleftrightarrow H_\omega^*(X_2)
&\longrightarrow&H^*(X_2,\mathbb C).
\end{array}
\]
左边箭头表示：当 $\omega$ 的虚部在 $X_1$ 的 Kähler 锥中趋于无穷大时，
$X_1$ 的 Batyrev 量子环趋于通常上同调环。
右边箭头类似，只是 $\omega$ 的虚部在 $X_2$ 的 Kähler 锥中趋于无穷大。
中间箭头则对应于 $\omega$ 穿过这堵墙。
只要保留相同的一维锥，式~\eqref{eq:quantum-sr-all-relations} 就允许改变扇，这使上述过渡成为可能。

可如下理解 $SR_\omega(\Sigma)$ 的两种描述：
式~\eqref{eq:quantum-sr-all-relations} 只用一维骨架，说明如何将 $H_\omega^*(X)$ 扩展到整个 GKZ 分解；
式~\eqref{eq:quantum-sr-primitive} 则通过本原集合使用整个扇，
说明如何在任意一个 Kähler 锥中取极限。进一步讨论与例子见 \SourceCite{Batyrev2}。

最后，在令 $\omega$ 的虚部趋于无穷大时，我们一直以很朴素的方式取理想 $SR_\omega(\Sigma)$ 的极限。
正确做法当然是使用单项式序与 Gröbner 基。
事实上，$\omega$ 的虚部给出一个单项式序，见 \SourceCite{Batyrev2} 定理 5.3 的证明。
此外，按所给单项式序将所有可能的 $\operatorname{Im}\omega$ 分组，
就得到 GKZ 分解的一个细分，称为 \emph{Gröbner 扇}。
\SourceCite{Sturmfels1} 描述了这个扇，\SourceCite{HLY1,HLY2} 将它用于镜像对称，研究超几何方程的解。
\index{Gröbner扇@Gröbner 扇|hyperpage}
\end{example}

\SourceExampleNumber{8.1.2}{3}
\begin{example}\label{ex:cy-small-quantum}
最后，考虑光滑 Calabi--Yau 三维簇 $V$ 的小量子上同调。
先证明，大多数情况下，小量子乘积与杯积一致。
\begin{lemma}\label{lem:cy-quantum-cup}
若 $V$ 是光滑 Calabi--Yau 三维簇，则只要 $\deg a\ne2$ 或 $\deg b\ne2$，就有 $a*b=a\cup b$。
\end{lemma}
\begin{proof}
先注意，当 $\beta\ne0$，且某个 $\deg\alpha_i\ne2$ 时，
$\langle I_{0,3,\beta}\rangle(\alpha_1,\alpha_2,\alpha_3)=0$。
\hyperref[subsec:cy-genus-zero]{第~\ref*{subsec:cy-genus-zero} 小节}已对五次三维簇证明了这一点，
该证明也适用于任意 Calabi--Yau 三维簇。引理随即得到。
\end{proof}
还需计算 $H^2(V,\mathbb C)$ 上的量子乘积。先设 $V$ 为五次三维簇。
如往常一样，用 $\langle I_{0,3,d}\rangle$ 代替 $\langle I_{0,3,\beta}\rangle$。
一旦知道 $H*H$，就唯一确定小量子乘积，其中
\SourcePage{249}
$H\in H^2(V,\mathbb Z)$ 是超平面类。
由推论~\ref{cor:cy-quantum-hodge}，$H*H\in H^4(V,\mathbb C)$，
然后定义~\ref{def:small-quantum-product} 与 $\int_VH^3=5$ 给出
\[
H*H=\frac15\sum_{d=0}^{\infty}\langle I_{0,3,d}\rangle(H,H,H)q^dH^2.
\]
还有由式~\eqref{eq:a-three-point-function}、\eqref{eq:three-point-quantum-integral} 定义的三点函数，
这里变成
\begin{equation}\label{eq:quintic-quantum-correlation}
\langle H,H,H\rangle=\int_VH*H*H
=\sum_{d=0}^{\infty}\langle I_{0,3,d}\rangle(H,H,H)q^d.
\end{equation}
\hyperref[subsec:cy-genus-zero]{第~\ref*{subsec:cy-genus-zero} 小节}中已看到，
此式右边就是\hyperref[chap:quintic]{第~\ref*{chap:quintic} 章}研究的 A \IndexLocation{indexloc:0187:249}\IndexLocation{indexloc:0188:249}\IndexLocation{indexloc:1183:249}模型\IndexLocation{indexloc:0195:249}关联函数 $\langle H,H,H\rangle$。
所以，五次三维簇的三点函数由小量子乘积给出。
如式~\eqref{eq:small-product-from-correlations} 所示，反过来也成立；结合上述两个公式，得到
\[
H*H=\langle H,H,H\rangle\frac15H^2.
\]
因此，量子乘积用三点函数表示出来。

对于任意 Calabi--Yau 三维簇 $V$，若限制到 $a,b,c\in H^2(V,\mathbb C)$，情形类似。
由式~\eqref{eq:a-three-point-function}、\eqref{eq:three-point-quantum-integral}，三点函数为
\begin{equation}\label{eq:cy-quantum-correlation-gw}
\begin{aligned}
\langle a,b,c\rangle&=\int_Va*b*c=\sum_\beta\langle I_{0,3,\beta}\rangle(a,b,c)q^\beta\\
&=\langle I_{0,3,0}\rangle(a,b,c)+\sum_{\beta\ne0}\langle I_{0,0,\beta}\rangle q^\beta
\int_\beta a\int_\beta b\int_\beta c.
\end{aligned}
\end{equation}
\hyperref[subsec:instanton-numbers]{第~\ref*{subsec:instanton-numbers} 小节}已介绍瞬子数 $n_\beta$，满足
\begin{equation}\label{eq:cy-instanton-mobius}
\langle I_{0,0,\beta}\rangle=\sum_{\beta=k\gamma}n_\gamma k^{-3}.
\end{equation}
如式~\eqref{eq:cy-a-correlation-instanton} 一样，由此得到
\begin{equation}\label{eq:cy-quantum-correlation-instanton}
\langle a,b,c\rangle=\int_Va*b*c=\int_Va\cup b\cup c
+\sum_{\substack{\beta\ne0\\\beta\text{ 有效}}}n_\beta\frac{q^\beta}{1-q^\beta}\int_\beta a\int_\beta b\int_\beta c.
\end{equation}
光滑 Calabi--Yau 三维簇三点函数的这一经典公式，最初对五次三维簇给出，见 \SourceCite{CdGP}，
并经常出现在文献中，见 \SourceCite{Witten5,AGM3,Morrison5,Voisin3}。
通过式~\eqref{eq:cy-instanton-mobius} 定义瞬子数 $n_\beta$ 时，这一公式具有明确意义。
但由\hyperref[subsec:instanton-numbers]{第~\ref*{subsec:instanton-numbers} 小节}可知，
$n_\beta$ 未必计数 $V$ 同调类为 $\beta$ 的有理曲线。
不过，猜想~\ref{conj:symplectic-instantons} 断言，$n_\beta$ 应与类 $\beta$ 中按定向带符号计数的 $J$ 全纯曲线数目密切相关。

能用量子上同调表达 Calabi--Yau 三维簇的三点函数，令人十分满意。
此外，如\hyperref[subsec:small-quantum-definition]{第~\ref*{subsec:small-quantum-definition} 小节}所见，
反过来也成立：可以用三点函数构造量子上同调。
\end{example}

\SourceExampleNumber{8.1.2}{4}
\begin{example}\label{ex:flop-quantum-invariance}
\index{翻越!量子上同调|hyperpage}
本例研究翻越对 Calabi--Yau 三维簇 $V$ 的\IndexLocation{indexloc:1129:249}\IndexLocation{indexloc:1134:249}\IndexLocation{indexloc:1139:249}\IndexLocation{indexloc:1142:249}\IndexLocation{indexloc:1202:249}小量子上同调的影响。
用\hyperref[subsec:threefold-kahler-moduli]{第~\ref*{subsec:threefold-kahler-moduli} 小节}的术语，设 $\phi:V\to\overline V$ 为本原 I 型收缩，
给出翻越 $V\dashrightarrow V'$，两边被收缩的曲线分别为 $C\subset V$ 与 $C'\subset V'$。
\SourcePage{250}
此时有自然识别 $H^2(V)\simeq H^2(V')$，在该识别下，Kähler 锥 $K(V)$ 与 $K(V')$ 沿共同的面相接。

为比较 $V,V'$ 的量子上同调，设可约曲线 $C$ 由同调类为 $\beta$、
彼此不相交且刚性嵌入的若干 $\mathbb P^1$ 构成。
进一步假设，$C$ 与所有其他类不为 $\beta$ 的曲线不相交。

在这一情形下，已知 $V,V'$ 有相同的 Hodge 数，但交形式可能不同，
所以通常上同调环可能不同构。不过，其小量子上同调环密切相关，见 \SourceCite{Witten5,AGM3,Morrison5}。
利用式~\eqref{eq:cy-quantum-correlation-instanton}，将 $V$ 上的三点函数写成
\[
\begin{aligned}
\langle a,b,c\rangle={}&\int_Va\cup b\cup c
+n_\beta\frac{q^\beta}{1-q^\beta}\int_\beta a\int_\beta b\int_\beta c\\
&+\sum_{\substack{\gamma\ne\beta,0\\\gamma\text{ 有效}}}n_\gamma\frac{q^\gamma}{1-q^\gamma}
\int_\gamma a\int_\gamma b\int_\gamma c.
\end{aligned}
\]
过渡到 $V'$ 时，右边第一项可能变化，第二项也可能变化，第三项则不变。
这是因为模空间 $\overline M_{0,3}(V,\gamma)$ 与 $V'$ 上相应的模空间同构，且这些同构保持虚拟基本类。
上面列出的文献中一个简单计算表明：只要通过 $H^2(V)\simeq H^2(V')$ 变换 $a,b,c$，
并通过该映射的 Poincaré 对偶变换 $\beta$，那么从 $V$ 过渡到 $V'$ 时，前两项之和
\[
\int_Va\cup b\cup c+n_\beta\frac{q^\beta}{1-q^\beta}\int_\beta a\int_\beta b\int_\beta c
\]
经相应解析延拓后\emph{保持不变}。

从形式角度看，在上述识别下，$V,V'$ 的三点函数给出相同的表达式；
容易看出，采用适当形式系数的小量子上同调环也相应一致。
不过，$V$ 的小量子乘积只有在 $\omega$ 的虚部足够深入 $K(V)$ 时才猜想收敛，$K(V')$ 的情形类似。
在 $H^2(V)\simeq H^2(V')$ 下，两个 Kähler 锥的内部不相交，
所以这个识别不再直接给出两者在数值系数下的量子上同调环同构，因为此时是在互不相交的 $\omega$ 集合上取值。
得到的是 $V$ 的三点函数到 $V'$ 的三点函数的解析延拓。
这与例题~\ref{ex:batyrev-quantum-ring} 研究翻越对环面簇量子上同调的影响时十分相似。
与那里一样，量子上同调使我们能将 $V$ 的通常上同调环形变为 $V'$ 的上同调环。

这些讨论都针对对应于 I 型收缩的 Kähler 锥的面。
\SourceCite{Wilson4} 研究了 III 型收缩的影响。
\end{example}

小量子乘积一直是活跃的研究领域。除刚才讨论的例子外，还能描述其他有趣\IndexLocation{indexloc:1140:250}簇的小量子上同调，包括：
\begin{itemize}
\item \textbf{Grassmann 簇。}Grassmann 簇的小量子环最初由 \SourceCite{Vafa2} 描述，
随后 \SourceCite{Witten6,ST} 给出了数学处理。
其他文献见 \SourceCite{MS,FP,Bertram,BCF}。
有关主题包括 Giambelli 与 Pieri 公式的量子版本、与 Verlinde 代数的联系，以及 Vafa--Intriligator 公式。
\SourcePage{251}
\item \textbf{旗流形。}旗流形的小量子环与 Toda 格有引人入胜的关系，最初在 \SourceCite{GK} 中描述。
其他文献包括 \SourceCite{AS,Kim1,Givental3,CiocanF,Kostant}。
\item \textbf{Fano 流形。}关于 Fano 流形的小量子环，已经有大量工作。
这类流形包括 Grassmann 簇、旗流形，以及某些射影超曲面与完全交。
相关论文包括 \SourceCite{ST,Beauville1,TX,CJ}。
\item \textbf{有理曲面。}\SourceCite{CM} 描述了一般有理曲面的小量子环。
对三次曲面，量子乘积的结合律基本上等价于存在 $27$ 条直线，且每条与另外 $10$ 条相交。
\end{itemize}

\subsection{系数}\label{subsec:formal-small-quantum}
\index{量子上同调!系数环|hyperpage}
\index{Novikov环@Novikov 环|hyperpage}
在迄今的讨论中，定义小量子乘积时使用复系数上同调。
这是因为，若 $\omega\in H^2(X,\mathbb C)$ 是复化 Kähler 类，则
$q^\beta=e^{2\pi i\int_\beta\omega}\in\mathbb C$。
但尚不知定义~\ref{def:small-quantum-product} 中的求和是否收敛。
避免这一问题的一种方法，是将 $q^\beta$ 视为形式变量，而非数值。
有几种方式可使这一点精确化。

对于 Fano 流形，由命题~\ref{prop:quantum-finite-fano}，定义~\ref{def:small-quantum-product} 中的和有限。
此时，可取系数环为半群环
\[
R=\mathbb Q[q^\beta:\beta\in H_2(X,\mathbb Z)].
\]
也就是说，每个 $\beta\in H_2(X,\mathbb Z)$ 给出形式符号 $q^\beta$，满足
$q^{\beta_1}q^{\beta_2}=q^{\beta_1+\beta_2}$；$R$ 的元素是各 $q^\beta$ 的有限有理系数和。
例如，对 $\mathbb P^r$，得到关于 $q$ 的 Laurent 多项式环 $R=\mathbb Q[q,q^{-1}]$。
对任意 Fano 流形，小量子乘积定义在 $H^*(X,R)$ 上。

若 $X$ 为任意射影流形，这种办法可能失效，因为定义~\ref{def:small-quantum-product} 中的和可能无限。
但由\IndexLocation{indexloc:0500:251}有效性公理，可假设 $\beta$ 有效。
更准确地说，可假设 $\beta$ 位于整 Mori 锥 $M(X)_{\mathbb Z}$ 中。
考虑
\[
\mathcal R=\mathbb Q[[q^\beta:\beta\in M(X)_{\mathbb Z}]],
\]
其中双括号表示允许无限和。两个这样的元素的乘积为
\[
\left(\sum_{\beta\in M(X)_{\mathbb Z}}a_\beta q^\beta\right)
\left(\sum_{\beta\in M(X)_{\mathbb Z}}b_\beta q^\beta\right)
=\sum_{\beta\in M(X)_{\mathbb Z}}
\left(\sum_{\beta=\beta_1+\beta_2}a_{\beta_1}b_{\beta_2}\right)q^\beta.
\]
这个和有意义，因为对任意 $\beta\in M(X)_{\mathbb Z}$，括号内的和有限：
Mori 锥的自由部分由强凸锥中的格点构成。
因此 $\mathcal R$ 是环，可视为整 Mori 锥的形式半群环。
以五次三维簇为例，$\mathcal R=\mathbb Q[[q]]$ 是关于 $q$ 的形式幂级数环。
一般地，容易看出小量子乘积定义在 $H^*(X,\mathcal R)$ 上。

最后，若 $(M,\omega)$ 为辛流形，可用定义~\ref{def:small-quantum-product} 定义小量子乘积，
但必须使用以下系数环的上同调：
\SourcePage{252}
\emph{Novikov 环} $\Lambda(\omega,\mathbb Q)$。
它由所有形式和 $\sum_{\beta\in H_2(M,\mathbb Z)}a_\beta q^\beta$，$a_\beta\in\mathbb Q$ 构成，
满足对任意 $C\in\mathbb R$，集合
\[
\left\{\beta:a_\beta\ne0\text{ 且 }\int_\beta\omega<C\right\}
\]
有限。直接可证，$\Lambda(\omega,\mathbb Q)$ 为环，
且小量子乘积定义在 $H^*(M,\Lambda(\omega,\mathbb Q))$ 上。

也可使用其他系数。在代数情形，可用 $\mathcal R$ 在适当自同构群下的不变子环代替 $\mathcal R$，见 \SourceCite{Morrison7}。
在辛情形，往往将 Novikov 环限制到球面同调类，见 \SourceCite{MS}，另见 \SourceCite{Getzler1}。

\section{大量子上同调}\label{sec:big-quantum-cohomology}
\index{量子上同调!大量子上同调|hyperpage}
\hyperref[sec:small-quantum-cohomology]{第~\ref*{sec:small-quantum-cohomology} 节}
讨论的量子上同调，并未用尽 Gromov--Witten 不变量 $\langle I_{0,n,\beta}\rangle$ 提供的枚举信息。
原因在于，小量子乘积由 $\langle I_{0,3,\beta}\rangle$ 定义，而更有趣的不变量常在 $n\geq3$ 时出现。
例如，\hyperref[chap:gromov-witten]{第~\ref*{chap:gromov-witten} 章}已看到，对 $\mathbb P^2$，数
\[
N_d=\langle I_{0,3d-1,d}\rangle(\underbrace{[pt],\ldots,[pt]}_{3d-1\text{ 次}})
\]
有一些显著性质。但 $\mathbb P^2$ 的小量子乘积只用到
$N_1=\langle I_{0,3,1}\rangle([pt],[pt],[\ell])=1$。
本节探讨一种称为\emph{大量子上同调}的量子上同调版本，它考虑全部 Gromov--Witten 不变量。

\subsection{作为超流形的 $H^*(X,\mathbb C)$}\label{subsec:cohomology-supermanifold}
\index{超流形|hyperpage}\index{超交换|hyperpage}
定义大量子乘积之前，需讨论 $H^*(X,\mathbb C)$ 的超流形结构。
固定有理上同调 $H^*(X,\mathbb Q)$ 的齐次基 $T_0=1,T_1,\ldots,T_m$。
对每个 $T_i$，引入与它次数相同的变量 $t_i$；这些变量是超流形 $H^*(X,\mathbb C)$ 的局部坐标。
$t_i$ 超交换，即 $t_it_j=(-1)^{\deg t_i\deg t_j}t_jt_i$；若 $t_i$ 的次数为奇数，则 $t_i^2=0$。
所以环 $\mathbb C[t_0,\ldots,t_m]$ 是一个关于偶变量的多项式环，与一个关于奇变量的外代数的张量积。
对不含 $t_i$ 的任意单项式 $t^\alpha$，定义偏导数算子 $\partial/\partial t_i$ 为
\[
\frac\partial{\partial t_i}(t_i^k\cdot t^\alpha)=kt_i^{k-1}\cdot t^\alpha.
\]
于是可以计算任意 $F\in\mathbb C[t_0,\ldots,t_m]$ 的 $\partial F/\partial t_i$，并且
\[
\frac{\partial^2F}{\partial t_i\partial t_j}
=(-1)^{\deg t_i\deg t_j}\frac{\partial^2F}{\partial t_j\partial t_i}.
\]
这些都容易扩展到形式幂级数环 $\mathbb C[[t_0,\ldots,t_m]]$。
\SourcePage{253}
关于超流形的更多内容，请参阅 \SourceCite{Manin3}。
如 \SourceCite[第 2 讲]{Alvarez} 所讨论，研究费米子系统时，超交换变量自然出现在量子物理中。

\subsection{Gromov--Witten 势}\label{subsec:gw-potential}
\index{Gromov--Witten势@Gromov--Witten 势|hyperpage}
定义大量子乘积的关键工具是 \emph{Gromov--Witten 势}。
如上，固定有理上同调 $H^*(X,\mathbb Q)$ 的基 $T_0=1,T_1,\ldots,T_m$，
置 $\gamma=\sum_{i=0}^mt_iT_i$，其中 $t_i$ 为超交换变量，$\deg t_i=\deg T_i$。
\begin{definition}\label{def:gw-potential}
设 $\omega$ 是光滑射影簇 $X$ 上的复化 Kähler 类。
则 \emph{Gromov--Witten 势}为形式和
\[
\Phi(\gamma)=\sum_{n=0}^{\infty}\sum_{\beta\in H_2(X,\mathbb Z)}
\frac1{n!}\langle I_{0,n,\beta}\rangle(\gamma^n)q^\beta,
\qquad q^\beta=e^{2\pi i\int_\beta\omega}.
\]
\end{definition}
在此和中，$\beta=0$ 时隐含 $n\geq3$，因为 $n\leq2$ 时 $\overline M_{0,n}(X,0)$ 不存在。
文献中 Gromov--Witten 势的定义略有差异：一些作者对\emph{所有} $\beta$ 都要求 $n\geq3$，从而截断级数。
当希望将 Gromov--Witten 势与 Gromov--Witten 类联系起来时，这一做法尤其合适。

为使定义有意义，需解释 $\langle I_{0,n,\beta}\rangle(\gamma^n)$ 的含义。
由于超交换性，这一点稍微微妙。思路是展开 $\gamma^n=(\sum_{i=0}^mt_iT_i)^n$，
将 $T_i$ 与 $t_i$ 都视为超交换变量。
采用多重指标记号：指数向量 $\alpha=(a_0,\ldots,a_m)$ 给出总次数
$|\alpha|=a_0+\cdots+a_m$ 的单项式 $T^\alpha,t^\alpha$。
还置 $\alpha!=a_0!\cdots a_m!$。
一个有用的观察是，$\gamma$ 的每项 $t_iT_i$ 都为偶，因此可将 $t_iT_i$ 作为整体交换。
于是，Gromov--Witten 不变量的符号规则给出
\begin{equation}\label{eq:gw-potential-super-expansion}
\frac1{n!}\langle I_{0,n,\beta}\rangle(\gamma^n)
=\sum_{|\alpha|=n}\epsilon(\alpha)\langle I_{0,n,\beta}\rangle(T^\alpha)\frac{t^\alpha}{\alpha!}.
\end{equation}
因子 $\epsilon(\alpha)=\pm1$ 是由等式
\begin{equation}\label{eq:gw-potential-super-sign}
(t_0T_0)^{a_0}\cdots(t_mT_m)^{a_m}
=\epsilon(\alpha)T_0^{a_0}\cdots T_m^{a_m}t_0^{a_0}\cdots t_m^{a_m}
\end{equation}
确定的符号。

将式~\eqref{eq:gw-potential-super-expansion} 代入 $\Phi$ 的定义，
希望将 $\Phi$ 视为\hyperref[subsec:cohomology-supermanifold]{第~\ref*{subsec:cohomology-supermanifold} 小节}
讨论的环 $\mathbb C[[t_0,\ldots,t_m]]$ 中的形式幂级数。
不过，小量子上同调的收敛问题在这里仍然存在。
为求严格，可能需要用\hyperref[subsec:formal-small-quantum]{第~\ref*{subsec:formal-small-quantum} 小节}
的环 $R,\mathcal R$ 或 $\Lambda(\omega,\mathbb Q)$ 代替 $\mathbb C$。
所以，在下面的讨论中，用 $\mathcal C$ 按需要表示
\[
\mathbb C[[t_0,\ldots,t_m]],\quad R[[t_0,\ldots,t_m]],\quad
\mathcal R[[t_0,\ldots,t_m]],\quad\Lambda(\omega,\mathbb Q)[[t_0,\ldots,t_m]]
\]
中的一个环。

由于 Gromov--Witten 势 $\Phi$ 包含全部可能的亏格零不变量，
可将它视为树级 Gromov--Witten 不变量的生成函数。
考虑到阶乘与不交换变量，更准确地说，$\Phi$ 是
\SourcePage{254}
$\langle I_{0,n,\beta}\rangle$ 的\emph{超交换指数型生成函数}。
从物理角度看，$\Phi$ 自然解释为亏格零自由能。

Gromov--Witten 势有时称为亏格零 Gromov--Witten 势。
\emph{亏格 $g$ 的 Gromov--Witten 势}可用亏格 $g$ 的不变量，以类似方式定义。

椭圆曲线 $E$ 给出一个简单的 Gromov--Witten 势例子。
由于不存在非恒定映射 $\mathbb P^1\to E$，只需考虑由点映射公理确定的
Gromov--Witten 不变量 $\langle I_{0,3,0}\rangle$。
取 $T_0=1\in H^0(E,\mathbb C)$、$T_3=[pt]\in H^2(E,\mathbb C)$，
再取 $T_1,T_2\in H^1(E,\mathbb C)$，满足 $\int_ET_1\cup T_2=1$。
于是计算得到
\begin{equation}\label{eq:elliptic-gw-potential}
\Phi=\tfrac12t_0^2t_3-t_0t_1t_2.
\end{equation}
负号来自式~\eqref{eq:gw-potential-super-expansion} 中的 $\epsilon(\alpha)$。
\hyperref[sec:computing-gw-ii]{第~\ref*{sec:computing-gw-ii} 节}将讨论一些不那么平凡的 Gromov--Witten 势例子。

利用 Gromov--Witten 势，如下定义大量子乘积。
\begin{definition}\label{def:big-quantum-product}
设 $\Phi$ 是光滑射影簇 $X$ 的 Gromov--Witten 势。定义
\[
T_i*T_j=\sum_k\frac{\partial^3\Phi}{\partial t_i\partial t_j\partial t_k}T^k,
\]
其中 $g^{kl}$ 如定义~\ref{def:small-quantum-product}，$T^k=\sum_lg^{kl}T_l$。
线性扩张后，得到上同调 $H^*(X,\mathcal C)$ 上的\emph{大量子乘积}。
\end{definition}
\index{量子乘积!大量子乘积|hyperpage}
注意，$n<3$ 的 Gromov--Witten 不变量是 $\Phi$ 中关于 $t_i$ 次数至多为二的项的系数，
所以在 Gromov--Witten 势的定义中保留或删去它们，并不影响大量子乘积的定义。

在椭圆曲线例子~\eqref{eq:elliptic-gw-potential} 中，$\Phi=\tfrac12t_0^2t_3-t_0t_1t_2$。
使用前面的上同调基 $T_0,T_1,T_2,T_3$，可计算量子乘积
\[
T_1*T_2=\frac{\partial^3\Phi}{\partial t_1\partial t_2\partial t_0}T_3
=\frac{\partial^3}{\partial t_1\partial t_2\partial t_0}
\left(\tfrac12t_0^2t_3-t_0t_1t_2\right)T_3=T_3,
\]
因为 $t_1,t_2$ 为奇变量。容易看出，这个例子中的大量子乘积与通常的杯积一致。
下面的引理有助于处理 Gromov--Witten 势的\IndexLocation{indexloc:0567:254}三阶偏导数。
\begin{lemma}\label{lem:gw-potential-third-derivatives}
对所有 $i,j,k$，
\[
\frac{\partial^3\Phi}{\partial t_i\partial t_j\partial t_k}
=\sum_{n=0}^{\infty}\sum_{\beta\in H_2(X,\mathbb Z)}
\frac1{n!}\langle I_{0,n+3,\beta}\rangle(T_i,T_j,T_k,\gamma^n)q^\beta.
\]
\end{lemma}
\begin{proof}
设指数向量 $\alpha=(a_0,\ldots,a_m)$ 满足 $|\alpha|=n$，
且 $\eta=\alpha-e_i-e_j-e_k$ 的各分量非负，$e_i,e_j,e_k$ 为标准基向量。
若能证明
\[
\frac{\partial^3}{\partial t_i\partial t_j\partial t_k}
\left(\epsilon(\alpha)\langle I_{0,n,\beta}\rangle(T^\alpha)\frac{t^\alpha}{\alpha!}\right)
=\epsilon(\eta)\langle I_{0,n,\beta}\rangle(T_i,T_j,T_k,T^\eta)\frac{t^\eta}{\eta!},
\]
\SourcePage{255}
则引理由式~\eqref{eq:gw-potential-super-expansion} 得到。
先设 $t_i,t_j,t_k$ 都为奇变量。可假设 $a_i=a_j=a_k=1$。
利用 $\langle I_{0,n,\beta}\rangle$ 与 $t_i$ 的符号规则，可写成
\[
\epsilon(\alpha)\langle I_{0,n,\beta}\rangle(T^\alpha)\frac{t^\alpha}{\alpha!}
=\epsilon(\eta+e_i+e_j+e_k)\langle I_{0,n,\beta}\rangle(T_k,T_j,T_i,T^\eta)
\,t_kt_jt_i\frac{t^\eta}{\eta!}.
\]
求三次偏导数，得到
\[
\frac{\partial^3}{\partial t_i\partial t_j\partial t_k}
\left(\epsilon(\alpha)\langle I_{0,n,\beta}\rangle(T^\alpha)\frac{t^\alpha}{\alpha!}\right)
=\epsilon(\eta+e_i+e_j+e_k)\langle I_{0,n,\beta}\rangle(T_k,T_j,T_i,T^\eta)\frac{t^\eta}{\eta!}.
\]
因此只需证明
\[
\epsilon(\eta+e_i+e_j+e_k)\langle I_{0,n,\beta}\rangle(T_k,T_j,T_i,T^\eta)
=\epsilon(\eta)\langle I_{0,n,\beta}\rangle(T_i,T_j,T_k,T^\eta).
\]
由式~\eqref{eq:gw-potential-super-sign}，$\epsilon(\eta)=(-1)^{s(s+1)/2}$，
其中 $s$ 是 $t^\eta$ 中奇变量的个数（可假设奇变量只出现一次幂）。
又由于 $t_i,t_j,t_k$ 为奇变量，式~\eqref{eq:gw-top-degree-condition} 给出 $s+3\equiv0\pmod2$。
所需等式随即成立。其他情形类似，留给读者。
\end{proof}

\subsection{WDVV 方程与结合律}\label{subsec:wdvv-associativity}
\index{WDVV方程@WDVV 方程|hyperpage}
下一个目标，是证明 $H^*(X,\mathcal C)$ 在大量子乘积下成为环。
与\hyperref[sec:small-quantum-cohomology]{第~\ref*{sec:small-quantum-cohomology} 节}一样，难点是结合律。
但这里出现了一个显著现象：大量子上同调的结合律等价于 $\Phi$ 满足的某个偏微分方程，
即所谓 WDVV（Witten--Dijkgraaf--Verlinde--Verlinde）方程，见 \SourceCite{Witten2,DVV1,DVV2}。
\begin{theorem}\label{thm:wdvv}
Gromov--Witten 势 $\Phi$ 对所有 $i,j,k,l$ 满足方程
\[
\begin{aligned}
\sum_{a,b}\frac{\partial^3\Phi}{\partial t_i\partial t_j\partial t_a}g^{ab}
\frac{\partial^3\Phi}{\partial t_b\partial t_k\partial t_l}
={}&(-1)^{\deg t_i(\deg t_j+\deg t_k)}\\
&\cdot\sum_{a,b}\frac{\partial^3\Phi}{\partial t_j\partial t_k\partial t_a}g^{ab}
\frac{\partial^3\Phi}{\partial t_b\partial t_i\partial t_l}.
\end{aligned}
\]
\end{theorem}
\begin{proof}
这里只概述主要思想，详细证明见 \SourceCite{KoM1}。
为简便，再假设 $H^*(X,\mathbb C)$ 只有偶数次上同调，这样就不必处理符号。

证明采用与 $\mathbb P^2$ 的 Gromov--Witten 不变量递推式~\eqref{eq:plane-recursion}
相同的策略。将式~\eqref{eq:plane-xi} 的类 $\xi$ 换成 Gromov--Witten 类
\[
I_{0,n+4,\beta}(\gamma^n,T_i,T_j,T_k,T_l)\in H^*(\overline M_{0,n+4},\mathcal C),
\]
其中 $\gamma=\sum_{i=0}^mt_iT_i$。
将在两个线性等价除子上对此类求值，下面描述它们。
用集合 $\mathcal I=\{1,\ldots,n,i,j,k,l\}$ 标记 $I_{0,n+4,\beta}$ 的 $n+4$ 个自变量
（这是对记号的轻微滥用）。
与式~\eqref{eq:plane-boundary-relation} 前的讨论一样，
将 $\mathcal I$ 分成不相交的 $A,B$，就得到一个除子
$\operatorname{Im}(\varphi_{A,B})\subset\overline M_{0,n+4}$。
\SourcePage{256}
定义
\[
D(i,j\mid k,l)=\sum_{\substack{i,j\in A\\k,l\in B}}\operatorname{Im}(\varphi_{A,B}),
\quad
D(j,k\mid i,l)=\sum_{\substack{j,k\in A\\i,l\in B}}\operatorname{Im}(\varphi_{A,B}).
\]
类似式~\eqref{eq:plane-boundary-relation} 的做法，可得到线性等价关系
\begin{equation}\label{eq:wdvv-boundary-equivalence}
D(i,j\mid k,l)\sim D(j,k\mid i,l).
\end{equation}
现考虑 WDVV 方程左边的和。利用 $T^k=\sum_lg^{kl}T_l$ 与引理~\ref{lem:gw-potential-third-derivatives}，可写为
\[
\sum\frac1{n_1!n_2!}\langle I_{0,n_1+3,\beta_1}\rangle(T_i,T_j,T_a,\gamma^{n_1})
g^{ab}\langle I_{0,n_2+3,\beta_2}\rangle(T_b,T_k,T_l,\gamma^{n_2})q^{\beta_1+\beta_2},
\]
其中求和遍历 $n_1,n_2,\beta_1,\beta_2,a,b$。
将此和分成较小的和，使每个小和中的 $n_1+n_2=n$ 与 $\beta_1+\beta_2=\beta$ 固定。
由交换性，对应于固定 $n,\beta$ 的小和可写成
\[
\sum\frac1{n_1!n_2!}\langle I_{0,n_1+3,\beta_1}\rangle(\gamma^{n_1},T_i,T_j,T_a)
g^{ab}\langle I_{0,n_2+3,\beta_2}\rangle(T_b,\gamma^{n_2},T_k,T_l)q^\beta,
\]
其中现在对满足 $n_1+n_2=n$、$\beta_1+\beta_2=\beta$ 的 $n_1,n_2,\beta_1,\beta_2,a,b$ 求和。
利用式~\eqref{eq:gw-potential-super-expansion} 展开并应用分裂公理，可证明此小和化为
\[
\int_{D(i,j\mid k,l)}\frac1{n!}I_{0,n+4,\beta}(\gamma^n,T_i,T_j,T_k,T_l)q^\beta,
\]
其中 $D(i,j\mid k,l)$ 是上面定义的除子。因此，WDVV 方程的左边为
\[
\sum_{n=0}^{\infty}\sum_{\beta\in H_2(X,\mathbb Z)}
\int_{D(i,j\mid k,l)}\frac1{n!}I_{0,n+4,\beta}(\gamma^n,T_i,T_j,T_k,T_l)q^\beta.
\]
对方程右边进行同样的处理，得到
\[
\sum_{n=0}^{\infty}\sum_{\beta\in H_2(X,\mathbb Z)}
\int_{D(j,k\mid i,l)}\frac1{n!}I_{0,n+4,\beta}(\gamma^n,T_i,T_j,T_k,T_l)q^\beta.
\]
由线性等价关系~\eqref{eq:wdvv-boundary-equivalence}，两者相等，定理立即成立。
\end{proof}
现在容易证明，大量子上同调构成环。
\begin{theorem}\label{thm:big-quantum-ring}
$H^*(X,\mathcal C)$ 在大量子乘积下是单位元为 $T_0=1$ 的交换环。
\end{theorem}
\begin{proof}
由\hyperref[sec:gw-properties]{第~\ref*{sec:gw-properties} 节}的\IndexLocation{indexloc:0501:256}等变性公理，大量子乘积超交换。
下面证明，$T_0=1$ 不仅是杯积的单位元，也是
\SourcePage{257}
大量子乘积的单位元。这很容易：在引理~\ref{lem:gw-potential-third-derivatives} 中取 $T_0=[X]$，得到
\begin{equation}\label{eq:gw-potential-unit-derivative}
\begin{aligned}
\frac{\partial^3\Phi}{\partial t_0\partial t_j\partial t_k}
&=\sum_{n=0}^{\infty}\sum_{\beta\in H_2(X,\mathbb Z)}
\frac1{n!}\langle I_{0,n+3,\beta}\rangle(T_0,T_j,T_k,\gamma^n)q^\beta\\
&=\langle I_{0,3,0}\rangle(T_0,T_j,T_k)=\int_XT_j\cup T_k=g_{jk}.
\end{aligned}
\end{equation}
第二、第三个等式分别由基本类公理与点映射公理得到。
此式给出 $T_0*T_j=T_j$，故 $T_0$ 为单位元。

最后证明结合律。利用 $T^k=\sum_lg^{kl}T_l$ 与大量子乘积的定义，得到
\[
(T_i*T_j)*T_k=\sum_{a,b}\sum_{l,s}
\frac{\partial^3\Phi}{\partial t_i\partial t_j\partial t_a}g^{ab}
\frac{\partial^3\Phi}{\partial t_b\partial t_k\partial t_l}g^{ls}T_s.
\]
以同样方式展开 $(T_j*T_k)*T_i$，并应用定理~\ref{thm:wdvv} 的 WDVV 方程，立即得到
\[
(T_i*T_j)*T_k=(-1)^{\deg T_i(\deg T_j+\deg T_k)}(T_j*T_k)*T_i.
\]
由于 $*$ 超交换，结合律随之成立。
\end{proof}
此定理说明，WDVV 方程推出大量子乘积的结合律。
反过来也成立，所以结合律等价于 WDVV 方程。
这是\hyperref[sec:dubrovin-formalism]{第~\ref*{sec:dubrovin-formalism} 节}将讨论的 Dubrovin 形式体系的一部分。

还应提到，大量子上同调与交乘积 $g(a,b)=\int_Xa\cup b$ 配合良好。
类似命题~\ref{prop:quantum-frobenius-pairing}，可以证明大量子乘积满足
\begin{equation}\label{eq:big-quantum-frobenius}
g(a*b,c)=g(a,b*c),\qquad a,b,c\in H^*(X,\mathcal C).
\end{equation}
所以 $H^*(X,\mathcal C)$ 在大量子乘积下是 Frobenius 代数，并有公式
\[
\int_XT_i*T_j*T_k=\frac{\partial^3\Phi}{\partial t_i\partial t_j\partial t_k}.
\]
大量子乘积有 Künneth 公式，见 \SourceCite{KoM2,Kaufmann}。
此外，\hyperref[subsec:gravitational-correlator-definition]{第~\ref*{subsec:gravitational-correlator-definition} 小节}将讨论一种称为\emph{引力量子乘积}的大量子乘积推广。

\section{计算 Gromov--Witten 不变量，II}\label{sec:computing-gw-ii}
现在重访\hyperref[sec:computing-gw-i]{第~\ref*{sec:computing-gw-i} 节}的若干例子，
看看它们与大量子上同调及 Gromov--Witten 势有何关系。
\subsection{Gromov--Witten 势的结构}\label{subsec:gw-potential-structure}
\index{Gromov--Witten势@Gromov--Witten 势!齐次性|hyperpage}\index{Gromov--Witten势@Gromov--Witten 势!经典部分和量子部分|hyperpage}
如下将看到，Gromov--Witten 势的显式公式可能相当复杂。
不过，$\Phi$ 具有一种结构，能使公式更容易理解。
\SourcePage{258}
首先说明：给各变量赋予适当的次数后，$\Phi$ 是齐次的。沿用上一节的记号，我们有上同调基 $T_i$ 与超交换变量 $t_i$。由式~\eqref{eq:gw-potential-super-expansion}，$\Phi$ 的一个典型项形如
\[
\langle I_{0,n,\beta}\rangle(T_0^{a_0},\ldots,T_m^{a_m})
\frac{t_0^{a_0}}{a_0!}\cdots\frac{t_m^{a_m}}{a_m!}q^\beta,
\]
其中 $a_0+\cdots+a_m=n$。另一方面，次数条件~\eqref{eq:gw-top-degree-condition} 给出
\[
a_0\deg T_0+\cdots+a_m\deg T_m
=2\dim X-2\int_\beta\omega_X+2(3\cdot0-3+n).
\]
利用 $a_0+\cdots+a_m=n$，可将它改写为
\[
a_0(\deg T_0-2)+\cdots+a_m(\deg T_m-2)+2\int_\beta\omega_X=2\dim X-6.
\]
因此，只要给 $t_0,\ldots,t_m$ 和 $q^\beta$ 赋予如下次数，Gromov--Witten 势 $\Phi$ 就是次数为 $2\dim X-6$ 的齐次式：
\begin{equation}\label{eq:gw-potential-grading}
\deg t_i=\deg T_i-2,\qquad \deg q^\beta=2\int_\beta\omega_X.
\end{equation}
也就是说，$t_i$ 的次数要从先前的次数中减去 $2$，因而奇偶性不变。至于 $q^\beta$，要这样赋予次数，就必须把它们看作形式符号，例如\hyperref[subsec:formal-small-quantum]{第~\ref*{subsec:formal-small-quantum} 小节}中的环 $R$、$\mathcal R$ 或 $\Lambda(\omega,\mathbb Q)$ 中的符号，而不能仅看作数值。

Gromov--Witten 势的齐次性在实际计算中十分有用。它说明，势中的许多项因次数不符而为零。应用于\IndexLocation{indexloc:1116:258}量子上同调时，齐次性还说明，$T_i*T_j$ 中 $T_\ell q^\beta$ 的系数具有次数 $\deg T_\ell-(\deg T_i+\deg T_j+2\int_\beta\omega_X)$。正如我们在命题~\ref{prop:small-quantum-degree} 的证明中对小量子乘积所注意到的那样。（不过，如果只关心量子乘积，通常直接使用次数条件比使用齐次性更方便。）

作为一个简单的例子，考虑式~\eqref{eq:elliptic-gw-potential} 中椭圆曲线的 Gromov--Witten 势 $\Phi=\tfrac12t_0^2t_3-t_0t_1t_2$。只要 $t_0,t_1,t_2,t_3$ 的次数分别为 $-2,-1,-1,0$，它就是次数为 $-4$ 的齐次式。

$\Phi$ 还有一个有用的性质：由\IndexLocation{indexloc:0505:258}点映射公理，对应于 $\beta=0$ 的各项化为 $\tfrac16\langle I_{0,3,0}\rangle(\gamma^3)=\tfrac16\int_X\gamma^3$。因此，可将 Gromov--Witten 势分解为“经典”项与“量子”项（这里采用 \SourceCite{FP} 的术语）：
\[
\begin{aligned}
\Phi&=\Phi_{\mathrm{classical}}+\Phi_{\mathrm{quantum}}\\
&=\frac16\int_X\gamma^3+
\sum_n\sum_{\beta\in H_2(X,\mathbb Z)\setminus\{0\}}
\frac1{n!}\langle I_{0,n,\beta}\rangle(\gamma^n)q^\beta.
\end{aligned}
\]
注意，由\IndexLocation{indexloc:0502:258}基本类公理，变量 $t_0$ 只出现在 $\Phi$ 的\IndexLocation{indexloc:0553:258}经典部分 $\Phi_{\mathrm{classical}}$ 中。

$\Phi$ 的另一个好性质是：由\IndexLocation{indexloc:0499:258}除子公理，对应于除子的变量在 $\Phi_{\mathrm{quantum}}$ 中以指数形式出现。为说明这一点，将基元素 $T_i$ 重新排列，使得 $T_1,\ldots,T_r$ 是 $H^2(X,\mathbb Q)$ 的一组基，并写成
\SourcePage{259}
$\gamma=t_0T_0+\delta+\epsilon$，其中 $\delta=\sum_{i=1}^rt_iT_i$，$\epsilon=\sum_{i=r+1}^mt_iT_i$。随后，利用二项式定理和除子公理，得到
\begin{equation}\label{eq:gw-potential-divisor-exponential}
\begin{aligned}
\Phi&=\frac16\int_X\gamma^3+\sum_n\sum_{\beta\ne0}
\frac1{n!}\langle I_{0,n,\beta}\rangle\bigl((\delta+\epsilon)^n\bigr)q^\beta\\
&=\frac16\int_X\gamma^3+\sum_{a,b}\sum_{\beta\ne0}
\frac1{a!b!}\langle I_{0,b,\beta}\rangle(\epsilon^b)
\left(\int_\beta\delta\right)^a q^\beta\\
&=\frac16\int_X\gamma^3+\sum_b\sum_{\beta\ne0}
\frac1{b!}\langle I_{0,b,\beta}\rangle(\epsilon^b)e^{\int_\beta\delta}q^\beta.
\end{aligned}
\end{equation}
\SourceCite{KoM1} 的命题 4.4 讨论了此公式；它对量子上同调还有一些推论，将在\hyperref[subsec:quantum-product-dependence]{第~\ref*{subsec:quantum-product-dependence} 小节}讨论。

\subsection{重访 $\mathbb P^2$}\label{subsec:plane-gw-potential}
\index{Gromov--Witten势@Gromov--Witten 势!射影平面|hyperpage}
对 $\mathbb P^2$，首先计算 Gromov--Witten 势。上同调基为 $T_0=[\mathbb P^2]=1$、$T_1=[\ell]$ 和 $T_2=[pt]$，因此，由式~\eqref{eq:gw-potential-grading}，$t_0,t_1,t_2$ 的次数为 $-2,0,2$，而 $q=q^\ell$ 的次数为 $-6$。照例将 Gromov--Witten 不变量写为 $\langle I_{0,n,d}\rangle$。用此记号，$N_d=\langle I_{0,3d-1,d}\rangle(T_2^{3d-1})$ 是通过 $\mathbb P^2$ 中 $3d-1$ 个一般位置点的 $d$ 次有理曲线数目。

我们知道，$\Phi$ 是次数为 $-2$ 的齐次式，其经典部分为 $\tfrac12(t_0^2t_2+t_0t_1^2)$。量子部分的一个典型项为
\[
\langle I_{0,a+b,d}\rangle(T_1^a,T_2^b)\frac{t_1^a}{a!}\frac{t_2^b}{b!}q^d,
\]
由齐次性，$2b+(-6)d=-2$。这给出 $b=3d-1$，所以该项可写成
\[
\langle I_{0,3d-1,d}\rangle(T_2^{3d-1})
\frac{(dt_1)^a}{a!}\frac{t_2^{3d-1}}{(3d-1)!}q^d
=N_d\frac{(dt_1)^a}{a!}\frac{t_2^{3d-1}}{(3d-1)!}q^d,
\]
其中因子 $d^a$ 来自除子公理。（这重新得到了\hyperref[subsec:plane-genus-zero]{第~\ref*{subsec:plane-genus-zero} 小节}中的一个计算，不过这里的推导更快。）因此，$\mathbb P^2$ 的 Gromov--Witten 势为
\[
\Phi=\frac12(t_0^2t_2+t_0t_1^2)+\sum_{d=1}^{\infty}
N_de^{dt_1}\frac{t_2^{3d-1}}{(3d-1)!}q^d.
\]
现在可以做些有趣的计算。对指标 $(i,j,k,l)=(1,1,2,2)$，定理~\ref{thm:wdvv} 的 WDVV 方程成为
\begin{equation}\label{eq:plane-wdvv}
\Phi_{222}+\Phi_{111}\Phi_{122}=\Phi_{112}^2,
\end{equation}
其中下标表示偏导数。因为式~\eqref{eq:gw-potential-unit-derivative} 给出 $\Phi_{0jk}=g_{jk}$，所以上式容易算出。为利用式~\eqref{eq:plane-wdvv}，注意，$N_d$ 在 $\Phi_{222}$ 中出现在如下项里：
\[
N_d\frac1{(3d-4)!}t_2^{3d-4}q^d.
\]
这提示我们考察 $\Phi_{112}^2$ 和 $\Phi_{111}\Phi_{122}$ 中含 $t_2^{3d-4}q^d$ 的项。在 $\Phi_{112}$ 中，不含 $t_1$ 的项为
\[
N_{d_1}d_1^2\frac1{(3d_1-2)!}t_2^{3d_1-2}q^{d_1},
\]
\SourcePage{260}
所以 $\Phi_{112}^2$ 中含 $t_2^{3d-4}q^d$ 的项为
\[
\sum_{d_1+d_2=d}N_{d_1}N_{d_2}d_1^2d_2^2
\frac1{(3d_1-2)!}\frac1{(3d_2-2)!}t_2^{3d-4}q^d.
\]
类似地，$\Phi_{111}\Phi_{122}$ 中含 $t_2^{3d-4}q^d$ 的项为
\[
\sum_{d_1+d_2=d}N_{d_1}N_{d_2}d_1^3d_2
\frac1{(3d_1-1)!}\frac1{(3d_2-3)!}t_2^{3d-4}q^d.
\]
这一计算留给读者验证。由此，式~\eqref{eq:plane-wdvv} 立即给出递推关系
\[
N_d=\sum_{d_1+d_2=d}N_{d_1}N_{d_2}
\left(d_1^2d_2^2\binom{3d-4}{3d_1-2}
-d_1^3d_2\binom{3d-4}{3d_1-1}\right).
\]

$\mathbb P^2$ 的大量子上同调也容易确定。例如，由于 $\Phi_{011}=1$，有
\[
T_1*T_1=T_2+\Phi_{111}T_1+\Phi_{112}T_0.
\]
另一方面，$\mathbb P^2$ 的小量子乘积满足 $T_1*T_1=T_2$，可见两种量子上同调有明显差别。还要注意，大量子乘积的结合律与\hyperref[sec:dubrovin-formalism]{第~\ref*{sec:dubrovin-formalism} 节}讨论的 WDVV 方程实质上等价。实际使用中，结合律往往更方便。例如，展开 $(T_1*T_1)*T_2=T_1*(T_1*T_2)$，并比较 $T_0$ 的系数，就直接得到式~\eqref{eq:plane-wdvv}。因此，$N_d$ 的递推关系是大量子乘积结合律的直接推论。正如 \SourceCite{FP} 所说：“量子形式体系使人不再非得想出巧妙的办法不可。只需写出结合律方程，便可从中读出计数信息。”\SourceCite{DI} 讨论了这样做的一种系统方法；对 $\mathbb P^3$ 和三维二次超曲面的应用见 \SourceCite{FP}。\SourceCite{Kresch} 还更深入地研究了结合律方程及其与 Kontsevich 第一重构定理（定理~\ref{thm:gw-reconstruction}）的关系。

对 $\mathbb P^2$，还可以将 $N_d$ 推广为 $N_d(a,b,c)$：它计数通过 $a$ 个一般位置点、与 $b$ 条一般位置直线相切，并在另外 $c$ 条一般位置直线上各自指定的一般位置点处与这些直线相切的 $d$ 次有理曲线；这里 $a+b+2c=3d-1$。\SourceCite{EK1} 研究了这些不变量，并导出了 $N_d(a,b,c)$ 的递推关系，推广了 \SourceCite{DI} 给出的 $N_d(a,b,0)$ 的递推关系。此外，可以用 $N_d(a,b,c)$ 定义 $\mathbb P^2$ 大量子上同调的一种推广，称为\emph{接触上同调环} \SourceCite{EK2}。\index{接触上同调环|hyperpage}

\subsection{重访五次三维簇}\label{subsec:quintic-gw-potential}
\index{Gromov--Witten势@Gromov--Witten 势!五次三维簇|hyperpage}
首先研究五次三维簇 $V$ 的 Gromov--Witten 势。$H^*(V,\mathbb C)$ 的一组上同调基为 $T_0=[V]$、$T_1=[H]$、$T_2=[\ell]$、$T_3=[pt]$ 与 $\gamma_1,\ldots,\gamma_{204}$，其中 $\gamma_i$ 是 $H^3(V,\mathbb C)$ 的一组基。对应的超交换变量记为 $t_0,t_1,t_2,t_3$ 和 $u_1,\ldots,u_{204}$。再令 $q_{ij}=\int_V\gamma_i\cup\gamma_j$。于是，$\Phi$ 的经典部分为
\[
\Phi_{\mathrm{classical}}=\frac12t_0^2t_3+t_0t_1t_2+
\frac56t_1^3-\sum_{i<j}q_{ij}t_0u_iu_j,
\]
其中负号来自式~\eqref{eq:gw-potential-super-expansion} 中的 $\epsilon(\alpha)$。
\SourcePage{261}
再看量子部分。我们知道，$\Phi$ 和 $q^\beta$ 的次数均为 $0$，而 $t_0$ 不出现在 $\Phi_{\mathrm{quantum}}$ 中。由式~\eqref{eq:gw-potential-grading}，$t_2,t_3$ 和 $u_i$ 的次数都是正的，所以 $\Phi_{\mathrm{quantum}}$ 只依赖于 $t_1$。又因为 $T_1$ 是除子（超平面截面 $H$）的类，式~\eqref{eq:gw-potential-divisor-exponential} 给出
\begin{equation}\label{eq:quintic-quantum-potential}
\Phi_{\mathrm{quantum}}=\sum_{d\ge1}\langle I_{0,0,d}\rangle e^{dt_1}q^d,
\end{equation}
其中照例用 $\langle I_{0,0,d}\rangle$ 代替 $\langle I_{0,0,\beta}\rangle$。因此，五次三维簇的 Gromov--Witten 势为
\[
\Phi=\frac12t_0^2t_3+t_0t_1t_2+\frac56t_1^3
-\sum_{i<j}q_{ij}t_0u_iu_j+\sum_{d\ge1}\langle I_{0,0,d}\rangle(e^{t_1}q)^d.
\]
由此，描述 $V$ 的大量子乘积相当简单。除 $T_1*T_1$ 外，所有基元素对的乘积都与杯积一致；而这个例外的乘积为
\[
T_1*T_1=\Phi_{111}T_2.
\]
此外，有
\[
\int_V T_1*T_1*T_1=\Phi_{111},
\]
并且，由于 $T_1$ 是 $H$ 的类，$\Phi_{111}$ 由如下公式给出：
\[
\Phi_{111}=5+\sum_{d\ge1}\langle I_{0,0,d}\rangle d^3(e^{t_1}q)^d
=5+\sum_{d\ge1}\langle I_{0,3,d}\rangle(H,H,H)(e^{t_1}q)^d.
\]
在这些公式中，将 $e^{t_1}q$ 换成 $q$，就恰好恢复五次三维簇的小量子上同调。

更一般地，对任意光滑 Calabi--Yau 三维簇，Gromov--Witten 势都可写成特别简洁的形式。首先注意，对于齐次插入类，$\beta\ne0$ 蕴含
\begin{equation}\label{eq:cy-nonzero-degree-insertions}
\langle I_{0,n,\beta}\rangle(\alpha_1,\ldots,\alpha_n)=0,
\quad\text{除非所有 $\alpha_i$ 的次数都为 $2$。}
\end{equation}
若选定的 $H^2(V,\mathbb C)$ 的基为 $T_1,\ldots,T_r$，式~\eqref{eq:cy-nonzero-degree-insertions} 说明 $\Phi_{\mathrm{quantum}}$ 是 $t_1,\ldots,t_r$ 的函数；由式~\eqref{eq:gw-potential-divisor-exponential}，这些变量以指数形式出现。因此，量子部分可写成
\[
\Phi_{\mathrm{quantum}}=\sum_{\beta\ne0}\langle I_{0,0,\beta}\rangle
 e^{\int_\beta\delta}q^\beta,
\]
其中 $\delta=\sum_{i=1}^rt_iT_i$。若去掉原来的因子 $q^\beta$，并令
$q^\beta=e^{\int_\beta\delta}=\prod_{i=1}^re^{t_i\int_\beta T_i}$，Gromov--Witten 势就由如下简洁公式给出：
\begin{equation}\label{eq:cy-total-gw-potential}
\Phi=\Phi_{\mathrm{classical}}+\sum_{\beta\ne0}\langle I_{0,0,\beta}\rangle q^\beta.
\end{equation}
利用式~\eqref{eq:cy-nonzero-degree-insertions} 和~\eqref{eq:cy-total-gw-potential}，得到
\begin{equation}\label{eq:cy-potential-three-derivatives}
\begin{aligned}
\frac{\partial^3\Phi}{\partial t_i\partial t_j\partial t_k}
&=\int_V T_i\cup T_j\cup T_k+
\sum_{\beta\ne0}\langle I_{0,0,\beta}\rangle
\int_\beta T_i\int_\beta T_j\int_\beta T_k\,q^\beta\\
&=\langle I_{0,3,0}\rangle(T_i,T_j,T_k)+
\sum_{\beta\ne0}\langle I_{0,3,\beta}\rangle(T_i,T_j,T_k)q^\beta\\
&=\langle T_i,T_j,T_k\rangle.
\end{aligned}
\end{equation}
\SourcePage{262}
其中第二个等式使用点映射公理和除子公理，第三个等式使用式~\eqref{eq:cy-quantum-correlation-gw}。这说明，将 Gromov--Witten 势改写为式~\eqref{eq:cy-total-gw-potential} 后，定义~\ref{def:big-quantum-product} 就成为小量子乘积的定义。因此，对 Calabi--Yau 三维簇，大量子上同调与小量子上同调实质上没有区别。还可以验证，此时 WDVV 方程归结为杯积的结合律。这与 $\mathbb P^2$ 的情形很不相同：在那里，大量子乘积与 WDVV 方程包含的计数信息远多于小量子上同调。

在\hyperref[subsec:quantum-product-dependence]{第~\ref*{subsec:quantum-product-dependence} 小节}将看到，对一般的簇，只要在 $T_i$ 的次数不为 $2$ 时令参数 $t_i=0$，大量子乘积与小量子乘积也实质上没有区别。

\section{Dubrovin 形式体系}\label{sec:dubrovin-formalism}
\index{Dubrovin 形式体系|hyperpage}\index{挠率|hyperpage}\index{平坦联络|hyperpage}
Gromov--Witten 势的代数结果是大量子乘积。本节将探讨 $\Phi$ 的主要几何结果，即 \emph{Dubrovin 联络}。\hyperref[sec:a-model-vhs]{第~\ref*{sec:a-model-vhs} 节}将利用 Dubrovin 形式体系，在 Calabi--Yau 流形的 Kähler 模空间上构造一种自然出现的 Hodge 结构变分。\index{Dubrovin 联络|hyperpage}

我们将 Dubrovin 联络的定义从 Gromov--Witten 势 $\Phi$ 推广到任意势函数 $F$；为此，依照 \SourceCite{Dubrovin1} 发展更一般的理论。

照例，取 $H^*(X,\mathbb C)$ 的基 $T_0=1,\ldots,T_m$，以及相应的超交换变量 $t_0,\ldots,t_m$。交配对的矩阵记为 $(g_{kl})$，其逆矩阵记为 $(g^{kl})$。令 $F\colon H^*(X,\mathbb C)\to\mathbb C$ 是关于 $t_i$ 的偶形式幂级数。在较好的情形，$F$ 在超流形 $H^*(X,\mathbb C)$ 的某个开集上收敛。给定 $F$ 后，定义张量 $A_{ijk}$ 为
\begin{equation}\label{eq:dubrovin-cubic-tensor}
A_{ijk}=\frac{\partial^3F}{\partial t_i\partial t_j\partial t_k}.
\end{equation}
若用下标表示偏导数，则须注意 $A_{ijk}=F_{kji}$。按物理学中通常的“升指标”约定，令
\begin{equation}\label{eq:dubrovin-raised-tensor}
A_{ij}^k=\sum_lA_{ijl}g^{lk}.
\end{equation}
现在可定义两个基本对象：二元运算
\[
T_i*T_j=\sum_kA_{ij}^kT_k,
\]
以及 \emph{Dubrovin 联络}
\[
\nabla^\lambda_{\frac{\partial}{\partial t_i}}
\left(\frac{\partial}{\partial t_j}\right)
=\lambda\sum_kA_{ij}^k\frac{\partial}{\partial t_k},
\]
其中 $\lambda$ 是非零复数。

二元运算 $*$ 与联络 $\nabla^\lambda$ 关系密切：一者的性质可以转化为另一者的性质。下面简要说明这种关系的三个重要方面。
\begin{itemize}
\item \textbf{挠率与交换律。}式~\eqref{eq:dubrovin-cubic-tensor} 和~\eqref{eq:dubrovin-raised-tensor} 容易推出恒等式
\[
A_{ij}^k=(-1)^{\deg t_i\deg t_j}A_{ji}^k.
\]
\SourcePage{263}
利用此式，不难证明联络 $\nabla^\lambda$ 的挠率为零（在超流形意义下），且二元运算 $*$ 超交换。
\item \textbf{曲率与结合律。}计算联络 $\nabla^\lambda$ 的曲率（在超流形意义下）可知，$(\nabla^\lambda)^2$ 是两项之和。由式~\eqref{eq:dubrovin-cubic-tensor} 和~\eqref{eq:dubrovin-raised-tensor}，有
\[
\frac{\partial}{\partial t_l}A_{ij}^k
=(-1)^{\deg t_i\deg t_l}\frac{\partial}{\partial t_i}A_{lj}^k,
\]
所以第一项为零；第二项为零，当且仅当
\[
\sum_aA_{ij}^aA_{ak}^l
=(-1)^{\deg t_i(\deg t_j+\deg t_k)}\sum_aA_{jk}^aA_{ai}^l.
\]
我们认出，这就是 WDVV 方程。于是容易得到
\[
\begin{aligned}
\text{$\nabla^\lambda$ 是平坦联络}
&\Longleftrightarrow\text{$F$ 满足 WDVV 方程}\\
&\Longleftrightarrow\text{$*$ 满足结合律。}
\end{aligned}
\]
\item \textbf{单位元。}我们希望 $T_0$ 是 $*$ 的单位元。容易得到等价关系
\[
\text{对所有 $i,j$，$A_{0ij}=g_{ij}$}
\Longleftrightarrow\text{对所有 $i,j$，$A_{0i}^j=\delta_{ij}$，}
\]
随后直接可证
\[
\begin{aligned}
\text{$T_0$ 是 $*$ 的单位元}
&\Longleftrightarrow\text{对所有 $i,j$，$A_{0ij}=g_{ij}$}\\
&\Longleftrightarrow\text{对所有 $i$，$\nabla^\lambda_{\frac{\partial}{\partial t_0}}
\left(\frac{\partial}{\partial t_i}\right)=\lambda\frac{\partial}{\partial t_i}$。}
\end{aligned}
\]
\end{itemize}
这些断言的证明见 \SourceCite{Dubrovin1,KoM1,Manin3}。

从现在起，始终假定 $F$ 满足 WDVV 方程，且对所有 $i,j$ 都有 $A_{0ij}=g_{ij}$。这意味着 Dubrovin 联络 $\nabla^\lambda$ 平坦且无挠，而 $\nabla^\lambda_{\frac{\partial}{\partial t_0}}$ 在 $\frac{\partial}{\partial t_i}$ 上的作用是乘以 $\lambda$。就二元运算 $*$ 而言，$H^*(X,\mathbb C)$ 有如下性质。
\begin{proposition}\label{prop:dubrovin-frobenius}
若 $F$ 满足 WDVV 方程，且对所有 $i,j$ 都有 $A_{0ij}=g_{ij}$，则 $H^*(X,\mathbb C)$ 在 $*$ 下是以 $T_0$ 为单位元的 Frobenius 代数。
\end{proposition}
\begin{proof}
假设说明，$*$ 给出以 $T_0$ 为单位元的超交换代数。因此只须证明它满足 Frobenius 性质。式~\eqref{eq:dubrovin-cubic-tensor} 还给出
\[
A_{ijk}=(-1)^{\deg t_i(\deg t_j+\deg t_k)}A_{jki}.
\]
我们又知道 $F$ 是关于 $t_i$ 的偶函数。这些事实蕴含
\[
g(T_i*T_j,T_k)=g(T_i,T_j*T_k).
\]
所以 $H^*(X,\mathbb C)$ 在 $*$ 下是 Frobenius 代数。
\end{proof}
此前研究的势函数 $\Phi$ 与 $\Phi_{\mathrm{small}}$ 都满足命题~\ref{prop:dubrovin-frobenius} 的条件。这就是大量子乘积与小量子乘积使 $H^*(X,\mathbb C)$ 成为 Frobenius 代数的原因。

Dubrovin 联络的构造在超流形 $M=H^*(X,\mathbb C)$ 的切丛上定义了 $\nabla^\lambda$。若将 $\frac{\partial}{\partial t_i}$ 与 $T_i$ 等同，就得到平凡丛
\SourcePage{264}
$H^*(X,\mathbb C)\times M$ 上的联络 $\nabla^\lambda$。容易验证，这个新联络由
\[
\nabla^\lambda_{\frac{\partial}{\partial t_i}}(T_j)=\lambda T_i*T_j
\]
给出。因此，可以将 Dubrovin 联络看作由量子乘法给出的联络。下一节将用到这种观点。

Dubrovin 形式体系还有若干方面，下面继续讨论。首先注意，到目前为止的理论并未使用\hyperref[subsec:gw-potential-structure]{第~\ref*{subsec:gw-potential-structure} 小节}所述 $\Phi$ 的齐次性。用 Dubrovin 联络的语言来表述，齐次性意味着存在一个偶向量场 $E$，称为 \emph{Euler 向量场}，它与 $\nabla^\lambda$ 有良好的相容关系。\index{Euler 向量场|hyperpage}

其次，$\nabla^\lambda$ 是由一个参数 $\lambda$ 参数化的联络族。若将 $\lambda$ 看作 $\mathbb P^1$ 的一个仿射坐标，\SourceCite{KoM1,Manin3} 说明，$\nabla^\lambda$ 可延拓为 $H^*(X,\mathbb C)\times\mathbb P^1$ 切丛上的联络 $\widehat\nabla$，在 $\lambda=0,\infty$ 处有奇点。当 $X=\mathbb P^2$ 时，$\widehat\nabla$ 的 Laplace 变换与第六 Painlevé 方程有一些令人意外的联系，见 \SourceCite{Manin3}。

超流形 $H^*(X,\mathbb C)$，连同度量 $g$ 及满足 WDVV 方程的偶势函数 $F$，给出的结构是 \emph{Frobenius 流形} 的一个例子 \SourceCite{Dubrovin1,Dubrovin2}。\index{Frobenius 流形|hyperpage}这有一些有趣的推论。例如，考虑 $F$ 的幂级数展开。忽略关于 $t_i$ 的次数不超过 $2$ 的项，可将 $F$ 写成
\[
F=\sum_{n=3}^{\infty}\sum_{|\alpha|=n}c_\alpha t^\alpha.
\]
随后，通过等式
\[
c_\alpha=\frac{\epsilon(\alpha)}{\alpha!}\langle I_n\rangle(T^\alpha),\qquad n\ge3,
\]
定义 $\langle I_n\rangle\colon H^*(X,\mathbb C)^{\otimes n}\to\mathbb C$，其中采用式~\eqref{eq:gw-potential-super-expansion} 的记号。于是可将 $F$ 写成
\[
F=\sum_{n=3}^{\infty}\frac1{n!}\langle I_n\rangle(\gamma^n),
\]
这是一个熟悉的形式。\SourceCite{Manin3} 中，Manin 证明了一个引人注目的结果：对 $n\ge3$，存在映射
\[
I_n\colon H^*(X,\mathbb C)^{\otimes n}\longrightarrow
H^*(\overline M_{0,n},\mathbb C),
\]
它们满足\IndexLocation{indexloc:0507:264}分裂公理和\IndexLocation{indexloc:0501:264}等变性公理，并按通常的方式确定 $\langle I_n\rangle$，即
\[
\langle I_n\rangle(\alpha_1,\ldots,\alpha_n)
=\int_{\overline M_{0,n}}I_n(\alpha_1,\ldots,\alpha_n).
\]
这些 $I_n$（$n\ge3$）组成 Manin 所称的\emph{上同调场论}，而 $\langle I_n\rangle$ 是它的\emph{\IndexLocation{indexloc:0127:264}\IndexLocation{indexloc:0183:264}关联函数}。\index{上同调场论|hyperpage}结论是：$H^*(X,\mathbb C)$ 上任意满足 WDVV 方程的偶函数 $F$，都会给出一种与 Gromov--Witten 类惊人相似的结构。还要指出，\hyperref[subsec:wdvv-associativity]{第~\ref*{subsec:wdvv-associativity} 小节}提到的量子上同调 Künneth 公式，更一般地也适用于上同调场论 \SourceCite{KoM2,Kaufmann}。还可以定义更高亏格的上同调场论 \SourceCite{Manin3,KK}。
\SourcePage{265}
最后，讨论 Dubrovin 形式体系时，务必记住，函数 $F$ 可以是形式幂级数。因此，联络 $\nabla^\lambda$ 也可能是形式联络；为此，它们有时称为\emph{形式 Dubrovin 联络}。这种情形下，必须使用\hyperref[subsec:formal-small-quantum]{第~\ref*{subsec:formal-small-quantum} 小节}讨论的一种系数环，才能定义上同调中相应的环结构。

\section{A 型 Hodge 结构变分}\label{sec:a-model-vhs}
\index{Hodge结构变分@Hodge 结构变分!A 型|hyperpage}
本节将说明，小量子上同调如何在\hyperref[chap:moduli]{第~\ref*{chap:moduli} 章}研究的 Kähler 模空间上给出自然的 Hodge 结构变分。这将使我们能够在\hyperref[sec:mirror-symmetry-revisited]{第~\ref*{sec:mirror-symmetry-revisited} 节}为镜像对建立精确的数学定义。
\subsection{大量子上同调与小量子上同调的关系}\label{subsec:quantum-product-dependence}
为研究\hyperref[subsec:a-model-connection]{第~\ref*{subsec:a-model-connection} 小节}定义的 A 模型联络，须先理解大量子上同调与小量子上同调之间的关系。

沿用通常的设定，取 $H^*(X,\mathbb Q)$ 的上同调基 $T_0=1,T_1,\ldots,T_m$，其中 $T_1,\ldots,T_r$ 张成 $H^2(X,\mathbb Q)$。定义~\ref{def:big-quantum-product} 中的大量子乘积是关于 $t_0,\ldots,t_m$ 和 $q^\beta=\exp(2\pi i\int_\beta\omega)$ 的形式幂级数。首先注意，变量 $t_1,\ldots,t_r$ 在此级数中以指数形式出现。更准确地说，定义~\ref{def:big-quantum-product}、引理~\ref{lem:gw-potential-third-derivatives} 和式~\eqref{eq:gw-potential-divisor-exponential} 的推导方法，容易给出如下引理。
\begin{lemma}\label{lem:big-quantum-divisor-exponential}
令 $\delta=\sum_{i=1}^rt_iT_i$ 与 $\epsilon=t_0T_0+\sum_{i=r+1}^mt_iT_i$，则大量子乘积为
\[
T_i*T_j=\sum_k\sum_{n=0}^{\infty}\sum_\beta\frac1{n!}
\langle I_{0,n+3,\beta}\rangle(T_i,T_j,T_k,\epsilon^n)
 e^{\int_\beta\delta}q^\beta T^k.
\]
\end{lemma}
利用此引理，可用几种方式比较大量子上同调与小量子上同调。最简单的是在 $T_i*T_j$ 的公式中令变量 $t_0,\ldots,t_m=0$。用引理~\ref{lem:big-quantum-divisor-exponential} 的记号，就是令 $\delta=\epsilon=0$，得到
\[
T_i*T_j\big|_{\delta=\epsilon=0}
=\sum_k\sum_\beta\langle I_{0,3,\beta}\rangle(T_i,T_j,T_k)q^\beta T^k.
\]
由于
\begin{equation}\label{eq:small-quantum-product-comparison}
T_i*_{\mathrm{small}}T_j=\sum_k\sum_\beta
\langle I_{0,3,\beta}\rangle(T_i,T_j,T_k)q^\beta T^k,
\end{equation}
可知
\[
T_i*T_j\big|_{\delta=\epsilon=0}=T_i*_{\mathrm{small}}T_j.
\]
因此，小量子乘积可看作大量子乘积的限制。

不过，为使用\hyperref[sec:dubrovin-formalism]{第~\ref*{sec:dubrovin-formalism} 节}的 Dubrovin 形式体系，我们采取稍有不同的方法，即“限制到 $H^2(X)$”。其含义如下：大量子乘积用 $t_0,\ldots,t_m$ 与 $q^\beta$ 在 $H^*(X)$ 上形式地定义，所以在 $T_i*T_j$ 的公式中“限制到 $H^2(X)$”是指令 $t_0=t_{r+1}=\cdots=
\SourcePage{266}
t_m=0$。这等价于在引理~\ref{lem:big-quantum-divisor-exponential} 中令 $\epsilon=0$，得到
\begin{equation}\label{eq:restricted-big-quantum-product}
T_i*T_j\big|_{\epsilon=0}=\sum_k\sum_\beta
\langle I_{0,3,\beta}\rangle(T_i,T_j,T_k)e^{\int_\beta\delta}q^\beta T^k.
\end{equation}
将它称为\emph{限制的大量子乘积}。它与小量子乘积~\eqref{eq:small-quantum-product-comparison} 相似，但尚不完全相同。注意到
\[
e^{\int_\beta\delta}q^\beta=e^{\int_\beta(\delta+2\pi i\omega)},
\]
这种相似性就更为明显。因此，式~\eqref{eq:restricted-big-quantum-product} 和~\eqref{eq:small-quantum-product-comparison} 中的指数项都涉及 $H^2(X)$ 中某个元素在 $\beta$ 上的积分，分别为 $\delta+2\pi i\omega$ 和 $2\pi i\omega$。

由此，可以如下使两个公式一致：令 $\omega=0$，并设
\[
\mathbf q^\beta=\exp\left(\int_\beta\delta\right)
=\exp\left(\int_\beta\sum_{i=1}^rt_iT_i\right),
\]
则限制的大量子乘积成为
\begin{equation}\label{eq:restricted-big-quantum-q}
T_i*T_j\big|_{\epsilon=0}=\sum_k\sum_\beta
\langle I_{0,3,\beta}\rangle(T_i,T_j,T_k)\mathbf q^\beta T^k.
\end{equation}
它与式~\eqref{eq:small-quantum-product-comparison} 很好地对应。从形式角度看，将式~\eqref{eq:restricted-big-quantum-q} 中的 $\mathbf q^\beta$ 看作适当系数环中依赖于 $\beta,t_1,\ldots,t_r$ 的形式符号，再使用将 $\mathbf q^\beta$ 映为 $q^\beta$ 的系数环同态，就得到与式~\eqref{eq:small-quantum-product-comparison} 相同的形式级数。

另一方面，将式~\eqref{eq:restricted-big-quantum-q} 和~\eqref{eq:small-quantum-product-comparison} 看作 $H^2(X)$ 上关于变量 $t_1,\ldots,t_r$ 的形式函数时，也可以联系起来。只要在式~\eqref{eq:small-quantum-product-comparison} 中令
\begin{equation}\label{eq:quantum-coordinate-rescaling}
\omega=\frac1{2\pi i}\sum_{i=1}^rt_iT_i,
\end{equation}
这些函数就相同。因此，式~\eqref{eq:small-quantum-product-comparison} 与~\eqref{eq:restricted-big-quantum-q} 的收敛假设也相同。采用这些约定后，无论在形式意义还是解析意义下，大量子乘积限制到 $H^2(X)$ 就是小量子乘积。

\subsection{A 模型联络}\label{subsec:a-model-connection}
\index{复化 Kähler 空间|hyperpage}
本章余下部分中，令 $V$ 是维数 $d\ge3$ 的 Calabi--Yau 流形。回顾\hyperref[chap:introduction]{第~\ref*{chap:introduction} 章}，$V$ 的复化 Kähler 空间是商空间
\[
K_{\mathbb C}(V)=\{\omega\in H^2(V,\mathbb C)\colon
\operatorname{Im}(\omega)\text{ 为 Kähler 类}\}/\operatorname{im}H^2(V,\mathbb Z),
\]
其中 $\operatorname{im}H^2(V,\mathbb Z)$ 是自然映射 $H^2(V,\mathbb Z)\to H^2(V,\mathbb C)$ 的像。下面说明，小量子乘积 $*_{\mathrm{small}}$ 如何在 $K_{\mathbb C}(V)$ 上给出一个有趣的联络。

照例，使用 $H^*(V,\mathbb Q)$ 的上同调基 $T_0,\ldots,T_m$。此外，假定其中位于 $H^2(V,\mathbb Q)$ 的基元素 $T_1,\ldots,T_r$ 满足：
\begin{itemize}
\item $T_1,\ldots,T_r$ 是 $\operatorname{im}H^2(V,\mathbb Z)$ 的一组基；
\item $T_1,\ldots,T_r$ 位于 Kähler 锥的闭包中。
\end{itemize}
由于 $H^{2,0}(V)=0$，容易看出可以找到这样的 $T_1,\ldots,T_r$。

现在引入对应于 $T_1,\ldots,T_r$ 的变量 $u_1,\ldots,u_r$，并令
\begin{equation}\label{eq:a-model-kahler-coordinates}
\omega=\sum_{i=1}^ru_iT_i\in H^2(V,\mathbb R)+iK(V).
\end{equation}
\SourcePage{267}
于是，对任意 $\beta\in H_2(V,\mathbb Z)$，函数 $q^\beta=e^{2\pi i\int_\beta\omega}$ 在 $K_{\mathbb C}(V)$ 上良好定义，因为 $T_i\in\operatorname{im}H^2(V,\mathbb Z)$。\hyperref[chap:moduli]{第~\ref*{chap:moduli} 章}曾用 $t_i$ 而非 $u_i$；很快就会明白这里使用 $u_i$ 的原因。

现在令 $*_{\mathrm{small}}$ 为 $H^*(V,\mathbb C)$ 上的小量子乘积。这需要一些收敛假设，留待稍后讨论。考虑 $K_{\mathbb C}(V)$ 上的平凡丛 $\mathcal H=H^*(V,\mathbb C)\times K_{\mathbb C}(V)$。按如下规则定义 $\mathcal H$ 上的 \emph{A 模型联络} $\nabla$：\index{A 模型联络|hyperpage}
\begin{equation}\label{eq:a-model-connection}
\nabla_{\frac{\partial}{\partial u_i}}(T_j)=T_i*_{\mathrm{small}}T_j,
\qquad 1\le i\le r,\quad0\le j\le m.
\end{equation}
因此，$\nabla$ 由小量子乘积构造，不过其中只用 $H^2(V,\mathbb C)$ 的元素进行乘法。用小量子乘积的定义展开 $\nabla$，得到
\begin{equation}\label{eq:a-model-connection-expanded}
\nabla_{\frac{\partial}{\partial u_i}}(T_j)
=\sum_k\sum_{\beta\in H_2(V,\mathbb Z)}
\langle I_{0,3,\beta}\rangle(T_i,T_j,T_k)q^\beta T^k.
\end{equation}
注意，此联络定义在 $K_{\mathbb C}(V)$ 上，且 $q^\beta$ 是这个空间上的非常值函数。因此，$\nabla$ 记录了当 $\omega$ 改变时，小量子乘积的变化。

严格地说，$\nabla$ 只是一个形式联络，因为式~\eqref{eq:a-model-connection-expanded} 中级数的收敛性仍是猜想。不过，为便于阐述，我们将把 $\nabla$ 看作真正的联络。\hyperref[subsec:a-model-asymptotics]{第~\ref*{subsec:a-model-asymptotics} 小节}将为 $\nabla$ 提出精确的收敛假设。

联络 $\nabla$ 有如下良好性质。
\begin{proposition}\label{prop:a-model-flat}
A 模型联络 $\nabla$ 平坦。
\end{proposition}
\begin{proof}
只要谨慎使用\hyperref[sec:dubrovin-formalism]{第~\ref*{sec:dubrovin-formalism} 节}的理论，这一结论就由 Dubrovin 联络 $\nabla^\lambda$ 的平坦性得到。

我们对取 $q^\beta=1$ 的 Gromov--Witten 势应用 Dubrovin 形式体系。这在 $H^*(V,\mathbb C)$ 上给出平坦联络 $\nabla^\lambda$，定义为
\begin{equation}\label{eq:dubrovin-big-quantum-connection}
\nabla^\lambda_{\frac{\partial}{\partial t_i}}(T_j)=\lambda T_i*T_j,
\end{equation}
其中 $*$ 为大量子乘积，$t_0,\ldots,t_m$ 为 $H^*(V,\mathbb C)$ 的通常变量。像\hyperref[subsec:quantum-product-dependence]{第~\ref*{subsec:quantum-product-dependence} 小节}那样，将 $\nabla^\lambda$ 限制到 $H^2(V,\mathbb C)$，则断言
\begin{equation}\label{eq:a-model-dubrovin-identification}
\nabla=\nabla^\lambda,\qquad\lambda=\frac1{2\pi i}.
\end{equation}
一旦证明此式，命题便由 $\nabla^\lambda$ 的平坦性立即得到。

为证明式~\eqref{eq:a-model-dubrovin-identification}，注意，由式~\eqref{eq:restricted-big-quantum-q}，限制的 Dubrovin 联络可写成
\[
\nabla^\lambda_{\frac{\partial}{\partial t_i}}(T_j)
=\lambda\sum_k\sum_\beta\langle I_{0,3,\beta}\rangle(T_i,T_j,T_k)\mathbf q^\beta T^k,
\]
其中 $t_1,\ldots,t_r$ 是 $H^2(V,\mathbb C)$ 的变量。由\hyperref[subsec:quantum-product-dependence]{第~\ref*{subsec:quantum-product-dependence} 小节}，在代换 $\mathbf q^\beta\mapsto q^\beta$ 下，这成为
\begin{equation}\label{eq:restricted-dubrovin-connection}
\nabla^\lambda_{\frac{\partial}{\partial t_i}}(T_j)=\lambda T_i*_{\mathrm{small}}T_j,
\end{equation}
前提是，像式~\eqref{eq:quantum-coordinate-rescaling} 那样，以
\[
\omega=\frac1{2\pi i}\sum_{i=1}^rt_iT_i
\]
定义 $q^\beta$。
\SourcePage{268}
A 模型联络用变量 $u_1,\ldots,u_r$ 定义；将上式与式~\eqref{eq:a-model-kahler-coordinates} 比较，得到 $u_j=(1/2\pi i)t_j$。用复化 Kähler 空间上的坐标 $q_1,\ldots,q_r$ 表示，就是
\[
q_j=\exp(2\pi iu_j)=\exp(t_j),
\]
因而
\begin{equation}\label{eq:a-model-coordinate-derivatives}
\frac{\partial}{\partial t_j}=q_j\frac{\partial}{\partial q_j},
\qquad
\frac{\partial}{\partial u_j}=2\pi iq_j\frac{\partial}{\partial q_j}
=2\pi i\frac{\partial}{\partial t_j}.
\end{equation}
于是，式~\eqref{eq:a-model-connection} 和~\eqref{eq:restricted-dubrovin-connection} 立即给出式~\eqref{eq:a-model-dubrovin-identification}，命题得证。
\end{proof}
第 \hyperref[chap:quantum-differential-equations]{\ref*{chap:quantum-differential-equations}}、\hyperref[chap:mirror-theorem]{\ref*{chap:mirror-theorem}} 章中，将使用 $t_1,\ldots,t_r$ 而非定义 A 模型联络时所用的 $u_1,\ldots,u_r$，因为这样更方便。不过，本章余下部分仍将使用 $u_1,\ldots,u_r$，因为它们与\hyperref[chap:moduli]{第~\ref*{chap:moduli} 章}的讨论联系更直接。

此外，\hyperref[chap:quantum-differential-equations]{第~\ref*{chap:quantum-differential-equations} 章}将遇到 $\nabla^\lambda$ 的一种扭曲形式，即 \emph{Givental 联络} $\nabla^g$，定义为\index{Givental联络@Givental 联络|hyperpage}
\[
\nabla^g_{\frac{\partial}{\partial t_i}}\left(\sum_{j=0}^ma_jT_j\right)
=\hbar\sum_{j=0}^m\frac{\partial a_j}{\partial t_i}T_j
-\sum_{j=0}^ma_jT_i*_{\mathrm{small}}T_j,
\]
其中 $\hbar$ 为参数。目前可将 $\hbar$ 看作非零复数，不过，\hyperref[subsec:equivariant-giventals-connection]{第~\ref*{subsec:equivariant-giventals-connection} 小节}将说明，$\hbar$ 在 $\mathbb C^*$ 等变上同调中有内在含义。用 Dubrovin 联络 $\nabla^\lambda$ 表示，则有
\[
\nabla^g=\hbar\nabla^\lambda,\qquad\lambda=-\hbar^{-1}.
\]
严格地说，由于因子 $\hbar$，$\nabla^g$ 不是联络，但讨论其平坦截面仍有意义；这些截面就是 $\nabla^{-\hbar^{-1}}$ 的平坦截面。

\hypertarget{source:268-batch6}{}
\subsection{A 模型联络的渐近性质}\label{subsec:a-model-asymptotics}\index{单值性!A 模型|hyperpage}
定义 A 模型联络 $\nabla$ 时，我们使用了位于 $V$ 的 Kähler 锥闭包内的整基 $T_1,\ldots,T_r$。为研究 $\nabla$ 的渐近性质，令 $\sigma$ 为这些类生成的单纯锥。容易看出 $\operatorname{Int}(\sigma)\subset K(V)$，因此
\[
D_\sigma=\bigl(H^2(V,\mathbb R)+i\operatorname{Int}(\sigma)\bigr)/\operatorname{im}H^2(V,\mathbb Z)\subset K_{\mathbb C}(V)
\]
是 $K_{\mathbb C}(V)$ 的开子集。

如\hyperref[subsec:large-radius-limits]{第~\ref*{subsec:large-radius-limits} 小节}所述，映射
\[
u_1T_1+\cdots+u_rT_r\longmapsto(q_1,\ldots,q_r)=(e^{2\pi iu_1},\ldots,e^{2\pi iu_r})
\]
诱导双全纯同构 $D_\sigma\simeq(\Delta^*)^r$。利用包含关系 $(\Delta^*)^r\subset\Delta^r$，按\hyperref[chap:moduli]{第~\ref*{chap:moduli} 章}的术语，可以把 $0\in\Delta^r$ 看作大半径极限点。此外，假设\hyperref[subsec:large-radius-limits]{第~\ref*{subsec:large-radius-limits} 小节}的\hyperref[conj:kahler-cone]{锥猜想}成立，那么这一大半径极限点自然位于 $V$ 的 Kähler 模空间某个光滑紧化的边界上。

对于小量子乘积，我们希望作如下收敛假设：只要 $\omega$ 在上述映射下的像位于 $0\in\Delta^r$ 的某个固定邻域 $U$ 中，式~\eqref{eq:a-model-connection-expanded} 的级数就收敛。因此，A 模型联络定义在 $U\cap(\Delta^*)^r$ 上。为使叙述简便，以下假设 $\nabla$ 定义在整个 $(\Delta^*)^r$ 上。
\SourcePage{269}
现在应用\hyperref[chap:hodge-yukawa]{第~\ref*{chap:hodge-yukawa} 章}的方法，研究 A 模型联络的单值性。在 $D_\sigma$ 的部分紧化 $\Delta^r$ 中，补集 $\Delta^r-D_\sigma=\Delta^r-(\Delta^*)^r$ 是正规交叉除子 $\bigcup_{j=1}^rD_j$。令 $\mathcal T_j$ 表示绕 $D_j$ 的单值性，它由围绕 $(\Delta^*)^r$ 第 $j$ 个因子的原点一周的回路确定。这些单值性变换可按如下方式计算。
\begin{proposition}\label{prop:a-model-regular-monodromy}\index{正则奇点|hyperpage}
Calabi--Yau 流形 $V$ 的 A 模型联络 $\nabla$ 沿 $\Delta^r-D_\sigma$ 具有正则奇点。此外，单值性变换 $\mathcal T_j$ 幺幂；在相差共轭的意义下，它的对数 $N_j=\log\mathcal T_j$ 由与 $-T_j$ 作杯积给出。
\end{proposition}
\begin{proof}
由式~\eqref{eq:a-model-coordinate-derivatives}，$\Delta^r$ 的坐标 $q_j=e^{2\pi iu_j}$ 满足
\[
\frac{\partial}{\partial u_j}=2\pi iq_j\frac{\partial}{\partial q_j}.
\]
于是，式~\eqref{eq:a-model-connection-expanded} 和 $T^l=\sum_sg^{ls}T_s$ 给出
\begin{equation}\label{eq:a-model-log-poles}
\nabla_{\frac{\partial}{\partial q_j}}(T_k)
=\frac1{2\pi iq_j}\sum_{l,s}\sum_{\beta\in H_2(V,\mathbb Z)}
\langle I_{0,3,\beta}\rangle(T_j,T_k,T_l)q^\beta g^{ls}T_s.
\end{equation}
因为 $\beta$ 有效，而且 $T_1,\ldots,T_r$ 位于 Kähler 锥的闭包内，
\[
q^\beta=q_1^{\int_\beta T_1}\cdots q_r^{\int_\beta T_r}
\]
的指数均非负。实际上，在固定 $j$ 的联络矩阵中，只要 $\beta\ne0$，就可以假设 $q_j$ 的指数为正。为说明这一点，设 $\int_\beta T_j=0$。因为 $\beta\ne0$，除子公理给出
\[
\langle I_{0,3,\beta}\rangle(T_j,T_k,T_l)
=\left(\int_\beta T_j\right)\langle I_{0,2,\beta}\rangle(T_k,T_l)=0.
\]
因此，可以把式~\eqref{eq:a-model-log-poles} 的右端看作一个级数，其中每个非零有效 $\beta$ 对应的非零项，在所固定的 $q_j$ 中具有正指数。

由式~\eqref{eq:a-model-log-poles}，$\nabla$ 相对于基 $T_0,\ldots,T_m$ 的联络矩阵，沿 $\Delta^r-D_\sigma$ 至多有对数极点。由 \SourceCite[定理 4.1]{Deligne1}，这些点为正则奇点。

现在固定 $1\le j\le r$。为计算单值性 $\mathcal T_j$，在 $\Delta^r$ 中取由 $q_i=\text{常数}$（$i\ne j$）给出的“切片”。这样就可把 $\nabla$ 看作以 $q_j$ 为坐标的 $\Delta^*$ 上的联络；其联络矩阵由固定 $j$ 后的式~\eqref{eq:a-model-log-poles} 给出。为计算该矩阵在 $q_j=0$ 处的留数，注意，若 $\beta\ne0$ 有效且相应项非零，则
\begin{equation}\label{eq:a-model-positive-vanishing}
q^\beta\longrightarrow0\qquad\text{当 }q_j\longrightarrow0,
\end{equation}
因为 $q^\beta$ 在 $q_j$ 中具有正指数。

于是，由式~\eqref{eq:a-model-positive-vanishing} 和~\eqref{eq:a-model-log-poles}，$\nabla_{\partial/\partial q_j}$ 在 $q_j=0$ 处的留数矩阵为
\[
\frac1{2\pi i}\left(\sum_l\langle I_{0,3,0}\rangle(T_j,T_k,T_l)g^{ls}\right)_{k,s}
=\frac1{2\pi i}\left(\sum_lg^{ls}\int_VT_j\cup T_k\cup T_l\right)_{k,s}.
\]
将此矩阵记为 $\operatorname{Res}_{q_j=0}(\nabla)$。由 $g^{ls}$ 的定义，容易看出，这就是与 $T_j$ 作杯积的线性映射所对应的矩阵的 $1/(2\pi i)$ 倍。

为完成证明，使用\hyperref[subsec:monodromy-pf]{第~\ref*{subsec:monodromy-pf} 小节}中式~\eqref{eq:connection-residue} 后第二组分项列出的事实。关键在于，与 $T_j$ 作杯积是幂零算子，故
\SourcePage{270}
$\operatorname{Res}_{q_j=0}(\nabla)$ 也幂零。其唯一特征值因而为 $0$，这意味着 $\mathcal T_j$ 的矩阵与 $\exp(-2\pi i\operatorname{Res}_{q_j=0}(\nabla))$ 共轭。因此，$\mathcal T_j$ 幺幂，它的对数 $N_j$ 与以 $-T_j$ 作杯积的算子共轭。
\end{proof}
尽管 $N_j$ 与以 $-T_j$ 作杯积的算子共轭，解释这一说法时仍须谨慎。理解 $\mathcal H$ 的典范延拓以后，情形就会更清楚。

因为 $\nabla$ 具有\IndexLocation{indexloc:1225:270}正则奇点，$(\Delta^*)^r$ 上的丛 $\mathcal H$ 在 $\Delta^r$ 上有典范延拓 $\overline{\mathcal H}$。理解 $\overline{\mathcal H}$ 的关键思想如下：若 $s$ 是 $\mathcal H$ 的一个 $\nabla$ 平坦多值截面，则 $\widetilde s=\exp(-\sum_ju_jN_j)s$ 是 $\mathcal H$ 的单值截面，并自然延拓为 $\overline{\mathcal H}$ 的截面 \SourceCite{CaK}。因而可用这些 $\widetilde s$ 在 $0\in\Delta^r$ 的某个邻域内平凡化 $\overline{\mathcal H}$。（我们的符号与 \SourceCite{CaK} 不同，因为这里使用逆时针的单值性生成元。）

一种常见的表述方式是使用联络 $\nabla^c$，其定义为
\begin{equation}\label{eq:a-model-canonical-connection}
\nabla^c=\nabla+\sum_jN_j\,du_j.
\end{equation}
简单计算表明，$\nabla^c$ 平坦截面恰为 $\exp(-\sum_ju_jN_j)s$，其中 $s$ 为 $\nabla$ 平坦截面。因此，$\nabla^c$ 延拓为 $\overline{\mathcal H}$ 上的联络；其平坦截面在边界点 $0\in\Delta^r$ 的某个邻域内平凡化典范延拓。

上同调基 $T_k$ 可以看作 $(\Delta^*)^r$ 上 $\mathcal H$ 的截面，并将此丛平凡化。接下来要把 $T_k$ 与由 $\nabla^c$ 平坦截面给出的平凡化联系起来。
\begin{proposition}\label{prop:a-model-canonical-flat-basis}\index{典范延拓|hyperpage}
对每个 $T_k$，存在唯一的 $\nabla^c$ 平坦截面 $\widetilde s_k$，满足 $\widetilde s_k=T_k+\text{次数更高的项}$ 且 $\widetilde s_k(0)=T_k$。此外，$N_j$ 作用于 $\widetilde s_k$ 的矩阵，等于以 $-T_j$ 与 $T_k$ 作杯积所得的矩阵。
\end{proposition}
\begin{proof}
先证明 $\dim V=3$ 的情形。为简便起见，仅考虑偶数次上同调，其基可写为
\[
T_0=1,T_1,\ldots,T_r,T^1,\ldots,T^r,T^0,
\]
其中 $T^j\in H^4(V,\mathbb Q)$ 与 $T_j\in H^2(V,\mathbb Q)$ 对偶（$1\le j\le r$），而 $T^0\in H^6(V,\mathbb Q)$ 与 $T_0$ 对偶。这些类都是整类。

先用 Gromov--Witten 势描述 $\nabla$ 的一组平坦截面基。\hyperref[subsec:quintic-gw-potential]{第~\ref*{subsec:quintic-gw-potential} 小节}末尾给出了 $\Phi$ 的一个简洁公式，在这里可写为
\begin{equation}\label{eq:a-model-potential}
\begin{aligned}
\Phi&=\frac{(2\pi i)^3}{6}\int_V\left(\sum_{j=1}^ru_jT_j\right)^3
+\sum_{\beta\ne0}\langle I_{0,0,\beta}\rangle q^\beta\\
&=\frac{(2\pi i)^3}{6}\int_V\left(\sum_{j=1}^ru_jT_j\right)^3+\Phi_{\mathrm{hol}},
\end{aligned}
\end{equation}
其中 $q^\beta=q_1^{\int_\beta T_1}\cdots q_r^{\int_\beta T_r}$，$q_j=\exp(2\pi iu_j)$。与命题~\ref{prop:a-model-regular-monodromy} 的证明一样，求和可限制在非零有效 $\beta$ 上，它们满足所有 $\int_\beta T_j\ge0$；因为 $\operatorname{Int}(\sigma)$ 包含 Kähler 类，这些指数不可能全为零。因此，$\Phi_{\mathrm{hol}}$ 关于 $q_j$ 全纯，并在 $0\in\Delta^r$ 处为零。
\SourcePage{271}
用 $\Phi$ 表示时，容易验证 A 模型联络由下列公式给出：
\begin{equation}\label{eq:a-model-threefold-connection}
\begin{aligned}
\nabla_{\frac{\partial}{\partial u_i}}(T^0)&=0,\\
\nabla_{\frac{\partial}{\partial u_i}}(T^k)&=\delta_{ik}T^0,&&1\le k\le r,\\
\nabla_{\frac{\partial}{\partial u_i}}(T_k)&=\frac1{(2\pi i)^3}\sum_l\frac{\partial^3\Phi}{\partial u_i\partial u_k\partial u_l}T^l,&&1\le k\le r,\\
\nabla_{\frac{\partial}{\partial u_i}}(T_0)&=T_i,
\end{aligned}
\end{equation}
其中用到了式~\eqref{eq:a-model-coordinate-derivatives}。直接计算可知，$\mathcal H$ 的一组 $\nabla$ 平坦基为
\begin{equation}\label{eq:a-model-flat-sections-threefold}
\begin{aligned}
s^0&=T^0,\\
s^k&=T^k-u_kT^0,&&1\le k\le r,\\
s_k&=T_k-\frac1{(2\pi i)^3}\sum_l\frac{\partial^2\Phi}{\partial u_k\partial u_l}T^l
+\frac1{(2\pi i)^3}\frac{\partial\Phi}{\partial u_k}T^0,&&1\le k\le r,\\
s_0&=T_0-\sum_lu_lT_l
-\frac1{(2\pi i)^3}\sum_l\left(\frac{\partial\Phi}{\partial u_l}-\sum_ju_j\frac{\partial^2\Phi}{\partial u_j\partial u_l}\right)T^l\\
&\quad+\frac1{(2\pi i)^3}\left(2\Phi-\sum_ju_j\frac{\partial\Phi}{\partial u_j}\right)T^0.
\end{aligned}
\end{equation}
这些截面在 $(\Delta^*)^r$ 上多值，因为 $u_j$ 出现在 $\Phi$ 的“拓扑”部分 $\frac{(2\pi i)^3}{6}\int_V(\sum_{j=1}^ru_jT_j)^3$ 中。

下一步计算单值性。基本观察是，沿 $(\Delta^*)^r$ 的第 $j$ 个切片绕一周，会把 $u_j$ 变成 $u_j+1$。将此代入上述平坦截面，并把 $\Phi$ 的拓扑部分的偏导数按 $u_j$ 的幂展开，即可计算 $\mathcal T_j$ 的作用。取对数得到
\begin{equation}\label{eq:a-model-monodromy-flat-basis}
\begin{aligned}
N_j(s^0)&=0,\\
N_j(s^k)&=-\delta_{jk}s^0,&&1\le k\le r,\\
N_j(s_k)&=-\sum_l\left.\frac{\partial^3(\Phi/(2\pi i)^3)}{\partial u_j\partial u_k\partial u_l}\right|_{q=0}s^l\\
&=-\sum_l\left(\int_VT_j\cup T_k\cup T_l\right)s^l,&&1\le k\le r,\\
N_j(s_0)&=-s_j.
\end{aligned}
\end{equation}
这些公式中的系数，显然恰为以 $-T_j$ 与给定上同调基作杯积所得矩阵的各项。

知道 $N_j$ 对 $\nabla$ 平坦截面的作用以后，便可计算施加 $\exp(-\sum_ju_jN_j)$ 的结果。这样得到下列 $\nabla^c$ 平坦截面：
\begin{equation}\label{eq:a-model-canonical-sections-threefold}
\begin{aligned}
\widetilde s^0&=T^0,\\
\widetilde s^k&=T^k,&&1\le k\le r,\\
\widetilde s_k&=T_k-\frac1{(2\pi i)^3}\sum_l\frac{\partial^2\Phi_{\mathrm{hol}}}{\partial u_k\partial u_l}T^l
+\frac1{(2\pi i)^3}\frac{\partial\Phi_{\mathrm{hol}}}{\partial u_k}T^0,&&1\le k\le r,\\
\widetilde s_0&=T_0-\frac1{(2\pi i)^3}\sum_l\frac{\partial\Phi_{\mathrm{hol}}}{\partial u_l}T^l
+\frac2{(2\pi i)^3}\Phi_{\mathrm{hol}}T^0.
\end{aligned}
\end{equation}
\SourcePage{272}
这些 $\nabla^c$ 平坦截面显然具有所需形式。在 $0\in\Delta^r$ 处也具有所需的值，因为 $\Phi_{\mathrm{hol}}$ 及其关于 $\partial/\partial u_k=2\pi iq_k\partial/\partial q_k$ 的各阶导数，都在 $0$ 处消失。最后，算子 $\exp(-\sum_ju_jN_j)$ 与 $N_j$ 交换，所以 $N_j$ 相对于式~\eqref{eq:a-model-flat-sections-threefold} 和~\eqref{eq:a-model-canonical-sections-threefold} 的矩阵相同。由式~\eqref{eq:a-model-monodromy-flat-basis}，$N_j$ 对式~\eqref{eq:a-model-canonical-sections-threefold} 的作用具有所需矩阵。

这就证明了 Calabi--Yau 三维簇偶数次上同调这一特殊情形。当 $V$ 的维数为任意 $d\ge3$ 时，需要使用\hyperref[chap:quantum-differential-equations]{第十章}\hyperref[subsec:a-model-gravitational-flat-sections]{第~\ref*{subsec:a-model-gravitational-flat-sections} 小节}对 $\nabla$ 平坦截面的描述。完整证明将在推论~\ref{cor:a-model-canonical-section-triangular} 中给出，故暂时推迟。
\end{proof}
此命题表明，在 $\Delta^r$ 上，可自然地将典范延拓 $\overline{\mathcal H}$ 等同于平凡丛 $H^*(V,\mathbb C)\times\Delta^r$。特别地，$0\in\Delta^r$ 上的纤维是 $H^*(V,\mathbb C)$。利用这一点，现在可以说明单值性对数 $N_j$ 的作用。由命题~\ref{prop:a-model-canonical-flat-basis}，$N_j$ 对 $\nabla^c$ 平坦基 $\widetilde s_k$ 的作用，与以 $-T_j$ 与 $T_k$ 作杯积相同。因为 $\widetilde s_k(0)=T_k$，在 $0$ 上的纤维中就有
\begin{equation}\label{eq:a-model-monodromy-boundary-cup}
N_j(\alpha)=-T_j\cup\alpha,\qquad\alpha\in H^*(V,\mathbb C).
\end{equation}
不过，离开 $0\in\Delta^r$ 后，此式可能不再成立。例如，当 $\dim V=3$ 时，命题~\ref{prop:a-model-canonical-flat-basis} 证明中的公式意味着，若把 $T_k$（$1\le k\le r$）看作 $\overline{\mathcal H}$ 的截面，则
\[
N_j(T_k)=-T_j\cup T_k-\frac1{(2\pi i)^3}\frac{\partial^2\Phi_{\mathrm{hol}}}{\partial u_j\partial u_k}T^0.
\]
在 $\Delta^r$ 的大多数点上，此式与 $-T_j\cup T_k$ 不同。

也可用命题~\ref{prop:a-model-canonical-flat-basis} 给 $\mathcal H$ 赋予整数结构。先取 $H^*(V,\mathbb C)$ 上的自然整数结构，并将其看作 $0\in\Delta^r$ 上的纤维中的整数结构。若 $\widetilde s=\exp(-\sum_ju_jN_j)s$ 在 $0$ 处的值是 $H^*(V,\mathbb C)$ 中的整类，就称 $\nabla$ 平坦截面 $s$ 为整截面。这些整平坦截面在 $(\Delta^*)^r$ 上形成局部常值层 $\mathcal H_{\mathbb Z}$。与 $\mathbb R$ 作张量积得到实向量空间的局部常值层 $\mathcal H_{\mathbb R}$。特别地，$\mathcal H$ 上有复共轭作用。

一个有趣的结果是，若把 $T_k$ 看作 $\mathcal H$ 的截面，那么在大多数纤维中，它既不是整的，也不是实的。例如，当 $\dim V=3$ 时，由所选上同调基可验证，式~\eqref{eq:a-model-flat-sections-threefold} 定义的 $\nabla$ 平坦截面都是整的。但这些公式也给出，例如
\[
T^k=s^k+u_ks^0,\qquad1\le k\le r.
\]
虽然 $s^k$ 和 $s^0$ 都整，$T^k$ 通常连实的都不是，更不用说整的。

还要注意，$\mathcal H_{\mathbb Z}$ 由 $\mathcal H$ 在 $0\in\Delta^r$ 上的行为决定。因此，$\mathcal H$ 的整数结构可能依赖大半径极限点的选择。我们猜想不会出现这种情形，即所有这些点都给出同一个 $\mathcal H_{\mathbb Z}$。只要作出保证 $\mathcal H_{\mathbb Z}$ 存在所需的适当收敛假设，证明应当不难。

既然已知单值性，也就可以确定 A 模型联络的单值性权滤过。
\begin{proposition}\label{prop:a-model-weight-filtration}\index{单值性权滤过|hyperpage}
对于维数 $d\ge3$ 的 Calabi--Yau 流形 $V$，$\nabla$ 在 $0\in\Delta^r$ 处的单值性权滤过 $W_\bullet$ 为
\[
W_i=\bigoplus_{j\ge2d-i}H^j(V,\mathbb C).
\]
\end{proposition}
\SourcePage{273}
\begin{proof}
滤过 $W_\bullet$ 由 $N=\sum_{j=1}^ra_jN_j$ 决定，其中 $N_j=\log\mathcal T_j$ 且 $a_j>0$。由式~\eqref{eq:a-model-monodromy-boundary-cup}，$N_j$ 在 $0$ 上的纤维中是与 $-T_j$ 作杯积的算子。因此，除去一个负号，$N$ 就是在此纤维中与 Kähler 类 $\omega=\sum_{j=1}^ra_jT_j$ 作杯积。由强 Lefschetz 定理，
\[
\cup\omega^k:H^{d-k}(V,\mathbb C)\simeq H^{d+k}(V,\mathbb C)
\]
是同构。因此，\SourceCite{Griffiths2} 中计算权滤过的过程，容易给出上述 $W_i$ 的描述。实际上，刻画单值性权滤过的一个关键性质，是同构 $N^k:W_{d+k}/W_{d+k-1}\simeq W_{d-k}/W_{d-k-1}$，这里正是强 Lefschetz 定理。
\end{proof}
$0$ 处的权滤过 $W_\bullet$ 可借助 $\nabla^c$ 传播为丛 $\overline{\mathcal H}$ 的滤过 $\mathcal W_\bullet$。这样得到如下结果。
\begin{corollary}\label{cor:a-model-weight-bundles}
$\nabla$ 在 $\Delta^r$ 上的单值性权滤过 $\mathcal W_\bullet$ 由平凡丛
\[
\mathcal W_i=\bigoplus_{j\ge2d-i}H^j(V,\mathbb C)\times\Delta^r
\]
给出。
\end{corollary}
\begin{proof}
在 $0$ 处，$W_i$ 由满足 $\deg T_k\ge2d-i$ 的 $T_k$ 决定。因此，在 $\Delta^r$ 上，$\mathcal W_i$ 由这些 $T_k$ 所对应的 $\widetilde s_k$ 决定（记号同命题~\ref{prop:a-model-canonical-flat-basis}）。因为 $\widetilde s_k=T_k+\text{次数更高的项}$，推论随即得到。
\end{proof}
此外，利用命题~\ref{prop:a-model-weight-filtration}，还可以按如下方式刻画 $\nabla$ 的单值性。
\begin{corollary}\label{cor:a-model-maximal-unipotent}\index{最大幺幂单值性!A 模型|hyperpage}
A 模型联络 $\nabla$ 在 $0\in\Delta^r$ 处\IndexLocation{indexloc:0010:273}具有最大幺幂单值性。
\end{corollary}
\begin{proof}
已知单值性幺幂，只须检验定义~\ref{def:maximal-unipotent-point} 的条件 (ii) 和 (iii)。因为 $H^1(V,\mathbb C)=0$，上一命题给出 $\dim W_0=\dim W_1=\dim H^{2d}(V,\mathbb C)=1$，以及 $\dim W_2=\dim(H^{2d}(V,\mathbb C)\oplus H^{2d-2}(V,\mathbb C))=1+r$。此外，若 $g_0\in H^{2d}(V,\mathbb C)$ 是一个点的 Poincaré 对偶，而 $g_1,\ldots,g_r\in H^{2d-2}(V,\mathbb C)$ 在杯积下与 $T_1,\ldots,T_r\in H^2(V,\mathbb C)$ 对偶，则在 $0\in\Delta^r$ 上的纤维中，
\[
N_i(g_j)=-T_i\cup g_j=-\left(\int_VT_i\cup g_j\right)g_0=-\delta_{ij}g_0,
\]
其中第一个等号来自式~\eqref{eq:a-model-monodromy-boundary-cup}。矩阵 $(\delta_{ij})$ 显然可逆，故单值性最大幺幂。
\end{proof}
这个推论有助于说明，为什么\hyperref[chap:hodge-yukawa]{第~\ref*{chap:hodge-yukawa} 章}如此强调最大幺幂单值性。在\IndexLocation{indexloc:0830:273}镜像对称下，A 模型联络应对应于描述镜像复结构变分的 Gauss--Manin 联络（精确表述见\hyperref[sec:mirror-symmetry-revisited]{第~\ref*{sec:mirror-symmetry-revisited} 节}）。A 模型联络具有良好的边界点，其单值性最大幺幂，故镜像的 Gauss--Manin 联络也应具有相同性质。

A 模型联络还满足一个与\hyperref[subsec:mirror-map-definition]{第~\ref*{subsec:mirror-map-definition} 小节}所讨论的整性猜想相类似的命题。为说明这一点，注意，在推论~\ref{cor:a-model-maximal-unipotent} 的证明中，若选取 $T_1,\ldots,T_r$ 为 $H^2(V,\mathbb Z)$ 的一组基，则 $g_0,g_1,\ldots,g_r$ 是 $H^{2d}(V,\mathbb Z)\oplus H^{2d-2}(V,\mathbb Z)$ 的一组基。由 $N(g_j)=-\delta_{ij}g_0$，整性猜想随即成立。不难预料，A 模型联络的这一行为也是镜像一侧整性猜想的动机之一。
\SourcePage{274}
关于 A 模型联络的最后一点说明，涉及 $H^*(V)$ 的基 $T_0,\ldots,T_m$。这些 $T_j$ 可以看作典范延拓 $\overline{\mathcal H}$ 的截面，但它们对于 $\nabla$ 和 $\nabla^c$ 都不平坦，而且除 $0\in\Delta^r$ 处以外也不整。不过，若考察它们在权滤过 $\mathcal W_\bullet$ 的分次部分上的行为，就会看到另一种情形。
\begin{corollary}\label{cor:a-model-graded-integral-basis}
若 $[T_j]$ 表示在 $\operatorname{Gr}^{\mathcal W}$ 上诱导的截面，则这些 $[T_j]$ 构成 $\operatorname{Gr}^{\mathcal W}$ 的典范延拓的一组整平坦截面基。
\end{corollary}
\begin{proof}
推论~\ref{cor:a-model-weight-bundles} 和 $\nabla$ 的定义容易推出，$[T_j]$ 对于 $\operatorname{Gr}^{\mathcal W}$ 上的诱导联络平坦。还可验证，单值性在 $\operatorname{Gr}^{\mathcal W}$ 上平凡，故 $\nabla$ 和 $\nabla^c$ 诱导同一个联络。结合 $T_j$ 在 $0\in\Delta^r$ 上的整数性，就知 $[T_j]$ 整。它们显然构成一组基，证明完成。
\end{proof}
此推论有趣之处如下。\hyperref[subsec:normalized-yukawa-connection]{第~\ref*{subsec:normalized-yukawa-connection} 小节}研究了 Calabi--Yau 三维簇族在最大幺幂边界点处的 Hodge 结构变分。回顾命题~\ref{prop:normalized-yukawa-pf} 的证明，会看到我们先取 $\operatorname{Gr}^{\mathcal W}$ 的一组整平坦截面基，再提升为 $\mathcal H$ 的一组基。由推论~\ref{cor:a-model-graded-integral-basis}，在 A 模型情形，这正是 $T_j$ 所起的作用。

最后两个推论表明，A 模型联络与 Calabi--Yau 三维簇族的 Gauss--Manin 联络十分相似。定义 A 模型 Hodge 滤过以后，这种相似性还将进一步加深；这就是接下来要讨论的主题。

\subsection{A 模型 Hodge 滤过}\label{subsec:a-model-hodge-filtration}
接下来要把 A 模型联络组成一个 Hodge 结构变分。为此，需要一个起 Hodge 滤过作用的对象。照例假设 $V$ 是维数 $\dim V=d\ge3$ 的 Calabi--Yau 流形，并令
\begin{equation}\label{eq:a-model-hodge-filtration}\index{Hodge 滤过!A 模型|hyperpage}
F^p=\bigoplus_{a\le d-p}H^{a,b}(V).
\end{equation}
这给出递减滤过 $H^*(V,\mathbb C)=F^0\supset F^1\supset\cdots$。在丛 $\mathcal H=H^*(V,\mathbb C)\times K_{\mathbb C}(V)$ 中，子丛
\[
\mathcal F^p=F^p\times K_{\mathbb C}(V)\subset\mathcal H
\]
确定 $\mathcal F^0=\mathcal H$ 的一个滤过。

我们的基本断言是，滤过 $\mathcal F^\bullet$ 与 A 模型联络 $\nabla$ 组成权为 $d$ 的 Hodge 结构变分，至少在复化 Kähler 空间 $K_{\mathbb C}(V)$ 的某个大半径极限点附近如此。为作出更精确的表述，回顾整类 $T_1,\ldots,T_r$ 生成的锥的内部 $\operatorname{Int}(\sigma)\subset K(V)$。这给出 $(\Delta^*)^r\subset K_{\mathbb C}(V)$，所讨论的大半径极限点就是 $0\in\Delta^r$。

\hyperref[subsec:a-model-asymptotics]{第~\ref*{subsec:a-model-asymptotics} 小节}在 $\mathcal H$ 上定义了整数结构 $\mathcal H_{\mathbb Z}$，从而特别给出了 $\mathcal H$ 上的复共轭作用。于是，证明 $(\mathcal H,\nabla,\mathcal H_{\mathbb Z},\mathcal F^\bullet)$ 是 Hodge 结构变分，归结为证明以下两个性质：
\begin{itemize}
\item 对每个 $p$，$\nabla(\mathcal F^p)\subset\mathcal F^{p-1}\otimes\Omega^1_{K_{\mathbb C}(V)}$。
\item 对每个 $p$，$\mathcal F^p\oplus\overline{\mathcal F^{d-p+1}}=\mathcal H$。
\end{itemize}
\SourcePage{275}
第一个性质容易证明，因为由\hyperref[cor:cy-quantum-hodge]{推论~\ref*{cor:cy-quantum-hodge}}，小量子乘积与 Hodge 分解相容。论证如下。给定 $\alpha\in H^{a,b}(V)$，其中 $a\le d-p$，可把 $\alpha$ 看作 $\mathcal F^p$ 的截面。因为 $T_i\in H^{1,1}(V)$，有
\[
\nabla_{\frac{\partial}{\partial u_i}}(\alpha)=T_i*_{\mathrm{small}}\alpha\in H^{a+1,b+1}(V).
\]
所以 $\nabla_{\partial/\partial u_i}(\alpha)$ 是 $\mathcal F^{p-1}$ 的截面。这就证明了 $\nabla$ 的 Griffiths 横截性。

不过，证明第二个性质 $\mathcal F^p\oplus\overline{\mathcal F^{d-p+1}}=\mathcal H$ 需要更多工作。为此，引入如下配对：若 $\alpha\in H^k(V,\mathbb C)$、$\beta\in H^l(V,\mathbb C)$，则
\begin{equation}\label{eq:a-model-polarizing-form}
S(\alpha,\beta)=\begin{cases}
(-1)^{k(k+1)/2}\displaystyle\int_V\alpha\cup\beta,&k+l=2d,\\
0,&k+l\ne2d.
\end{cases}
\end{equation}
若把 $\alpha$、$\beta$ 看作 $\mathcal H$ 的截面，上式仍有意义，因此可把 $S$ 看作定义在 $\mathcal H$ 上的配对。

配对 $S$ 具有下列有用性质。
\begin{lemma}\label{lem:a-model-pairing-properties}
设 $S$ 为上述 $\mathcal H$ 上的双线性型，则：
\begin{enumerate}[label=(\roman*)]
\item $S$ 的奇偶性为 $(-1)^d$，即 $S(\alpha,\beta)=(-1)^dS(\beta,\alpha)$。
\item $S$ 关于 $\nabla$ 平坦，即对 $\mathcal H$ 的任意截面 $\alpha,\beta$，以及所有 $1\le i\le r$，有
\[
\frac{\partial}{\partial u_i}S(\alpha,\beta)
=S\left(\nabla_{\frac{\partial}{\partial u_i}}\alpha,\beta\right)
+S\left(\alpha,\nabla_{\frac{\partial}{\partial u_i}}\beta\right).
\]
\item 单值性对数 $N_j$ 是 $S$ 的无穷小自同构，即对任意 $\alpha,\beta\in H^*(V,\mathbb C)$，有
\[
S(N_j(\alpha),\beta)+S(\alpha,N_j(\beta))=0.
\]
\item $S$ 在 $\mathcal H_{\mathbb Z}$ 上取整数值。
\item 对所有 $p$，$S(\mathcal F^p,\mathcal F^{d-p+1})=0$。
\end{enumerate}
\end{lemma}
\begin{proof}
对于第一部分，设 $\alpha\in H^k(V,\mathbb C)$、$\beta\in H^l(V,\mathbb C)$ 且 $k+l=2d$。则 $\alpha\cup\beta=(-1)^{kl}\beta\cup\alpha$，由 $S$ 的定义容易得到所需等式 $S(\alpha,\beta)=(-1)^dS(\beta,\alpha)$。

为证明第二部分，可设 $\alpha=T_j$、$\beta=T_l$。用 $\nabla_{\partial/\partial u_i}$ 与 $T_i$ 的量子乘积之间的定义关系，问题归结为证明
\[
S(T_i*_{\mathrm{small}}T_j,T_l)+S(T_j,T_i*_{\mathrm{small}}T_l)=0.
\]
可设 $\deg T_j=k$、$\deg T_l=2d-k-2$，否则上式平凡成立。于是 $\deg(T_i*_{\mathrm{small}}T_j)=k+2$，故
\[
\begin{aligned}
S(T_i*_{\mathrm{small}}T_j,T_l)
&=(-1)^{(k+2)(k+3)/2}\int_V(T_i*_{\mathrm{small}}T_j)\cup T_l\\
&=-(-1)^{k(k+1)/2}\int_V(T_i*_{\mathrm{small}}T_j)*_{\mathrm{small}}T_l,
\end{aligned}
\]
其中第二个等号使用了命题~\ref{prop:quantum-frobenius-pairing} 给出的 $\int_V\alpha*_{\mathrm{small}}\beta=\int_V\alpha\cup\beta$。类似地，
\[
S(T_j,T_i*_{\mathrm{small}}T_l)=(-1)^{k(k+1)/2}\int_VT_j*_{\mathrm{small}}(T_i*_{\mathrm{small}}T_l).
\]
\SourcePage{276}
因为 $T_i$ 的次数为偶数，$T_i*_{\mathrm{small}}T_j=T_j*_{\mathrm{small}}T_i$。结合小量子乘积的结合律，便得到所需恒等式。

对于第三部分，注意 (ii) 意味着对任意 $a,b\in H^*(V,\mathbb C)$，有 $S(\mathcal T_j(a),\mathcal T_j(b))=S(a,b)$，因为 $\mathcal T_j(a)$ 和 $\mathcal T_j(b)$ 分别由 $a$、$b$ 经平行移动得到。幺幂的保双线性型变换的对数，是无穷小的保双线性型变换，所需结论由此得到。

为证明 (iv)，令 $s_1,s_2$ 为整 $\nabla$ 平坦截面。由 (ii)，$S(s_1,s_2)$ 为常数。令 $\widetilde s_i=\exp(-\sum_ju_jN_j)s_i$；由 (iii)，$S(\widetilde s_1,\widetilde s_2)=S(s_1,s_2)$ 也为常数，所以 $S(s_1,s_2)=S(\widetilde s_1(0),\widetilde s_2(0))$。由 $\mathcal H_{\mathbb Z}$ 的定义，$\widetilde s_i(0)$ 是整类，故 $S(s_1,s_2)\in\mathbb Z$。

最后，设 $\alpha\in H^{a,b}(V)$、$a\le d-p$，以及 $\beta\in H^{e,f}(V)$、$e\le p-1=d-(d-p+1)$。则 $\alpha$、$\beta$ 分别给出 $\mathcal F^p$、$\mathcal F^{d-p+1}$ 的截面。$\alpha\cup\beta$ 的 Hodge 型为 $(a+e,b+f)$，而 $a+e\le d-1$，所以它不是 $(d,d)$ 型，$S(\alpha,\beta)=0$ 随即得到。
\end{proof}
证明 $\nabla$ 和 $\mathcal F^\bullet$ 给出 Hodge 结构变分的关键步骤之一，是证明 $S$ 极化某个混合 Hodge 结构。所讨论的结构来自式~\eqref{eq:a-model-hodge-filtration} 定义的滤过 $F^\bullet$，以及命题~\ref{prop:a-model-weight-filtration} 给出的权滤过。容易验证
\begin{equation}\label{eq:a-model-mixed-hodge-pieces}
F^p\cap\overline{F^{k-p}}\cap W_k=H^{d-p,d-k+p}(V).
\end{equation}
因为 $\operatorname{Gr}^W_k=H^{2d-k}(V,\mathbb C)$，可给 $\operatorname{Gr}^W_k$ 赋予权为 $k$ 的纯 Hodge 结构，使 $(\operatorname{Gr}^W_k)^{p,k-p}=H^{d-p,d-k+p}(V)$。因此，$(W_\bullet,F^\bullet)$ 是混合 Hodge 结构。

还要回顾，在 $0\in\Delta^r$ 处，单值性对数 $N_j$ 是与 $-T_j$ 作杯积的算子。因此，$N=\sum_ja_jN_j$（$a_j>0$）是与 $-\omega$ 作杯积的算子，其中 $\omega=\sum_ja_jT_j$ 为 Kähler 类。命题~\ref{prop:a-model-weight-filtration} 的证明已说明，$W_\bullet$ 是 $N$ 的单值性权滤过。于是有 \SourceCite{CKS2} 中的下列观察。
\begin{proposition}\label{prop:a-model-polarized-mixed-hodge}\index{混合 Hodge 结构!极化|hyperpage}
对任意 $N=\sum_ja_jN_j$，若所有 $a_j>0$，则混合 Hodge 结构 $(W_\bullet,F^\bullet)$ 由 $-N$ 和 $S$ 极化。
\end{proposition}
\begin{proof}
先回顾 \SourceCite{CaK} 中 $(W_\bullet,F^\bullet)$ 被 $-N$ 和 $S$ 极化的含义。定义首先要求 $W_\bullet$ 是 $N$ 的单值性权滤过（已在上文说明），$N(F^p)\subset F^{p-1}$（由定义显然成立），且 $N$ 是 $S$ 的无穷小自同构（由引理~\ref{lem:a-model-pairing-properties} 的证明得到）。此外，若 $k\ge d$，“本原部分” $P_k\subset\operatorname{Gr}^W_k$ 定义为
\[
P_k=\ker\bigl(N^{k-d+1}:\operatorname{Gr}^W_k\longrightarrow\operatorname{Gr}^W_{2d-k-2}\bigr).
\]
极化混合 Hodge 结构定义的最后一项要求，$P_k$ 上的纯 Hodge 结构由 $S(\,\cdot\,,(-N)^{k-d}(\,\cdot\,))$ 极化（这里的负号源于我们的 $N$ 是 \SourceCite{CaK} 所用 $N$ 的相反数）。因此，只须证明，对任意非零 $\alpha\in P_k^{p,k-p}$，有
\begin{equation}\label{eq:a-model-primitive-positivity}
S\bigl(i^{p-(k-p)}\alpha,(-N)^{k-d}(\overline\alpha)\bigr)>0.
\end{equation}
关键在于，以上安排使此式成为本原上同调的通常极化性质。为说明这一点，回顾 $\operatorname{Gr}^W_k=H^{2d-k}(V,\mathbb C)$，且 $(\operatorname{Gr}^W_k)^{p,k-p}=H^{d-p,d-k+p}(V)$。
\SourcePage{277}
与命题~\ref{prop:a-model-weight-filtration} 的证明一样，因 $N$ 是与 $-\omega$ 作杯积的算子，容易看出 $P_k\subset\operatorname{Gr}^W_k$ 恰为关于 $\omega$ 的通常本原上同调 $H^{2d-k}_0(V,\mathbb C)$。因此，若 $\alpha\in P_k^{p,k-p}=H^{d-p,d-k+p}_0(V)$ 非零，通常的 Hodge--Riemann 双线性关系~\eqref{eq:hodge-polarization} 给出
\[
(-1)^{(2d-k)(2d-k-1)/2}\int_V\omega^{d-(2d-k)}\wedge i^{(d-p)-(d-k+p)}\alpha\wedge\overline\alpha>0.
\]
利用 $S$ 的定义，经过简单代数运算，此式化为式~\eqref{eq:a-model-primitive-positivity}。
\end{proof}
现在终于可以证明，我们得到了 Hodge 结构变分。
\begin{theorem}\label{thm:a-model-hodge-variation}\index{极化!A 模型|hyperpage}
设 $V$ 为维数 $d\ge3$ 的 Calabi--Yau 流形。则存在 $0\in\Delta^r$ 的邻域 $U$，使 $(\mathcal H,\nabla,\mathcal H_{\mathbb Z},\mathcal F^\bullet)$ 在 $U\cap(\Delta^*)^r$ 上是由双线性型 $S$ 极化的权 $d$ Hodge 结构变分。
\end{theorem}
\begin{proof}
证明思路与定义 $\mathcal H_{\mathbb Z}$ 时相似：当时用 $0\in\Delta^r$ 上的整数结构确定 $\mathcal H$ 的整数结构。这里用命题~\ref{prop:a-model-polarized-mixed-hodge} 在 $0\in\Delta^r$ 上给出的极化混合 Hodge 结构，证明 $(\mathcal H,\nabla,\mathcal H_{\mathbb Z},\mathcal F^\bullet)$ 在 $0$ 附近是 Hodge 结构变分。

证明将使用 \SourceCite{CaK,CKS1} 的强有力工具，并沿用其中的记号。我们感谢 \SourceCite{CaK} 的作者解释了他们的工作与此情形的联系。令 $\check D$ 为分类 $H^*(V,\mathbb C)$ 上滤过 $\widetilde F^\bullet$ 的旗流形，这些滤过对所有 $p$ 都满足 $\dim\widetilde F^p=\sum_{a\le d-p}h^{a,b}(V)$ 和 $S(\widetilde F^p,\widetilde F^{d-p+1})=0$。在 $\check D$ 内有开集 $D$，分类那些在 $H^*(V,\mathbb C)$ 上诱导权 $d$ 极化 Hodge 结构的滤过。

式~\eqref{eq:a-model-hodge-filtration} 定义的 $F^\bullet$ 位于 $\check D$ 中，因为引理~\ref{lem:a-model-pairing-properties} 给出 $S(F^p,F^{d-p+1})=0$。再考虑滤过族 $\exp(-\sum_ju_jN_j)F^\bullet$；因为 $N_j$ 是 $S$ 的无穷小自同构，这个族也位于 $\check D$ 中。按 \SourceCite{CaK} 的定义，只要满足
\begin{itemize}
\item $N_j(F^p)\subset F^{p-1}$；
\item 当 $\operatorname{Im}(u_j)\gg0$ 时，$\exp(-\sum_ju_jN_j)F^\bullet\in D$，
\end{itemize}
这个族就是幂零轨道。

第一条已在滤过 $F^\bullet$ 的定义~\eqref{eq:a-model-hodge-filtration} 后的讨论中证明。对于第二条，令 $N=\sum_ja_jN_j$，其中 $a_j>0$。则 $N$ 给出命题~\ref{prop:a-model-weight-filtration} 的权滤过 $W_\bullet$，而由命题~\ref{prop:a-model-polarized-mixed-hodge}，$(W_\bullet,F^\bullet)$ 被 $-N$ 和 $S$ 极化。由 \SourceCite[命题 (4.66)]{CKS1}，$\exp(-\sum_ju_jN_j)F^\bullet$ 是幂零轨道。（应当说明，\SourceCite{CaK,CKS1} 将幂零轨道写成 $\exp(\sum_ju_jN_j)F^\bullet$，因为他们的 $N_j$ 是我们的 $-N_j$。）

更具体地说，刚定义的幂零轨道给出 $\mathcal H$ 的滤过 $\widetilde{\mathcal F}^\bullet$。\SourceCite[第 74 页]{CaK} 阐述的关键思想如下：把 $H^*(V,\mathbb C)$ 的滤过 $F^\bullet$ 看作位于\IndexLocation{indexloc:0005:277}典范延拓 $\overline{\mathcal H}$ 在 $0\in\Delta^r$ 上的纤维中（这里沿用\hyperref[subsec:a-model-asymptotics]{第~\ref*{subsec:a-model-asymptotics} 小节}的记号）。用 $\overline{\mathcal H}$ 的 $\nabla^c$ 平坦截面传播 $F^p$，得到子丛 $\widetilde{\mathcal F}^p$，其中 $\nabla^c$ 由式~\eqref{eq:a-model-canonical-connection} 定义。换言之，$\nabla^c$ 平坦截面 $\widetilde s$ 属于 $\widetilde{\mathcal F}^p$，当且仅当 $\widetilde s(0)\in F^p$。

因此，在 $(\Delta^*)^r$ 中充分靠近 $0$ 的地方，$(\mathcal H,\nabla,\mathcal H_{\mathbb Z},\widetilde{\mathcal F}^\bullet)$ 是来自幂零轨道的 Hodge 结构变分。另一方面，滤过 $\mathcal F^\bullet$ 满足 Griffiths 横截性，并在 $0\in\Delta^r$ 处与 $\widetilde{\mathcal F}^\bullet$ 一致。正如 Deligne 在 \SourceCite[第 2.3 节]{Deligne2} 中指出的，这些条件与 \SourceCite[定理 (2.8)]{CaK} 一起，意味着 $(\mathcal H,\nabla,\mathcal H_{\mathbb Z},\mathcal F^\bullet)$ 在 $(\Delta^*)^r$ 中 $0$ 的某个邻域内是 Hodge 结构变分。定理得证。
\end{proof}
\SourcePage{278}
按 \SourceCite{Morrison7} 的术语，称 $(\mathcal H,\nabla,\mathcal H_{\mathbb Z},\mathcal F^\bullet)$ 为 \emph{A 模型 Hodge 结构变分}。由我们的构造方式，直接可证，A 模型变分的极限 Hodge 结构恰为命题~\ref{prop:a-model-polarized-mixed-hodge} 中的混合 Hodge 结构 $(W_\bullet,F^\bullet)$。\index{A 模型 Hodge 结构变分|hyperpage}

为研究镜像对称，需要考虑 A 模型 Hodge 结构变分的一个有趣子变分。令 $\mathcal H^{\mathrm{middle}}$ 为 $\bigoplus_{p=0}^dH^{p,p}(V)$ 给出的丛；这是 $H^*(V,\mathbb C)$ 的 Hodge 分解的中间部分。由\hyperref[cor:cy-quantum-hodge]{推论~\ref*{cor:cy-quantum-hodge}}，小量子乘积与 Hodge 分解相容，故 $\nabla$ 保持 $\mathcal H^{\mathrm{middle}}$。将 $\nabla$ 在其上的限制记为 $\nabla^{\mathrm{middle}}$。还要注意，$\bigoplus_{p=0}^dH^{p,p}(V)$ 具有自然的实结构；通过\hyperref[subsec:a-model-asymptotics]{第~\ref*{subsec:a-model-asymptotics} 小节}的过程，得到实结构 $\mathcal H^{\mathrm{middle}}_{\mathbb R}$。最后，略微滥用记号，仍用 $\mathcal F^\bullet$ 表示 A 模型变分在 $\mathcal H^{\mathrm{middle}}$ 上诱导的滤过。

由这些定义以及定理~\ref{thm:a-model-hodge-variation}，容易得到实 Hodge 结构变分
\begin{equation}\label{eq:middle-a-model-variation}
(\mathcal H^{\mathrm{middle}},\nabla^{\mathrm{middle}},\mathcal H^{\mathrm{middle}}_{\mathbb R},\mathcal F^\bullet).
\end{equation}
这是\emph{中间 A 模型 Hodge 结构变分}，不过通常也将式~\eqref{eq:middle-a-model-variation} 简称为 A 模型变分。上述所有结果都适用于它。特别地，它由交配对
\begin{equation}\label{eq:middle-a-model-polarization}
S(\alpha,\beta)=(-1)^p\int_V\alpha\cup\beta,
\qquad\alpha\in H^{p,p}(V),\quad\beta\in H^{d-p,d-p}(V)
\end{equation}
极化。这是将式~\eqref{eq:a-model-polarizing-form} 限制到 $\bigoplus_{p=0}^dH^{p,p}(V)$ 的结果。还应说明，当 $d=3$ 时，$\mathcal H^{\mathrm{middle}}$ 具有自然的整数结构，因为这时 $H^{1,1}(V)=H^2(V,\mathbb C)$、$H^{2,2}(V)=H^4(V,\mathbb C)$。因此，三维簇情形可将式~\eqref{eq:middle-a-model-variation} 替换为 $(\mathcal H^{\mathrm{middle}},\nabla^{\mathrm{middle}},\mathcal H^{\mathrm{middle}}_{\mathbb Z},\mathcal F^\bullet)$。

还可描述式~\eqref{eq:middle-a-model-variation} 在 $0\in\Delta^r$ 处的行为。由推论~\ref{cor:a-model-maximal-unipotent}，这里的单值性最大幺幂；$0\in\Delta^r$ 上的极限混合 Hodge 结构为 $(W_\bullet,F^\bullet)$，其中
\[
W_k=\bigoplus_{2j\ge2d-k}H^{j,j}(V),\qquad
F^p=\bigoplus_{j\le d-p}H^{j,j}(V).
\]
这是分别限制命题~\ref{prop:a-model-weight-filtration} 和式~\eqref{eq:a-model-hodge-filtration} 对 $W_\bullet$、$F^\bullet$ 的描述所得的结果。此外，这些公式意味着，对所有 $i$，有 $W_{2i}=W_{2i+1}$；对所有 $p$，有 $F^p\oplus W_{2p-2}=\bigoplus_{j=0}^dH^{j,j}(V)$。它们恰为混合 Hodge 结构是 Hodge--Tate 结构的条件~\eqref{eq:mum-hodge-tate-splitting}。因此得到如下命题。
\begin{proposition}\label{prop:a-model-hodge-tate}\index{Hodge--Tate结构@Hodge--Tate 结构|hyperpage}
变分 $(\mathcal H^{\mathrm{middle}},\nabla^{\mathrm{middle}},\mathcal H^{\mathrm{middle}}_{\mathbb R},\mathcal F^\bullet)$ 的极限混合 Hodge 结构为 Hodge--Tate 结构。
\end{proposition}
本节最后讨论 Calabi--Yau 三维簇的三点函数及其与 A 模型联络的关系。
\SourceExampleNumber{8.5.4}{1}
\begin{example}\label{ex:a-model-yukawa-connection}
设 $V$ 是 Calabi--Yau 三维簇。\hyperref[subsec:small-quantum-definition]{第~\ref*{subsec:small-quantum-definition} 小节}对 $a,b,c\in H^2(V,\mathbb C)$ 定义了三点函数（也称 A \IndexLocation{indexloc:0188:278}模型\IndexLocation{indexloc:0195:278}关联函数）$\langle a,b,c\rangle$，并证明它可表示为
\[
\langle a,b,c\rangle=\int_Va*_{\mathrm{small}}b*_{\mathrm{small}}c
\]
\SourcePage{279}
（见式~\eqref{eq:a-three-point-function} 和~\eqref{eq:three-point-quantum-integral}）。我们将从两个方面讨论 $\langle a,b,c\rangle$ 与 A 模型联络 $\nabla^{\mathrm{middle}}$ 的关系，以下将后者简记为 $\nabla$。

先说明如何用 A 模型联络表示 $\langle a,b,c\rangle$。关键观察是，由 $\nabla$ 的定义，
\[
\nabla_{\frac{\partial}{\partial u_i}}(T_j*_{\mathrm{small}}T_k)
=T_i*_{\mathrm{small}}T_j*_{\mathrm{small}}T_k
+H^*(V,\mathbb C)\text{ 中次数较低的项}.
\]
因为 $1=T_0$，容易得到
\[
\int_V\nabla_{\frac{\partial}{\partial u_i}}\nabla_{\frac{\partial}{\partial u_j}}\nabla_{\frac{\partial}{\partial u_k}}1
=\int_VT_i*_{\mathrm{small}}T_j*_{\mathrm{small}}T_k=\langle T_i,T_j,T_k\rangle.
\]
此式可写成
\begin{equation}\label{eq:a-model-yukawa-derivatives}
\langle T_i,T_j,T_k\rangle
=\int_V1\cup\nabla_{\frac{\partial}{\partial u_i}}\nabla_{\frac{\partial}{\partial u_j}}\nabla_{\frac{\partial}{\partial u_k}}1,
\end{equation}
它与\hyperref[sec:yukawa-couplings]{第~\ref*{sec:yukawa-couplings} 节}定义的 B 模型关联函数，也就是 Yukawa 耦合的公式
\[
\int_V\Omega\wedge\nabla_{\theta_1}\nabla_{\theta_2}\nabla_{\theta_3}\Omega
\]
十分相似。在两种情形中，关联函数都以完全相同的方式，由联络（A 模型联络或 Gauss--Manin 联络）、$\mathcal F^3$ 的一个截面（$1$ 或 $\Omega$），以及交配对构造。很快会看到，这种相似性是镜像对称的重要组成部分。

第二个方面是，$\nabla$ 不仅如上所述确定 $\langle a,b,c\rangle$，反过来也成立。这来自命题~\ref{prop:a-model-canonical-flat-basis} 的证明。用式~\eqref{eq:a-model-potential} 给出的 Gromov--Witten 势 $\Phi$，将式~\eqref{eq:cy-potential-three-derivatives} 替换为
\begin{equation}\label{eq:a-model-yukawa-potential}
\frac1{(2\pi i)^3}\frac{\partial^3\Phi}{\partial u_i\partial u_j\partial u_k}=\langle T_i,T_j,T_k\rangle,
\end{equation}
因而式~\eqref{eq:a-model-threefold-connection} 可写为
\begin{equation}\label{eq:a-model-yukawa-full-connection}
\begin{aligned}
\nabla_{\frac{\partial}{\partial u_i}}(T^0)&=0,\\
\nabla_{\frac{\partial}{\partial u_i}}(T^k)&=\delta_{ik}T^0,&&1\le k\le r,\\
\nabla_{\frac{\partial}{\partial u_i}}(T_k)&=\sum_l\langle T_i,T_k,T_l\rangle T^l,&&1\le k\le r,\\
\nabla_{\frac{\partial}{\partial u_i}}(T_0)&=T_i.
\end{aligned}
\end{equation}
所以 $\nabla$ 完全由 A 模型关联函数 $\langle a,b,c\rangle$ 决定。

一个特殊情形是 $h^2(V)=1$。令 $H=T_1$ 生成 $H^2(V,\mathbb C)$，$C=T^1\in H^4(V,\mathbb C)$ 为对偶生成元，并令 $u=u_1$。推广例题~\ref{ex:cy-small-quantum} 的计算，得到
\[
H*_{\mathrm{small}}H=\langle H,H,H\rangle C.
\]
记三点函数为 $Y(q)=\langle H,H,H\rangle$，则上式成为
\begin{equation}\label{eq:one-parameter-quantum-yukawa}
H*_{\mathrm{small}}H=Y(q)C.
\end{equation}
\SourcePage{280}
由式~\eqref{eq:a-model-yukawa-full-connection}，A 模型联络的矩阵为
\[
\begin{pmatrix}0&0&0&0\\1&0&0&0\\0&Y(q)&0&0\\0&0&1&0\end{pmatrix}.
\]
这个矩阵~\eqref{eq:normalized-yukawa-matrix} 曾出现在命题~\ref{prop:normalized-yukawa-pf} 的证明中，当时研究的是复模空间一维的 Calabi--Yau 三维簇。特别地，该命题说明，在最大幺幂边界点处，只要 $Y$ 是归一化 Yukawa 耦合，并以镜像映射作为局部坐标，Gauss--Manin 联络就具有上述形式。与前面一样，这种相似性并非偶然，而是镜像猜想的一部分。

此例也略微预示了\hyperref[chap:quantum-differential-equations]{第十章}的内容。上述矩阵意味着，$1=T_0$ 满足方程
\[
\nabla_{\frac{\partial}{\partial u}}^2
\left(\frac1{Y(q)}\nabla_{\frac{\partial}{\partial u}}^2\right)1=0.
\]
按\hyperref[sec:quantum-cohomology-relations]{第~\ref*{sec:quantum-cohomology-relations} 节}的语言，这是一条\emph{量子微分方程}。在例题~\ref{ex:cy-quantum-equation-relation} 中将看到，它意味着 $V$ 的小量子上同调环中的关系
\[
H*_{\mathrm{small}}H*_{\mathrm{small}}\left(\frac{H*_{\mathrm{small}}H}{Y(q)}\right)=0.
\]
这个关系是平凡的，因为左端属于 $H^8(V)$，而该群显然为零。另一方面，从上述量子微分方程可提取 $Y(q)$（相差一个常数），所以量子微分方程包含了量子上同调的十分非平凡的信息。\hyperref[chap:quantum-differential-equations]{第十章}将证明，一般的量子微分方程同样给出量子上同调中的关系；还将通过例题说明，其中许多关系非平凡。
\end{example}

\section{镜像猜想}\label{sec:mirror-symmetry-revisited}
\hyperref[chap:quantum-cohomology]{第八章}最后要给出镜像对称的一个数学版本。为简便起见，主要考虑局部版本。其基本思想是，在复模空间与 Kähler 模空间中适当选取的边界点附近，\hyperref[chap:moduli]{第~\ref*{chap:moduli} 章}讨论的镜像映射应当把 $V$ 的 Gauss--Manin 联络变成镜像 $V^\circ$ 的 A 模型联络 $\nabla^{\mathrm{middle}}$。先在三维情形给出其精确含义。

\subsection{Calabi--Yau 三维簇的镜像对称}\label{subsec:threefold-mirror-symmetry}
设 $V$、$V^\circ$ 是光滑 Calabi--Yau 三维簇。与\hyperref[sec:mirror-map]{第~\ref*{sec:mirror-map} 节}一样，还假设 $V$ 的复模空间 $\mathcal M(V)$ 和 $V^\circ$ 的 Kähler 模空间 $\mathcal{KM}(V^\circ)$ 具有适当相容的紧化。也就是说，具有：
\begin{itemize}
\item $V$ 的复模空间的某个光滑紧化 $\overline{\mathcal M}(V)$ 上的最大幺幂边界点 $p_0$。
\item 镜像 $V^\circ$ 的 Kähler 模空间的相容光滑紧化 $\overline{\mathcal{KM}}(V^\circ)$ 上的\IndexLocation{indexloc:0696:280}大半径极限点 $q_0$。
\item 镜像映射~\eqref{eq:mirror-map-local}，将 $p_0\in\overline{\mathcal M}(V)$ 的某个邻域映到 $q_0\in\overline{\mathcal{KM}}(V^\circ)$ 的某个邻域。假设 $V$ 满足\hyperref[subsec:higher-dimensional-mum]{第~\ref*{subsec:higher-dimensional-mum} 小节}的整性猜想；由\hyperref[subsec:mirror-map-definition]{第~\ref*{subsec:mirror-map-definition} 小节}，这就唯一确定镜像映射。
\end{itemize}
\SourcePage{281}
如\hyperref[subsec:mirror-map-definition]{第~\ref*{subsec:mirror-map-definition} 小节}所述，将 $p_0$ 与 $q_0$ 配对以及确定镜像映射时，需要作出若干选择。以下假设这些选择已完成。

这里还有两个极化 Hodge 结构变分。在 $V$ 的复模空间上，由 $H^3(V,\mathbb C)$ 给出的丛 $\mathcal H^V$ 具有 Gauss--Manin 联络及 Hodge 滤过。为避免混淆，将 Gauss--Manin 联络记为 $\nabla^{\mathrm{GM}}$。在镜像一侧，$V^\circ$ 的 Kähler 模空间上，由 $H^{\mathrm{middle}}(V^\circ)=\bigoplus_{p=0}^3H^{p,p}(V^\circ)=\bigoplus_{p=0}^3H^{2p}(V^\circ,\mathbb C)$ 给出的丛 $\mathcal H^{V^\circ}$，具有 A 模型联络 $\nabla^{\mathrm{middle}}$ 及 Hodge 滤过。此外，$\mathcal H^V$ 和 $\mathcal H^{V^\circ}$ 分别有整数结构 $\mathcal H^V_{\mathbb Z}$、$\mathcal H^{V^\circ}_{\mathbb Z}$。最后，还有如下极化：
\begin{itemize}
\item 因为每个 $\alpha\in H^3(V,\mathbb C)$ 都本原，由式~\eqref{eq:hodge-polarization}，双线性型
\[
Q(\alpha,\beta)=(-1)^{3(3-1)/2}\int_V\alpha\cup\beta=-\int_V\alpha\cup\beta
\]
极化 $(\mathcal H^V,\nabla^{\mathrm{GM}},\mathcal H^V_{\mathbb Z},\mathcal F^\bullet)$。
\item 式~\eqref{eq:middle-a-model-polarization} 在 $\bigoplus_{p=0}^3H^{p,p}(V^\circ)$ 上定义了双线性型
\[
S(\alpha,\beta)=(-1)^p\int_{V^\circ}\alpha\cup\beta,\qquad
\alpha\in H^{p,p}(V^\circ),\quad\beta\in H^{3-p,3-p}(V^\circ).
\]
由定理~\ref{thm:a-model-hodge-variation}，$S$ 极化 $(\mathcal H^{V^\circ},\nabla^{\mathrm{middle}},\mathcal H^{V^\circ}_{\mathbb Z},\mathcal F^\bullet)$。
\end{itemize}
除此以外，还有几个特定截面：
\begin{itemize}
\item 对于 $V^\circ$，$\mathcal H^{V^\circ}$ 具有 $W_0$ 的整生成元 $[pt]$，其中 $[pt]\in H^{3,3}(V^\circ,\mathbb Z)$ 为一个点的类。还有特定截面 $1=T_0\in H^0(V^\circ,\mathbb C)$，由式~\eqref{eq:a-model-hodge-filtration}，它生成 $\mathcal F^3$。注意 $S(1,[pt])=1$。
\item 对于 $V$，丛 $\mathcal H^V$ 具有最大幺幂边界点定义中的截面 $g_0$。它是 $W_0$ 的整生成元，且在相差 $\pm1$ 的意义下唯一。此外，还有 $\mathcal F^3$ 的截面 $\Omega$。为得到典范的对象，令 $y_0=\langle g_0,\Omega\rangle$。一旦固定 $g_0$，$\Omega/y_0$ 就是满足 $Q(\Omega/y_0,g_0)=1$ 的唯一 $\mathcal F^3$ 截面。
\end{itemize}
于是终于可以如下精确定义 Calabi--Yau 三维簇的镜像对称。
\begin{definition}\label{def:mathematical-mirror-pair}
给定如上的 Calabi--Yau 三维簇 $V$、$V^\circ$，若镜像映射可提升为 $p_0\in\overline{\mathcal M}(V)$ 和 $q_0\in\overline{\mathcal{KM}}(V^\circ)$ 的邻域上，丛 $\mathcal H^V$ 与 $\mathcal H^{V^\circ}$ 的同构，就称 $(V,V^\circ)$ 为\emph{数学镜像对}。此外，此同构应保持来自 $\nabla^{\mathrm{GM}}$ 和 $\nabla^{\mathrm{middle}}$ 的极化 Hodge 结构变分，并将 $\mathcal H^V$ 的截面 $\Omega/y_0$、$g_0$ 分别映到 $\mathcal H^{V^\circ}$ 的截面 $1$、$[pt]$。
\end{definition}
\index{数学镜像对|hyperpage}\index{镜像对!数学镜像对|hyperpage}
借用 \SourceCite{Morrison7} 的术语，数学镜像对 $(V,V^\circ)$ 也可称为\emph{基于 Hodge 理论的镜像对}。定义~\ref{def:mathematical-mirror-pair} 只是局部表述；利用\hyperref[chap:moduli]{第~\ref*{chap:moduli} 章}讨论的复模空间和 Kähler 模空间，也可以给出更具整体性的定义。

接下来讨论定义~\ref{def:mathematical-mirror-pair}，并给出一些等价表述，随后再讨论镜像猜想的含义。

首先，定义~\ref{def:mathematical-mirror-pair} 假设 $V$ 满足整性猜想。\hyperref[subsec:a-model-asymptotics]{第~\ref*{subsec:a-model-asymptotics} 小节}末尾已指出，$V^\circ$ 上的 A 模型 Hodge 结构变分总满足整性猜想。因此，对 $V$ 的这一假设
\SourcePage{282}
是合理的，因为定义~\ref{def:mathematical-mirror-pair} 中的丛同构保持整数结构。

其次，与定义~\ref{def:mathematical-mirror-pair} 相容的丛同构 $\mathcal H^V\simeq\mathcal H^{V^\circ}$ 至多有一个。为说明这一点，注意，由 \SourceCite{BG}，反复对 $\widetilde\Omega=\Omega/y_0$ 施加 $\nabla^{\mathrm{GM}}_{\delta_i}$ 就能生成 $\mathcal H^V$。镜像一侧的 Lefschetz 分解表明，$H^{\mathrm{middle}}(V^\circ)$ 在杯积下由 $H^{1,1}(V^\circ)$ 生成；至少在 $q_0$ 附近，对小量子乘积也如此。这容易推出，反复对 $T_0=1$ 施加 $\nabla^{\mathrm{middle}}_{\delta_i}$ 就能生成 $\mathcal H^{V^\circ}$。所需的唯一性随即得到。

下面用更具体的方式讨论定义~\ref{def:mathematical-mirror-pair}。可以使用关联函数，也可以使用势函数，这将分别导出下面的定理~\ref{thm:mirror-correlation-equivalence} 和推论~\ref{cor:mirror-potential-equivalence}。

先\IndexLocation{indexloc:0757:282}用关联函数表述数学镜像对的含义。$V$ 上有 B 模型关联函数，也称归一化 Yukawa 耦合。回顾命题~\ref{prop:normalized-yukawa-pf} 后的讨论，在 $p_0$ 处存在局部坐标 $q_1,\ldots,q_r$，使 $\nabla^{\mathrm{GM}}$ 具有特别简单的形式~\eqref{eq:multimoduli-connection}；例题~\ref{ex:mirror-map-integral-connection} 表明，$q_j$ 确定镜像映射。若 $q_j=\exp(2\pi iu_j)$，则由式~\eqref{eq:multimoduli-yukawa}，$V$ 的归一化 Yukawa 耦合为
\begin{equation}\label{eq:b-model-normalized-correlations}\index{Yukawa耦合@Yukawa 耦合!归一化|hyperpage}
\begin{aligned}
Y_{ijk}&=-\int_V\widetilde\Omega\wedge\nabla_{\frac{\partial}{\partial u_i}}\nabla_{\frac{\partial}{\partial u_j}}\nabla_{\frac{\partial}{\partial u_k}}\widetilde\Omega\\
&=Q\left(\widetilde\Omega,\nabla_{\frac{\partial}{\partial u_i}}\nabla_{\frac{\partial}{\partial u_j}}\nabla_{\frac{\partial}{\partial u_k}}\widetilde\Omega\right),
\end{aligned}
\end{equation}
其中 $\widetilde\Omega=\Omega/y_0$ 为归一化三形式，$\nabla=\nabla^{\mathrm{GM}}$ 为 Gauss--Manin 联络，最后一个等号来自极化 $Q$ 的定义。这些就是 $V$ 的 B 模型关联函数。

在镜像一侧，存在 $T_1,\ldots,T_r\in H^2(V^\circ,\mathbb Z)$，使
\[
(q_1,\ldots,q_r)=(e^{2\pi iu_1},\ldots,e^{2\pi iu_r})
\longmapsto\left[\sum_ju_jT_j\right]\in K_{\mathbb C}(V^\circ)
\]
给出大半径极限点 $q_0$ 处的局部坐标。利用这些坐标，式~\eqref{eq:three-point-quantum-integral} 给出 A 模型关联函数
\begin{equation}\label{eq:a-model-mirror-correlations}
\langle T_i,T_j,T_k\rangle=\int_{V^\circ}T_i*_{\mathrm{small}}T_j*_{\mathrm{small}}T_k.
\end{equation}
用 A 模型联络 $\nabla=\nabla^{\mathrm{middle}}$ 表示时，式~\eqref{eq:a-model-yukawa-derivatives} 给出
\begin{equation}\label{eq:a-model-correlations-polarization}
\begin{aligned}
\langle T_i,T_j,T_k\rangle
&=\int_{V^\circ}1\cup\nabla_{\frac{\partial}{\partial u_i}}\nabla_{\frac{\partial}{\partial u_j}}\nabla_{\frac{\partial}{\partial u_k}}1\\
&=S\left(1,\nabla_{\frac{\partial}{\partial u_i}}\nabla_{\frac{\partial}{\partial u_j}}\nabla_{\frac{\partial}{\partial u_k}}1\right),
\end{aligned}
\end{equation}
其中最后一个等号使用了极化 $S$ 的定义。
\begin{theorem}\label{thm:mirror-correlation-equivalence}
若 $V$ 满足整性猜想，且镜像映射如上定义，则 $(V,V^\circ)$ 是定义~\ref{def:mathematical-mirror-pair} 意义下的数学镜像对，当且仅当 $V$ 的 B 模型关联函数 $Y_{ijk}$ 和 $V^\circ$ 的 A 模型关联函数 $\langle T_i,T_j,T_k\rangle$ 满足
\[
Y_{ijk}=\langle T_i,T_j,T_k\rangle\qquad\text{对所有 }i,j,k.
\]
\end{theorem}
\SourcePage{283}
\begin{proof}
一个方向平凡：若有同构将 $\widetilde\Omega$ 映到 $1$、$\nabla^{\mathrm{GM}}$ 映到 $\nabla^{\mathrm{middle}}$、$Q$ 映到 $S$，则关联函数相等由式~\eqref{eq:b-model-normalized-correlations} 和~\eqref{eq:a-model-correlations-polarization} 立即得到。

反过来，设关联函数按所需方式相同。用 $q_j$ 同时作为 $V$ 的复模空间与 $V^\circ$ 的 Kähler 模空间的局部坐标。由命题~\ref{prop:normalized-yukawa-pf} 后的讨论，丛 $\mathcal H^V$ 具有基 $e_0,e_j$（$1\le j\le r$）、$e^j$（$1\le j\le r$）、$e^0$，其中 $e_0=\widetilde\Omega$ 是特定截面。类似地，由命题~\ref{prop:a-model-canonical-flat-basis} 的证明，$\mathcal H^{V^\circ}$ 具有基 $T_0,T_j$（$1\le j\le r$）、$T^j$（$1\le j\le r$）、$T^0$，其中 $T_0=1$ 为特定截面。令 $e_i\mapsto T_i$、$e^i\mapsto T^i$，显然定义了一个丛同构。

为完成证明，须检验联络相容，且极化与整数结构都被保持。关联函数相同的假设，结合式~\eqref{eq:a-model-threefold-connection} 和~\eqref{eq:multimoduli-connection}，立即给出联络相容。极化也同样容易处理：使用双线性型 $Q$、$S$ 的定义，以及式~\eqref{eq:multimoduli-integral-pairings} 中 $e_i,e^i$ 的描述和命题~\ref{prop:a-model-canonical-flat-basis} 中 $T_i,T^i$ 的描述即可。

还需考虑整数结构。这一点稍微微妙，因为上述基不是整的。不过，这些基可延拓到 $\Delta^r$ 上的典范延拓，并在 $0\in\Delta^r$ 上成为整的。\hyperref[subsec:a-model-asymptotics]{第~\ref*{subsec:a-model-asymptotics} 小节}定义 $\mathcal H^{V^\circ}$ 的整数结构时规定：$\nabla^{\mathrm{middle}}$ 平坦截面 $s$ 为整截面，当且仅当 $\exp(-\sum_ju_jN_j)s$ 在 $0$ 上整。对于 $\mathcal H^V$ 也如此，于是容易得到，同构保持整数结构。

\end{proof}
这个定理很好地将数学镜像对的定义与\hyperref[chap:introduction]{第~\ref*{chap:introduction} 章}引入的“经典”镜像对称概念联系起来。这种方法的一个不足是，当 $V$ 的复模空间维数 $r$ 较大时，可能需要考虑很多关联函数，计算会变得相当复杂。使用势函数是一种简化办法。

回顾这一过程。在 $V$ 上，引理~\ref{lem:yukawa-potential} 证明了势函数 $\Phi^{\mathrm{GM}}$ 的存在。如\hyperref[subsec:normalized-yukawa-mirror]{第~\ref*{subsec:normalized-yukawa-mirror} 小节}所述，它是变量 $u_j$ 的函数，其中镜像映射为 $q_j=e^{2\pi iu_j}$。在这些变量中，引理~\ref{lem:yukawa-potential} 说明，归一化 Yukawa 耦合~\eqref{eq:b-model-normalized-correlations} 为
\begin{equation}\label{eq:b-model-potential-correlations}
Y_{ijk}=\frac1{(2\pi i)^3}\frac{\partial^3\Phi^{\mathrm{GM}}}{\partial u_i\partial u_j\partial u_k}
\qquad\text{对所有 }i,j,k.
\end{equation}
转到 $V^\circ$，由式~\eqref{eq:a-model-potential}，Gromov--Witten 势 $\Phi^{\mathrm{GW}}$ 等于
\[
\Phi^{\mathrm{GW}}=\frac{(2\pi i)^3}{6}\int_{V^\circ}\left(\sum_{j=1}^ru_jT_j\right)^3
+\sum_{\beta\ne0}\langle I_{0,0,\beta}\rangle q^\beta,
\]
其中 $q^\beta=q_1^{\int_\beta T_1}\cdots q_r^{\int_\beta T_r}$、$q_j=\exp(2\pi iu_j)$。类似地，式~\eqref{eq:a-model-yukawa-potential} 给出
\[
\langle T_i,T_j,T_k\rangle=\frac1{(2\pi i)^3}\frac{\partial^3\Phi^{\mathrm{GW}}}{\partial u_i\partial u_j\partial u_k}
\qquad\text{对所有 }i,j,k.
\]
因此，定理~\ref{thm:mirror-correlation-equivalence} 立即给出数学镜像对的势函数表述。
\SourcePage{284}
\begin{corollary}\label{cor:mirror-potential-equivalence}\index{镜像对称!势函数|hyperpage}
若 $V$ 满足整性猜想，且镜像映射如上定义，则 $(V,V^\circ)$ 是数学镜像对，当且仅当
\[
\Phi^{\mathrm{GM}}=\Phi^{\mathrm{GW}},
\]
其中允许两边相差关于 $u_j$ 的至多二次项。
\end{corollary}
\hyperref[subsec:quintic-mirror-symmetry]{第~\ref*{subsec:quintic-mirror-symmetry} 小节}将说明定理~\ref{thm:mirror-correlation-equivalence} 和推论~\ref{cor:mirror-potential-equivalence} 如何适用于五次三维簇。

在上述对定义~\ref{def:mathematical-mirror-pair} 及其后果的讨论中，应始终记住：要使 A 模型联络 $\nabla^{\mathrm{middle}}$ 成为真正的联络，需要收敛假设。有两种方式可以避开这一困难。第一种方式是注意，定理~\ref{thm:mirror-correlation-equivalence} 中的等式
\[
Y_{ijk}=\langle T_i,T_j,T_k\rangle
\]
作为关于形式幂级数的陈述完全有意义。因此，若能证明此等式，因为已知 $Y_{ijk}$ 收敛，$\langle T_i,T_j,T_k\rangle$ 的收敛性随即得到。于是\hyperref[sec:a-model-vhs]{第~\ref*{sec:a-model-vhs} 节}的全部理论都适用。类似地，在推论~\ref{cor:mirror-potential-equivalence} 的情形，只须证明形式幂级数等式 $\Phi^{\mathrm{GM}}=\Phi^{\mathrm{GW}}$，其余结论便都成立。

另一种可能更有趣的方式，是按 \SourceCite[第 7.2 节]{Morrison7} 的主张，建立“形式退化的 Hodge 结构变分”理论。这要求将\hyperref[sec:a-model-vhs]{第~\ref*{sec:a-model-vhs} 节}的一切发展为形式版本，再将定义~\ref{def:mathematical-mirror-pair} 替换为相应的形式陈述。

既然已有数学镜像对的严格定义，下一步似乎应当给出一个\emph{镜像猜想}，断言每个 Calabi--Yau 三维簇 $V$ 都有一个镜像 $V^\circ$，使 $(V,V^\circ)$ 满足定义~\ref{def:mathematical-mirror-pair}。遗憾的是，这行不通。刚性 Calabi--Yau 三维簇 $V$（只有平凡形变，故 $H^{2,1}(V)=0$）不可能具有该定义意义下的镜像 $V^\circ$，因为这会导致 $H^{1,1}(V^\circ)\simeq H^{2,1}(V)=0$，这是不可能的。因此，最多只能希望一个镜像猜想断言某些 Calabi--Yau 三维簇具有镜像。从这个角度看，\hyperref[chap:mirror-constructions]{第~\ref*{chap:mirror-constructions} 章}的每一种镜像构造，都导出一个镜像猜想。\hyperref[subsec:toric-mirror-conjecture]{第~\ref*{subsec:toric-mirror-conjecture} 小节}将详细探讨 Batyrev 镜像构造中的情形。

以推论~\ref{cor:mirror-potential-equivalence} 提示的一个有趣问题结束数学镜像对的讨论。给定任意函数 $\Phi(u_1,\ldots,u_r)$，使
\[
Y_{ijk}=\frac1{(2\pi i)^3}\frac{\partial^3\Phi}{\partial u_i\partial u_j\partial u_k}
\]
关于 $q_j=\exp(2\pi iu_j)$ 全纯，就可以用式~\eqref{eq:multimoduli-connection} 在某个向量丛上定义联络，并得到显然的 Hodge 滤过。利用\hyperref[subsec:a-model-hodge-filtration]{第~\ref*{subsec:a-model-hodge-filtration} 小节}的方法，可证明这给出 Hodge 结构变分。可能的 $\Phi$ 很多，故以此方式得到的变分也很多。但其中哪些来自几何，也就是来自 Calabi--Yau 三维簇族？换言之，Calabi--Yau 变分的势函数有什么特殊之处？镜像对称给出了部分答案：若 $\Phi$ 来自 Calabi--Yau 三维簇族，就应当是镜像的 Gromov--Witten 势
\SourcePage{285}
$\Phi^{\mathrm{GW}}$。这意味着势函数 $\Phi$ 具有非常特殊的形式。因此，镜像对称也帮助我们理解，哪些抽象的 Hodge 结构变分来自几何这一一般问题。

\subsection{五次三维簇的镜像对称}\label{subsec:quintic-mirror-symmetry}
接下来比较\hyperref[chap:quintic]{第~\ref*{chap:quintic} 章}给出的五次三维簇镜像对称的初步版本，与定义~\ref{def:mathematical-mirror-pair} 中更精细的表述。

\hyperref[chap:quintic]{第~\ref*{chap:quintic} 章}将五次三维簇的镜像对称表述为
\begin{equation}\label{eq:quintic-mirror-correlations-revisited}
5+\sum_{d=1}^\infty n_dd^3\frac{q^d}{1-q^d}
=\frac5{(1+5^5z)y_0(z)^2}\left(\frac qz\frac{dz}{dq}\right)^3,
\end{equation}
其中 $n_d$ 是五次三维簇的瞬子数，$z$ 是五次镜像的模坐标（\hyperref[chap:quintic]{第二章}记为 $x$），使五次镜像族在 $z=0$ 处的单值性最大幺幂。此外，
\[
\begin{aligned}
y_0(z)&=\sum_{n=0}^\infty\frac{(5n)!}{(n!)^5}(-1)^nz^n,\\
y_1(z)&=y_0(z)\log(-z)+5\sum_{n=1}^\infty\frac{(5n)!}{(n!)^5}
\left[\sum_{j=n+1}^{5n}\frac1j\right](-1)^nz^n,\\
q&=\exp(y_1/y_0).
\end{aligned}
\]
于是有如下良好事实。
\begin{theorem}\label{thm:quintic-mirror-correlation-equivalence}
取 $V^\circ$ 为五次三维簇、$V$ 为五次镜像时，定义~\ref{def:mathematical-mirror-pair} 与式~\eqref{eq:quintic-mirror-correlations-revisited} 等价。
\end{theorem}
\begin{proof}
根据现有知识，这容易看出。由式~\eqref{eq:cy-quantum-correlation-instanton}，式~\eqref{eq:quintic-mirror-correlations-revisited} 左端是五次三维簇的 A 模型关联函数 $\langle H,H,H\rangle$；例题~\ref{ex:quintic-normalized-yukawa} 则证明，右端是五次镜像的归一化 Yukawa 耦合。因此，该等式意味着定义 A 模型联络的级数收敛；如前所述，由 \SourceCite{Morrison2}，五次镜像族满足整性猜想。定理遂由定理~\ref{thm:mirror-correlation-equivalence} 得到。
\end{proof}
换言之，一旦证明\hyperref[chap:quintic]{第二章}中的公式，定义~\ref{def:mathematical-mirror-pair} 给出的五次三维簇镜像对称的完整版本，就成为直接推论。不过，不应因刚才的证明简短而低估其中的工作。对式~\eqref{eq:quintic-mirror-correlations-revisited} 两端所作的论断，背后都有大量工作。对于 A 模型，定义 $n_d$ 需要\hyperref[chap:gromov-witten]{第七章}的 Gromov--Witten 不变量定义，而理解其计数意义也并非平凡，见\hyperref[subsec:instanton-numbers]{第~\ref*{subsec:instanton-numbers} 小节}。随后还须完成\hyperref[sec:a-model-vhs]{第~\ref*{sec:a-model-vhs} 节}的 A 模型 Hodge 理论。对于 B 模型，\hyperref[chap:hodge-yukawa]{第五章}研究了镜像族的 Hodge 理论；例题~\ref{ex:quintic-normalized-yukawa} 表明，为证明右端是归一化 Yukawa 耦合，需要非常具体的单值性信息。

下面按推论~\ref{cor:mirror-potential-equivalence} 的方式，说明五次三维簇的镜像对称与势函数的关系。这很重要，因为\hyperref[chap:mirror-theorem]{第十一章}正是以此方式证明五次三维簇的镜像对称。

先描述 A 模型和 B 模型的势函数。它们是 $q$ 的多值函数，不过不再用 $u_1=(1/2\pi i)\log q$ 表示，而改用 $t_1=\log q$。如\hyperref[subsec:quantum-product-dependence]{第~\ref*{subsec:quantum-product-dependence} 小节}所述，变量
\SourcePage{286}
$t_1=2\pi iu_1$ 更适合\hyperref[chap:mirror-theorem]{第十一章}的讨论。对于 A 模型，这意味着将式~\eqref{eq:a-model-potential} 替换为 Gromov--Witten 势
\[
\Phi^{\mathrm{GW}}=\frac56t_1^3+\sum_{d=1}^\infty\langle I_{0,0,d}\rangle q^d.
\]
对于 B 模型，令 $Y$ 为式~\eqref{eq:quintic-mirror-correlations-revisited} 右端的归一化 Yukawa 耦合。若把 $Y$ 看作 $t_1=\log q$ 的函数，则势函数为三阶导数等于 $Y$ 的函数 $\Phi^{\mathrm{GM}}(t_1)$。为描述 $\Phi^{\mathrm{GM}}$，令 $z=-x$。将 $y_0,y_1$ 看作 $x$ 的函数时，它们由式~\eqref{eq:quintic-frobenius-basis} 给出；该式还给出 $y_2,y_3$，使 $y_0,\ldots,y_3$ 构成五次镜像 Picard--Fuchs 方程（以 $x$ 为变量）的一组解基。势函数由如下简洁公式给出。
\begin{proposition}\label{prop:quintic-mirror-potential}
五次镜像族的势函数为
\[
\Phi^{\mathrm{GM}}=\frac52\left(\frac{y_1}{y_0}\frac{y_2}{y_0}-\frac{y_3}{y_0}\right),
\]
其中借助镜像映射 $q=\exp(y_1/y_0)$，将其看作 $t_1=\log q$ 的函数。
\end{proposition}
\begin{proof}
须证明
\[
\frac{d^3}{dt_1^3}\frac52\left(\frac{y_1}{y_0}\frac{y_2}{y_0}-\frac{y_3}{y_0}\right)=Y.
\]
\hyperref[chap:quintic]{第二章}已证明此式，为完整起见，再回顾论证。$y_i$ 满足例题~\ref{ex:quintic-pf-method} 推导的 Picard--Fuchs 方程。将 $y_i$ 替换为 $y_i/y_0$，将 $x$ 替换为 $q=e^{t_1}$，命题~\ref{prop:normalized-yukawa-pf} 表明该方程成为
\[
\frac{d^2}{dt_1^2}\left(\frac{d^2y/dt_1^2}{Y}\right)=0.
\]
还要注意，由镜像映射的定义，$t_1=y_1/y_0$。利用此方程以及 $Y$、$y_i$ 的幂级数展开，式~\eqref{eq:period-derivative-identities} 中已证明
\[
5\frac{d^2}{dt_1^2}\frac{y_2}{y_0}=Y,\qquad
5\frac{d^2}{dt_1^2}\frac{y_3}{y_0}=Yt_1.
\]
于是，式~\eqref{eq:lly-yukawa} 后的论证给出所需等式。
\end{proof}
结合推论~\ref{cor:mirror-potential-equivalence}，得到如下定理。以下在镜像映射的源端用 $t_1=\log x$，在目标端用 $\Psi(t_1)=\log q=y_1/y_0$。
\begin{theorem}\label{thm:quintic-mirror-potential-equivalence}
取 $V^\circ$ 为五次三维簇、$V$ 为五次镜像时，定义~\ref{def:mathematical-mirror-pair} 等价于方程
\[
\Phi^{\mathrm{GW}}(\Psi(t_1))=\frac52\left(\frac{y_1}{y_0}\frac{y_2}{y_0}-\frac{y_3}{y_0}\right),
\]
其中镜像映射为 $\Psi(t_1)=y_1/y_0$，势函数取上述归一化。
\end{theorem}
采用这种形式表述定理~\ref{thm:quintic-mirror-potential-equivalence}，是因为上述方程就是 \SourceCite{LLY} 中的“\IndexLocation{indexloc:0864:286}\IndexLocation{indexloc:0867:286}镜像定理”，也是\hyperref[chap:mirror-theorem]{第十一章}定理~\ref{thm:lly-quintic-mirror} 证明五次三维簇\IndexLocation{indexloc:1197:286}镜像对称所用的形式。

\SourcePage{287}
\subsection{更高维情形}\label{subsec:higher-dimensional-mirror}\index{镜像对称!高维情形|hyperpage}
下面考虑 $V$、$V^\circ$ 是维数 $d>3$ 的光滑 Calabi--Yau 流形的情形。与前面一样，假设有对应的点 $p_0\in\overline{\mathcal M}(V)$、$q_0\in\overline{\mathcal{KM}}(V^\circ)$。困难出现在考察两个极化 Hodge 结构变分时。将会看到，情形不如三维簇那么理想。

先看 $V$。虽然 $H^d(V,\mathbb C)$ 上有良好的 Hodge 结构变分，但极化~\eqref{eq:hodge-polarization} 只适用于本原上同调，而本原上同调不一定是整个 $H^d(V,\mathbb C)$。为处理这一点，固定 Kähler 形式 $\omega$，并考虑 Lefschetz 分解。每个 $\alpha\in H^d(V,\mathbb C)$ 都可写成 $\sum_j\omega^j\cup\alpha_{d-2j}$，其中 $\alpha_{d-2j}\in H^{d-2j}_0(V,\mathbb C)$ 本原。于是有双线性型
\[
Q(\alpha,\beta)=\sum_j(-1)^{d(d-1)/2}(-1)^j
\int_V\omega^{2j}\cup\alpha_{d-2j}\cup\beta_{d-2j}.
\]
一般而言，此型只在 $\mathbb R$ 上定义（因为 $\omega$ 未必整），但由式~\eqref{eq:hodge-polarization}，它仍是极化。因此，在 $V$ 上选定 Kähler 结构以后，就得到极化实 Hodge 结构变分 $(\mathcal H^V,\nabla^{\mathrm{GM}},\mathcal H^V_{\mathbb R},\mathcal F^\bullet)$。

镜像一侧要考虑 $\bigoplus_{p=0}^dH^{p,p}(V^\circ)$ 上的 A 模型 Hodge 结构变分。虽然已有式~\eqref{eq:middle-a-model-polarization} 的良好极化型 $S$，但关键是：为定义 $\nabla^{\mathrm{middle}}$，须知道 $H^{p,p}(V^\circ)$ 如何嵌入 $H^{2p}(V^\circ,\mathbb C)$，因为量子上同调定义在后者上。这意味着必须固定 $V^\circ$ 的复结构。此外，$\bigoplus_{p=0}^dH^{p,p}(V^\circ)$ 一般只在 $\mathbb R$ 上定义。因此，选定 $V^\circ$ 的复结构以后，就得到极化实 Hodge 结构变分 $(\mathcal H^{V^\circ},\nabla^{\mathrm{middle}},\mathcal H^{V^\circ}_{\mathbb R},\mathcal F^\bullet)$。

还有 $\mathcal H^V$ 的特定截面 $\widetilde\Omega,g_0$，以及 $\mathcal H^{V^\circ}$ 的特定截面 $1=T_0,[pt]$。据此如下定义高维镜像对。
\begin{definition}\label{def:higher-mathematical-mirror-pair}
给定维数 $d>3$ 的 Calabi--Yau 流形 $V$、$V^\circ$，若能在 $V$ 上找到 Kähler 结构、在 $V^\circ$ 上找到复结构，使镜像映射提升为 $p_0\in\overline{\mathcal M}(V)$ 与 $q_0\in\overline{\mathcal{KM}}(V^\circ)$ 的邻域上丛 $\mathcal H^V$ 与 $\mathcal H^{V^\circ}$ 的同构，就称 $(V,V^\circ)$ 为\emph{数学镜像对}。此外，此同构保持来自 $\nabla^{\mathrm{GM}}$ 与 $\nabla^{\mathrm{middle}}$ 的极化实 Hodge 结构变分，并将 $\mathcal H^V$ 的截面 $\widetilde\Omega,g_0$ 分别映到 $\mathcal H^{V^\circ}$ 的截面 $1,[pt]$。
\end{definition}
与三维情形一样，数学镜像对有时也称为基于 Hodge 理论的镜像对。此定义有趣之处在于所需的额外数据：$V$ 上的 Kähler 结构与 $V^\circ$ 上的复结构。将这些数据加到 $\mathcal M(V)$ 与 $\mathcal{KM}(V^\circ)$ 上，意味着使用的是完整 SCFT 模空间上的镜像映射，如\hyperref[chap:introduction]{第一章}的式~\eqref{eq:mirror-foliations} 所示。还要注意，与 $p_0,q_0$ 的十分具体的选择相比，此定义对额外数据如何选择说得比较含糊。因此，或许最好将定义~\ref{def:higher-mathematical-mirror-pair} 看作初步版本。

另一个不确定之处是，A 模型变分在\IndexLocation{indexloc:0696:287}大半径极限点 $q_0$ 总为 Hodge--Tate，而\IndexLocation{indexloc:0772:287}最大幺幂单值性的定义，并不意味着 $H^d(V,\mathbb C)$ 的 Hodge 结构变分在对应点 $p_0$ 处为 Hodge--Tate。因此，定义~\ref{def:higher-mathematical-mirror-pair} 可能还需要加强。
\SourcePage{288}
尽管有这些限制，此定义与 $d$ 点函数仍有良好关系。$d$ 点函数定义如下。设镜像映射由 $q_1,\ldots,q_r$ 给出，照例同时用这些 $q_j$ 作为复模空间与 Kähler 模空间的局部坐标。令 $q_j=\exp(2\pi iu_j)$；为简化记号，令 $\delta_j=\partial/\partial u_j$。在这一约定下，$V$ 的 B 模型关联函数为
\[
Y_{i_1\cdots i_d}=\frac{(-1)^{d(d-1)/2}}{(2\pi i)^d}\int_V\widetilde\Omega\wedge\nabla_{\delta_{i_1}}\cdots\nabla_{\delta_{i_d}}\widetilde\Omega,
\qquad1\le i_1,\ldots,i_d\le r.
\]
类似地，$V^\circ$ 的 A 模型关联函数为
\[
\langle T_{i_1},\ldots,T_{i_d}\rangle=\int_{V^\circ}T_{i_1}*_{\mathrm{small}}\cdots*_{\mathrm{small}}T_{i_d}.
\]
这些关联函数具有如下关系。
\begin{proposition}\label{prop:higher-mirror-correlations}
若 $(V,V^\circ)$ 是定义~\ref{def:higher-mathematical-mirror-pair} 意义下的数学镜像对，则
\[
Y_{i_1\cdots i_d}=\langle T_{i_1},\ldots,T_{i_d}\rangle
\qquad\text{对所有 }1\le i_1,\ldots,i_d\le r.
\]
\end{proposition}
\begin{proof}
与定理~\ref{thm:mirror-correlation-equivalence} 的证明一样，定义~\ref{def:higher-mathematical-mirror-pair} 给出
\begin{equation}\label{eq:higher-mirror-pairing-equality}
Q(\widetilde\Omega,\nabla_{\delta_{i_1}}\cdots\nabla_{\delta_{i_d}}\widetilde\Omega)
=S(1,\nabla_{\delta_{i_1}}\cdots\nabla_{\delta_{i_d}}1),
\end{equation}
其中左端的 $\nabla$ 表示 Gauss--Manin 联络，右端表示 A 模型联络。

为理解左端，先注意 $\widetilde\Omega$ 是 $(d,0)$ 形式，故本原；类似地，$\nabla_{\delta_{i_1}}\cdots\nabla_{\delta_{i_d}}\widetilde\Omega$ 的 $(0,d)$ 分量也本原。由 $Q$ 的定义，不难得到式~\eqref{eq:higher-mirror-pairing-equality} 左端为 B 模型关联函数 $Y_{i_1\cdots i_d}$。

再看右端，因为所有 $T_{i_j}$ 的次数均为 $2$，
\[
\nabla_{\delta_{i_1}}\cdots\nabla_{\delta_{i_d}}1
=T_{i_1}*_{\mathrm{small}}\cdots*_{\mathrm{small}}T_{i_d}
+\text{次数小于 }2d\text{ 的项}.
\]
又因 $1\in H^0(V^\circ,\mathbb C)$，由 $S$ 的定义，右端为
\[
\begin{aligned}
\int_{V^\circ}1\cup\nabla_{\delta_{i_1}}\cdots\nabla_{\delta_{i_d}}1
&=\int_{V^\circ}\nabla_{\delta_{i_1}}\cdots\nabla_{\delta_{i_d}}1\\
&=\int_{V^\circ}T_{i_1}*_{\mathrm{small}}\cdots*_{\mathrm{small}}T_{i_d},
\end{aligned}
\]
恰为 A 模型关联函数 $\langle T_{i_1},\ldots,T_{i_d}\rangle$。
\end{proof}
虽然 $d>3$ 时关联函数有良好对应，但保持这些关联函数，不再等价于保持 Hodge 结构变分。这与三维情形很不同。通过一个例题说明其精确含义。
\SourceExampleNumber{8.6.3}{1}
\begin{example}\label{ex:fourfold-mirror-correlations}
设 $V$、$V^\circ$ 为 Calabi--Yau 四维簇。在 $V^\circ$ 上，只考虑 $\bigoplus_{p=0}^4H^{p,p}(V^\circ)$ 中的类。选取上同调基 $T_i\in H^{1,1}(V^\circ)$（$1\le i\le r$）、$T_k\in H^{2,2}(V^\circ)$（$r+1\le k\le s$）。令 $T^l\in H^{2,2}(V^\circ)$ 为对偶基，使 $\langle T_k,T^l\rangle=\delta_{kl}$。\IndexLocation{indexloc:0498:288}次数公理、\IndexLocation{indexloc:0502:288}基本类公理及\IndexLocation{indexloc:0499:288}除子公理表明，对于 $\beta\ne0$，只需考虑的非零 Gromov--Witten 不变量为 $\langle I_{0,3,\beta}\rangle(T_i,T_j,T_k)$，其中 $1\le i,j\le r$、$r+1\le k\le s$。定义
\[
Y^A_{ijk}=\sum_\beta\langle I_{0,3,\beta}\rangle(T_i,T_j,T_k)q^\beta,
\]
\SourcePage{289}
其中上标 $A$ 表示“A 模型”。于是，全部计数信息都包含在
\[
\nabla_{\delta_i}(T_j)=T_i*_{\mathrm{small}}T_j=\sum_{k=r+1}^sY^A_{ijk}T^k
\]
中，$\nabla$ 为 A 模型联络。用上述四点函数 $\langle T_i,T_j,T_a,T_b\rangle$ 表示，注意到
\begin{equation}\label{eq:fourfold-four-point-factorization}
\begin{aligned}
\langle T_i,T_j,T_a,T_b\rangle
&=\int_{V^\circ}T_i*_{\mathrm{small}}T_j*_{\mathrm{small}}T_a*_{\mathrm{small}}T_b\\
&=\int_{V^\circ}(T_i*_{\mathrm{small}}T_j)\cup(T_a*_{\mathrm{small}}T_b)\\
&=\sum_{k,l=r+1}^sY^A_{ijk}Y^A_{abl}\langle T^k,T^l\rangle,
\end{aligned}
\end{equation}
其中中间的等号来自命题~\ref{prop:quantum-frobenius-pairing}。

现在转向四维簇 $V$ 上的 B 模型。假设 $V$ 在所选的最大幺幂边界点处满足整性猜想，并且极限混合 Hodge 结构为 Hodge--Tate。与前面一样，可找到坐标 $q_1,\ldots,q_r$，使取 $e_0=\widetilde\Omega$ 后，$e_i=\nabla_{\delta_i}(e_0)$（$1\le i\le r$）给出 $\mathcal W_6\cap\mathcal F^3$ 的一组基。再令 $e^k$ 为 $\mathcal W_4\cap\mathcal F^2$ 的某组基，就可写出
\[
\nabla_{\delta_i}(e_j)=\sum_{k=r+1}^sY^B_{ijk}e^k
\]
，其中 $Y^B_{ijk}$ 为一些函数，$B$ 当然表示“B 模型”。

设 $(V,V^\circ)$ 为数学镜像对。由\hyperref[subsec:mirror-map-derivative]{第~\ref*{subsec:mirror-map-derivative} 小节}，镜像映射的导数为 $e_i\mapsto T_i$（$1\le i\le r$）。应用命题~\ref{prop:higher-mirror-correlations} 得到四点函数相等，从而可由 B 模型计算 A 模型关联函数 $\langle T_i,T_j,T_a,T_b\rangle$。但如式~\eqref{eq:fourfold-four-point-factorization} 所示，这不能再直接给出 $V^\circ$ 的 Gromov--Witten 不变量；这些不变量编码在 $Y^A_{ijk}$ 中。

如 \SourceCite{Morrison7} 指出的，真正希望得到的是等式 $Y^A_{ijk}=Y^B_{ijk}$，它能告诉我们如何用 B 模型找到 Gromov--Witten 不变量。但这需要完整同构 $\mathcal H^V\simeq\mathcal H^{V^\circ}$，特别要用到 $e^k\mapsto T^k$（$r+1\le k\le s$）。后者等价于知道同构 $H^{2,2}(V)\simeq H^{2,2}(V^\circ)$。因此，需要\hyperref[chap:introduction]{第一章}镜像对称同构~\eqref{eq:mirror-cohomology-forms} 的四维版本。

不过，若能自然选取同构 $H^{2,2}(V)\simeq H^{2,2}(V^\circ)$，应当就能给出定义~\ref{def:higher-mathematical-mirror-pair} 的精确版本，从而推出 $Y^B_{ijk}=Y^A_{ijk}$。这样就可根据 $V$ 的 Hodge 理论预测 $V^\circ$ 的 Gromov--Witten 不变量。\SourceCite{Morrison7} 讨论了这种四维镜像对称方法，并给出一些例题。此外，若同构 $H^{2,2}(V)\simeq H^{2,2}(V^\circ)$ 保持极限混合 Hodge 结构的实结构，那么定义~\ref{def:higher-mathematical-mirror-pair} 等价于对所有 $i,j,k$ 都有 $Y^B_{ijk}=Y^A_{ijk}$。因此，四维簇情形也有定理~\ref{thm:mirror-correlation-equivalence} 的类似命题，只是稍微复杂一些。
\end{example}
\SourcePage{290}
有了 $d>3$ 时镜像对的定义，就可以像 \SourceCite{Morrison7} 那样，尝试提出一般的镜像猜想。不过，由于刚性 Calabi--Yau 流形的存在，而且缺少镜像的一般描述，这样的猜想必然比较含糊。因此，这里不提出一般猜想，下面只讨论环面情形。

\subsection{环面镜像猜想}\label{subsec:toric-mirror-conjecture}
刚才指出，提出一般镜像猜想的困难之一，是给定 $V$ 后不知道如何找到 $V^\circ$。环面情形更为理想，因为可以使用\hyperref[chap:mirror-constructions]{第四章}的镜像构造。为简便起见，只考虑环面超曲面，此时 $V^\circ$ 为 $V$ 的 Batyrev 镜像。将会看到，得到的猜想在某些方面更具体，在另一些方面却不那么具体。

假设处于\hyperref[sec:batyrev-mirror]{第~\ref*{sec:batyrev-mirror} 节}的情形。固定维数至少为 $4$ 的自反多胞形 $\Delta$，令 $\Sigma$、$\Sigma^\circ$ 分别为 $\Delta$、$\Delta^\circ$ 的法扇的极大射影细分。得到极小 Calabi--Yau 环面超曲面 $V\subset X_\Sigma$ 及其 Batyrev 对偶 $V^\circ\subset X_{\Sigma^\circ}$。

为提出猜想，\hyperref[subsec:threefold-mirror-symmetry]{第~\ref*{subsec:threefold-mirror-symmetry}} 和\hyperref[subsec:higher-dimensional-mirror]{第~\ref*{subsec:higher-dimensional-mirror} 小节}提示，应从 $V$、$V^\circ$ 为三维簇的情形开始。此时得到如下\emph{环面三维簇镜像猜想}。
\begin{conjecture}\label{conj:toric-threefold-mirror}
若 $V\subset X_\Sigma$ 及其 Batyrev 对偶 $V^\circ\subset X_{\Sigma^\circ}$ 如上，且 $\dim V=\dim V^\circ=3$，则 $(V,V^\circ)$ 为定义~\ref{def:mathematical-mirror-pair} 意义下的数学镜像对。
\end{conjecture}
这是一个良好的猜想，但实际中，它与通常所谓环面镜像猜想不完全相同。原因在于，环面情形使用不同的模空间及镜像映射定义。

回顾这一过程，现在设 $\dim V=d\ge3$ 任意。\hyperref[chap:moduli]{第六章}在\hyperref[subsec:toric-complex-moduli]{第~\ref*{subsec:toric-complex-moduli} 小节}构造了简化模空间 $\overline{\mathcal M}_{\mathrm{simp}}(V)$，并在\hyperref[subsec:enlarged-kahler-moduli]{第~\ref*{subsec:enlarged-kahler-moduli} 小节}构造了环面 Kähler 模空间 $\overline{\mathcal{KM}}_{\mathrm{toric}}(V^\circ)$。为得到边界点，取由 $H^2_{\mathrm{toric}}(V^\circ,\mathbb Z)$ 的基 $T_1,\ldots,T_{\widetilde r}$ 生成的锥 $\sigma$，这些基类位于 Kähler 锥的闭包内。\hyperref[chap:moduli]{第六章}说明，$\sigma$ 给出最大幺幂边界点 $p_0\in\overline{\mathcal M}_{\mathrm{simp}}(V)$（假设猜想~\ref{conj:toric-mum-boundary} 成立），以及大半径极限点 $q_0\in\overline{\mathcal{KM}}_{\mathrm{toric}}(V^\circ)$。锥 $\sigma$ 还分别确定 $p_0$、$q_0$ 附近的显式局部坐标 $z_1,\ldots,z_{\widetilde r}$、$q_1,\ldots,q_{\widetilde r}$。此外，如\hyperref[subsec:toric-mirror-map]{第~\ref*{subsec:toric-mirror-map} 小节}所述，可在 $p_0,q_0$ 附近完全显式地确定环面版本的镜像映射。

为提出猜想，再回到 $V$、$V^\circ$ 为三维簇的情形。由命题~\ref{prop:toric-cy-construction} 证明后的说明，它们光滑。为进一步简化，先假设
\begin{equation}\label{eq:toric-divisor-surjectivity}
H^{1,1}(X_{\Sigma^\circ})\longrightarrow H^{1,1}(V^\circ)
\end{equation}
满射。按\hyperref[subsec:batyrev-hodge]{第~\ref*{subsec:batyrev-hodge} 小节}的记号，这意味着 $H^{1,1}_{\mathrm{toric}}(V^\circ)=H^{1,1}(V^\circ)$；由定理~\ref{thm:batyrev-hodge} 的证明，它又等价于 $H^{2,1}_{\mathrm{poly}}(V)=H^{2,1}(V)$。因此，上一段中的模空间正是\hyperref[subsec:threefold-mirror-symmetry]{第~\ref*{subsec:threefold-mirror-symmetry} 小节}所考虑的空间 $\overline{\mathcal{KM}}(V^\circ)$ 与 $\overline{\mathcal M}(V)$。

仍有联络 $\nabla^{\mathrm{GM}}$、$\nabla^{\mathrm{middle}}$ 及相应的 Hodge 结构变分。此外，如\hyperref[subsec:threefold-mirror-symmetry]{第~\ref*{subsec:threefold-mirror-symmetry} 小节}所述，镜像映射的导数、截面 $\widetilde\Omega$ 和 $1$、以及极化，给出唯一同构 $\mathcal H^V\simeq\mathcal H^{V^\circ}$。但即使这个同构保持联络 $\nabla^{\mathrm{GM}}$
\SourcePage{291}
和 $\nabla^{\mathrm{middle}}$，也不能断言它保持整数结构 $\mathcal H^V_{\mathbb Z}$ 和 $\mathcal H^{V^\circ}_{\mathbb Z}$。问题如下：定义~\ref{def:mathematical-mirror-pair} 中的镜像映射要求整性猜想，因而几乎由定义就保持整数结构。而\hyperref[subsec:toric-mirror-map]{第~\ref*{subsec:toric-mirror-map} 小节}给出的环面镜像映射可能与它不同，特别可能不保持整数结构。换句话说，尚不知道\hyperref[subsec:batyrev-hodge]{第~\ref*{subsec:batyrev-hodge}} 和\hyperref[subsec:toric-mirror-map]{第~\ref*{subsec:toric-mirror-map} 小节}的单项式—除子镜像映射是否保持整数结构。实际上，据我们所知，它甚至可能不保持实结构。

当然，猜想认为两种镜像映射相同，且 $V$ 满足整性猜想。但为了提出不依赖这些假设的猜想，只考虑仅涉及复数的结构。因此，对于 $V$，有 $(\mathcal H^V,\nabla^{\mathrm{GM}},\mathcal F^\bullet)$ 给出的\emph{复 Hodge 结构变分}。它由型 $Q$“极化”，这意味着第一条 Hodge--Riemann \IndexLocation{indexloc:0610:291}双线性关系的加强形式，即对所有 $p$，$(\mathcal F^p)^\perp=\mathcal F^{d-p+1}$。类似地，对于 $V^\circ$，有 $(\mathcal H^{V^\circ},\nabla^{\mathrm{middle}},\mathcal F^\bullet)$，它如\hyperref[subsec:threefold-mirror-symmetry]{第~\ref*{subsec:threefold-mirror-symmetry} 小节}所述被 $S$ 极化。

还要注意，由命题~\ref{prop:a-model-weight-filtration}，$\mathcal H^{V^\circ}$ 上的\IndexLocation{indexloc:0907:291}单值性权滤过在 $0\in\Delta^r$ 处满足 $\operatorname{Gr}^{\mathcal W}_4=H^2(V^\circ,\mathbb C)=H^{1,1}(V^\circ)$。类似地，在最大幺幂边界点处，$\mathcal H^V$ 上有 $\operatorname{Gr}^{\mathcal W}_4=H^{2,1}(V)$。这是混合 Hodge 结构为 Hodge--Tate 的结果；如式~\eqref{eq:mum-hodge-tate-splitting} 附近所述，在三维情形，它来自最大幺幂单值性。

在这种情形，得到\emph{基于 Hodge 理论的\IndexLocation{indexloc:0791:291}\IndexLocation{indexloc:0795:291}\IndexLocation{indexloc:0797:291}环面镜像猜想}的下列特殊情形。
\begin{conjecture}\label{conj:toric-hodge-mirror-surjective}
设 $V$、$V^\circ$ 如上为互为 Batyrev 镜像的三维 Calabi--Yau 环面超曲面，并假设式~\eqref{eq:toric-divisor-surjectivity} 满射。则镜像映射提升为 $p_0\in\overline{\mathcal M}(V)$ 和 $q_0\in\overline{\mathcal{KM}}(V^\circ)$ 的邻域上，丛 $\mathcal H^V$ 与 $\mathcal H^{V^\circ}$ 的同构。此外，此同构保持来自 $\nabla^{\mathrm{GM}}$ 与 $\nabla^{\mathrm{middle}}$ 的极化复 Hodge 结构变分，并将 $\mathcal H^V$ 的截面 $\widetilde\Omega,g_0$ 分别映到 $\mathcal H^{V^\circ}$ 的截面 $1,[pt]$。最后，在 $p_0,q_0$ 处的 $\operatorname{Gr}^{\mathcal W}_4$ 上，诱导同构就是单项式—除子镜像映射。
\end{conjecture}
\index{镜像猜想!环面情形|hyperpage}
与定理~\ref{thm:mirror-correlation-equivalence} 类似，此时又可用 A 模型、B 模型关联函数表述镜像对称。与前面一样，$V^\circ$ 上有通常的 A 模型关联函数 $\langle T_i,T_j,T_k\rangle$。在 $V$ 上，定义\emph{环面 B 模型关联函数}为
\[
Y_{ijk}=\int_V\widetilde\Omega\wedge\nabla_{\delta_i}\nabla_{\delta_j}\nabla_{\delta_k}\widetilde\Omega,
\]
其中 $\delta_j=2\pi iq_j\partial/\partial q_j$，$\widetilde\Omega=\Omega/y_0$ 为归一化三形式。它与\hyperref[subsec:normalized-yukawa-mirror]{第~\ref*{subsec:normalized-yukawa-mirror} 小节}定义的归一化 Yukawa 耦合有两点不同：
\begin{itemize}
\item 这里的 $y_0$ 是 $\Omega$ 的 Picard--Fuchs 方程在最大幺幂边界点处的任意非零全纯解。因此，只在相差常数的意义下知道 $y_0$，从而也只在此意义下知道 $\widetilde\Omega$。
\item 使用的是环面版本的镜像映射，它可能与定义归一化 Yukawa 耦合时所用的镜像映射不同。
\end{itemize}
因此，$Y_{ijk}$ 只在相差乘法常数的意义下确定；与定义~\ref{def:normalized-yukawa} 相比，上式没有负号，原因就在这里。不过，如\hyperref[subsec:normalized-yukawa-mirror]{第~\ref*{subsec:normalized-yukawa-mirror} 小节}指出的，除非极其详细地研究单值性，我们的计算方法也只能
\SourcePage{292}
在相差常数的意义下确定 $Y_{ijk}$。幸而，此时仍有势函数（引理~\ref{lem:yukawa-potential} 的证明仍适用），故本质上只须处理一个常数。

由此容易给出定理~\ref{thm:mirror-correlation-equivalence} 的一个版本，断言猜想~\ref{conj:toric-hodge-mirror-surjective} 等价于
\begin{equation}\label{eq:toric-mirror-correlations}
Y_{ijk}=\langle T_i,T_j,T_k\rangle\qquad\text{对所有 }i,j,k,
\end{equation}
其中允许将所有 $Y_{ijk}$ 同乘一个与 $i,j,k$ 无关的常数 $c$，使等式成立。这个证明直接，细节略去。

以下例题说明，所涉及模空间为二维时，猜想~\ref{conj:toric-hodge-mirror-surjective} 具有怎样的形式。
\SourceExampleNumber{8.6.4}{1}
\begin{example}\label{ex:octic-toric-mirror-correlations}
一个多次考虑过的例子，是 $\mathbb P(1,1,2,2,2)$ 中八次超曲面的环面解消 $V$ 及其 Batyrev 镜像 $V^\circ$。已知 $V$ 的环面 Kähler 模空间为二维（这里也等于全部 Kähler 模空间的维数），故式~\eqref{eq:toric-divisor-surjectivity} 满射。先写出 $V$ 的一个 A 模型关联函数。

在例题~\ref{ex:weighted-extremal-transition} 中看到，$V$ 的 Kähler 锥由记为 $D_3,D_1$ 的类生成。由 \SourceCite{CdFKM}，这些类可描述如下：$D_3=H$ 是 $\mathbb P(1,1,2,2,2)$ 上由二次方程定义的丰沛类的严格变换；$D_1=L$ 是由一次方程定义的除子的严格变换。若 $E\subset V$ 为映射 $V\to\overline V$ 的例外轨迹，则 $H\sim2L+E$。如该例指出的，$H^3=8$、$L^3=0$。

由 $H,L$ 生成的锥 $\sigma$（这里是整个 Kähler 锥）确定 $V$ 的大半径极限点和 $V^\circ$ 的最大幺幂边界点。$V$ 的 Kähler 模空间的局部坐标为 $q_j=\exp(2\pi iu_j)$，其中 $\sigma=\{u_1H+u_2L:u_1,u_2\ge0\}$。若 $h,\ell$ 为 $H_2(V,\mathbb Z)$ 中的对偶基，则 $\beta\in H_2(V,\mathbb Z)$ 可写为 $\beta=ah+b\ell$。因为 $H^3=8$、$\int_\beta H=a$，式~\eqref{eq:cy-quantum-correlation-instanton} 给出
\begin{equation}\label{eq:octic-a-model-instants}
\langle H,H,H\rangle=8+\sum_{(a,b)\ne(0,0)}n(a,b)\frac{a^3q_1^aq_2^b}{1-q_1^aq_2^b},
\end{equation}
其中 $n(a,b)$ 为\hyperref[subsec:instanton-numbers]{第~\ref*{subsec:instanton-numbers} 小节}定义的瞬子数 $n_{ah+b\ell}$。粗略地说，$n(a,b)$ 是 $V$ 中同调类为 $ah+b\ell$ 的有理曲线数；但例题~\ref{ex:quintic-instanton-enumerative} 表明，瞬子数与有理曲线的关系比初看时更微妙。

利用猜想~\ref{conj:toric-hodge-mirror-surjective} 求出 $\langle H,H,H\rangle$ 的猜想公式，就能同时预测所有瞬子数 $n(a,b)$。先注意，按这里的记号，使用环面镜像猜想时交换了 $V$ 与 $V^\circ$ 的角色。希望这不会造成太多混淆。以下用 $\nabla$ 表示 $V^\circ$ 的 Gauss--Manin 联络。由式~\eqref{eq:toric-mirror-correlations}，环面镜像猜想预测
\[
\langle H,H,H\rangle=Y_{111}=\int_{V^\circ}\widetilde\Omega\wedge\nabla_{\delta_1}\nabla_{\delta_1}\nabla_{\delta_1}\widetilde\Omega,
\]
其中 $\delta_1=2\pi iq_1\partial/\partial q_1$。

为计算右端积分，先用 $V^\circ$ 的模坐标 $z_1,z_2$。在例题~\ref{ex:weighted-yukawa} 中，已利用这些变量计算 Yukawa
\SourcePage{293}
耦合
\[
K^{ij}=\int_{V^\circ}\Omega\wedge\nabla_{z_1\frac{\partial}{\partial z_1}}^i\nabla_{z_2\frac{\partial}{\partial z_2}}^j\Omega,
\qquad i+j=3.
\]
现在须将 $\Omega$ 除以 $y_0$ 归一化，再转到镜像坐标 $q_1,q_2$。因为
\[
\delta_1=2\pi iq_1\frac{\partial}{\partial q_1}
=2\pi i\left(\frac{q_1}{z_1}\frac{\partial z_1}{\partial q_1}z_1\frac{\partial}{\partial z_1}
+\frac{q_1}{z_2}\frac{\partial z_2}{\partial q_1}z_2\frac{\partial}{\partial z_2}\right),
\]
代入 $\langle H,H,H\rangle$ 的公式，得到
\[
\begin{aligned}
\langle H,H,H\rangle=\frac{(2\pi i)^3}{y_0^2}\Biggl[&
\left(\frac{q_1}{z_1}\frac{\partial z_1}{\partial q_1}\right)^3K^{30}
+3\left(\frac{q_1}{z_1}\frac{\partial z_1}{\partial q_1}\right)^2
\left(\frac{q_1}{z_2}\frac{\partial z_2}{\partial q_1}\right)K^{21}\\
&+3\left(\frac{q_1}{z_1}\frac{\partial z_1}{\partial q_1}\right)
\left(\frac{q_1}{z_2}\frac{\partial z_2}{\partial q_1}\right)^2K^{12}
+\left(\frac{q_1}{z_2}\frac{\partial z_2}{\partial q_1}\right)^3K^{03}\Biggr].
\end{aligned}
\]
例题~\ref{ex:weighted-yukawa} 已显式计算 $K^{ij}$，只差一个共同常数 $c\ne0$；例题~\ref{ex:weighted-hypergeometric-mirror-map} 又将 $y_0,q_2$ 显式计算为 $z_1,z_2$ 的函数。$q_1$ 也可这样计算，因此上述公式能将 $\langle H,H,H\rangle$ 表示为 $q_1,q_2$ 的函数。

先确定常数 $c$。由镜像映射的环面归一化~\eqref{eq:toric-mirror-normalization}，$q_j=z_j+\cdots$；由例题~\ref{ex:weighted-yukawa}，
\[
K^{30}=\frac c{(1-256z_1)^2-512^2z_1^2z_2},
\]
故 $K^{30}=c+\cdots$。又因 $y_0=1+\cdots$，上述公式给出
\[
\langle H,H,H\rangle=(2\pi i)^3c+\cdots.
\]
而 $\langle H,H,H\rangle=8+\cdots$，故必须有
\[
c=\frac8{(2\pi i)^3},
\]
这与\hyperref[chap:quintic]{第二章}出现的情形相似。取此 $c$，得到展开
\[
\begin{aligned}
\langle H,H,H\rangle={}&8+640\frac{q_1}{1-q_1}
+10032\frac{2^3q_1^2}{1-q_1^2}+640\frac{q_1q_2}{1-q_1q_2}\\
&+288384\frac{3^3q_1^3}{1-q_1^3}
+72224\frac{2^3q_1^2q_2}{1-q_1^2q_2}
+10979984\frac{4^3q_1^4}{1-q_1^4}+\cdots.
\end{aligned}
\]
只要有耐心及足够计算能力，便可计算任意多项。与式~\eqref{eq:octic-a-model-instants} 比较，得到瞬子数 $n(a,b)$。可证明，$a\ge1$ 时有 $n(a,b)=n(a,a-b)$，例如 $n(1,1)=n(1,0)=640$；$b>a$ 时，除 $n(0,1)=4$ 外，都有 $n(a,b)=0$。因此，$\langle H,H,H\rangle$ 的级数确定所有 $(a,b)\ne(0,1)$ 对应的 $n(a,b)$。此例及其他 Calabi--Yau 环面超曲面例子的进一步计算，见 \SourceCite{CdFKM,HKTY1}。

\hyperref[chap:mirror-theorem]{第十一章}的例题~\ref{ex:octic-toric-mirror-theorem} 将说明，上述公式来自定理~\ref{thm:givental-toric-mirror}。此外，利用 \SourceCite{CdFKM} 中详细的单值性计算，可以实际证明，此例\IndexLocation{indexloc:1031:293}满足\IndexLocation{indexloc:0636:293}整性猜想，且取 $c=8/(2\pi i)^3$ 后，$Y_{111}$ 就是\hyperref[subsec:normalized-yukawa-mirror]{第~\ref*{subsec:normalized-yukawa-mirror} 小节}严格意义下的归一化 Yukawa 耦合。
\end{example}
\SourcePage{294}
下面考虑 $V$ 及其 Batyrev 镜像 $V^\circ$ 仍为三维簇，但不再假设映射~\eqref{eq:toric-divisor-surjectivity} 满射的情形。这时须处理\hyperref[subsec:toric-complex-moduli]{第~\ref*{subsec:toric-complex-moduli} 小节}的简化模空间 $\overline{\mathcal M}_{\mathrm{simp}}(V)$ 和\hyperref[subsec:enlarged-kahler-moduli]{第~\ref*{subsec:enlarged-kahler-moduli} 小节}的环面 Kähler 模空间 $\overline{\mathcal{KM}}_{\mathrm{toric}}(V^\circ)$。

一种提出猜想的方式如下。完整模空间上有丛 $\mathcal H^V$、$\mathcal H^{V^\circ}$；限制后得到
\[
\begin{aligned}
\mathcal H^V_{\mathrm{simp}}&=\mathcal H^V\text{ 在 }\overline{\mathcal M}_{\mathrm{simp}}(V)\text{ 上的限制},\\
\mathcal H^{V^\circ}_{\mathrm{toric}}&=\mathcal H^{V^\circ}\text{ 在 }\overline{\mathcal{KM}}_{\mathrm{toric}}(V^\circ)\text{ 上的限制}.
\end{aligned}
\]
这些限制丛继承自然的极化 Hodge 结构变分。于是，基于 Hodge 理论的环面镜像猜想会断言，存在保持极化复 Hodge 结构变分的丛同构 $\mathcal H^V_{\mathrm{simp}}\simeq\mathcal H^{V^\circ}_{\mathrm{toric}}$，如猜想~\ref{conj:toric-hodge-mirror-surjective} 所述。不过，不将它正式列为猜想，因为这样的丛同构需要知道同构 $H^{2,1}(V)\simeq H^{1,1}(V^\circ)$，而环面情形只有
\[
H^{2,1}_{\mathrm{poly}}(V)\simeq H^{1,1}_{\mathrm{toric}}(V^\circ),
\]
它来自\hyperref[subsec:batyrev-hodge]{第~\ref*{subsec:batyrev-hodge}} 和\hyperref[subsec:toric-mirror-map]{第~\ref*{subsec:toric-mirror-map} 小节}的单项式—除子镜像映射。回顾，$H^{2,1}_{\mathrm{poly}}(V)$ 是 $H^{2,1}(V)\simeq H^1(V,\Theta_V)$ 的子空间，对应于改变 $V\subset X_\Sigma$ 的定义方程得到的形变。因此，问题在于还不理解如何定义 $\mathcal H^V_{\mathrm{simp}}\simeq\mathcal H^{V^\circ}_{\mathrm{toric}}$。

不过，若将这两个丛替换为略小的丛，就可给出相当精确的环面镜像猜想版本。先从 $V^\circ$ 开始。令 $i:V^\circ\hookrightarrow X_{\Sigma^\circ}$ 为包含映射，定义
\[
H^*_{\mathrm{toric}}(V^\circ)=\operatorname{im}\bigl(i^*:H^*(X_{\Sigma^\circ})\longrightarrow H^*(V^\circ)\bigr).
\]
因为 $p\ne q$ 时 $H^{p,q}(X_{\Sigma^\circ})=0$，$H^*_{\mathrm{toric}}(V^\circ)$ 由 $(p,p)$ 型上同调类组成。因此，限制到该空间得到子丛
\[
\begin{aligned}
H^*_{\mathrm{toric}}(V^\circ)\times\overline{\mathcal{KM}}_{\mathrm{toric}}(V^\circ)
&\subset H^{\mathrm{middle}}(V^\circ)\times\overline{\mathcal{KM}}_{\mathrm{toric}}(V^\circ)\\
&=\mathcal H^{V^\circ}_{\mathrm{toric}}.
\end{aligned}
\]
将此子丛记为 $\mathcal T\mathcal H^{V^\circ}_{\mathrm{toric}}$，其中 $\mathcal T$ 表示所用的环面子环。

当 $X_{\Sigma^\circ}$ 凸（例如 $X_{\Sigma^\circ}=\mathbb P^4$）且 $V^\circ$ 丰沛时，\SourceCite[命题 4]{Pandharipande3} 的论证表明，$H^*_{\mathrm{toric}}(V^\circ)$ 是 $H^*(V^\circ)$ 关于小量子乘积的子环。一般情形是否如此尚不知道，因此进一步假设
\begin{equation}\label{eq:toric-quantum-subring-assumption}
H^*_{\mathrm{toric}}(V^\circ)\text{ 是 }H^*(V^\circ)\text{ 关于 }*_{\mathrm{small}}\text{ 的子环}.
\end{equation}
在此假设下，A 模型联络 $\nabla^{\mathrm{middle}}$ 在 $\mathcal T\mathcal H^{V^\circ}_{\mathrm{toric}}$ 上诱导联络。断言实际上得到极化 Hodge 结构变分。为证明这一点，假设 $X_{\Sigma^\circ}$ 光滑，则有如下引理。
\begin{lemma}\label{lem:toric-cohomology-pairing-nondegenerate}
若 $V^\circ$ 是光滑环面簇 $X_{\Sigma^\circ}$ 中的 Calabi--Yau 环面超曲面，则杯积配对在 $H^*_{\mathrm{toric}}(V^\circ)$ 上的限制非退化。
\end{lemma}
\SourcePage{295}
\begin{proof}
只须证明
\[
\int_{V^\circ}i^*(\alpha)\cup i^*(\beta)=0\quad\text{对所有 }\beta\in H^*(X_{\Sigma^\circ})
\quad\Longrightarrow\quad i^*(\alpha)=0.
\]
由投影公式，
\[
\int_{V^\circ}i^*(\alpha)\cup i^*(\beta)=\int_{X_{\Sigma^\circ}}\alpha\cup\beta\cup[V^\circ].
\]
而 \SourceCite[(3.36)]{HLY1} 意味着，对于 $\alpha\in H^*(X_{\Sigma^\circ})$，有
\[
\alpha\cup[V^\circ]=0\text{ 于 }H^*(X_{\Sigma^\circ})
\quad\Longleftrightarrow\quad
 i^*(\alpha)=0\text{ 于 }H^*_{\mathrm{toric}}(V^\circ).
\]
这两个事实容易推出所需的非退化性。
\end{proof}
假设式~\eqref{eq:toric-quantum-subring-assumption} 成立，便可将 $\mathcal H^{V^\circ}_{\mathrm{toric}}$ 替换为 $\mathcal T\mathcal H^{V^\circ}_{\mathrm{toric}}$，重复\hyperref[sec:a-model-vhs]{第~\ref*{sec:a-model-vhs} 节}的构造。极化来自刚证明的引理；如\hyperref[subsec:a-model-asymptotics]{第~\ref*{subsec:a-model-asymptotics} 小节}所述，还得到 $\mathcal T\mathcal H^{V^\circ}_{\mathrm{toric}}$ 的整数结构，因为 $H^*_{\mathrm{toric}}(V^\circ)$ 自然在整数上定义。

镜像一侧需要找到对应的 Hodge 结构变分。回顾定理~\ref{thm:batyrev-hodge} 的证明，其中有正合列
\[
\longrightarrow\bigoplus_iH^1(V\cap D_i,\mathbb C)\longrightarrow H^3(V,\mathbb C)
\longrightarrow\operatorname{Gr}^W_3H^3(Z_f,\mathbb C)\longrightarrow0.
\]
这里 $D_i$ 是 $X_\Sigma$ 的环面除子，$Z_f=T\cap V$ 是环面 $T\subset X_\Sigma$ 中的仿射超曲面。在该证明中，最终得到
\[
h^{2,1}\bigl(\operatorname{Gr}^W_3H^3(Z_f,\mathbb C)\bigr)=h^{2,1}_{\mathrm{poly}}(V).
\]
这提示存在自然同构
\begin{equation}\label{eq:polynomial-affine-hodge-assumption}
H^{2,1}_{\mathrm{poly}}(V)\simeq H^{2,1}\bigl(\operatorname{Gr}^W_3H^3(Z_f,\mathbb C)\bigr),
\end{equation}
由 $H^{2,1}_{\mathrm{poly}}(V)\subset H^{2,1}(V)\to H^{2,1}(\operatorname{Gr}^W_3H^3(Z_f,\mathbb C))$ 诱导，最后一个映射来自上述正合列。因为其证明尚未写出，这里直接假设式~\eqref{eq:polynomial-affine-hodge-assumption} 是同构。

在此假设下，利用 $H^{2,1}_{\mathrm{poly}}(V)$ 构造 Hodge 结构变分的自然办法，就是使用模空间 $\overline{\mathcal M}_{\mathrm{simp}}(V)$ 上 $\operatorname{Gr}^W_3H^3(Z_f,\mathbb C)$ 的自然极化 Hodge 结构变分。将所得丛记为 $\mathcal T\mathcal H^V_{\mathrm{simp}}$；这时 $\mathcal T$ 表示取交所用的环面。现在可以表述三维簇的基于 Hodge 理论的\IndexLocation{indexloc:0791:295}\IndexLocation{indexloc:0795:295}\IndexLocation{indexloc:0797:295}环面镜像猜想。
\begin{conjecture}\label{conj:toric-hodge-mirror-general}\index{镜像猜想!Hodge 理论|hyperpage}
设 $V$、$V^\circ$ 如上为互为 Batyrev 镜像的三维 Calabi--Yau 环面超曲面，并假设式~\eqref{eq:toric-quantum-subring-assumption}、\eqref{eq:polynomial-affine-hodge-assumption} 成立，且 $X_{\Sigma^\circ}$ 光滑。则镜像映射提升为 $p_0\in\overline{\mathcal M}_{\mathrm{simp}}(V)$ 和 $q_0\in\overline{\mathcal{KM}}_{\mathrm{toric}}(V^\circ)$ 的邻域上，丛 $\mathcal T\mathcal H^V_{\mathrm{simp}}$ 与 $\mathcal T\mathcal H^{V^\circ}_{\mathrm{toric}}$ 的同构。此外，此同构保持来自 $\nabla^{\mathrm{GM}}$ 与 $\nabla^{\mathrm{middle}}$ 的极化复 Hodge 结构变分，并将 $\widetilde\Omega,g_0$ 分别映到 $1,[pt]$。最后，在 $p_0,q_0$ 处的 $\operatorname{Gr}^{\mathcal W}_4$ 上，诱导同构就是单项式—\IndexLocation{indexloc:0910:295}除子镜像映射。
\end{conjecture}
因为此猜想的表述包含许多假设，应将其看作“尚在推进中的工作”。但希望猜想~\ref{conj:toric-hodge-mirror-general} 已相当接近真实情形。一个令人鼓舞的事实是，式~\eqref{eq:toric-divisor-surjectivity} 满射时，此猜想化为猜想~\ref{conj:toric-hodge-mirror-surjective}。
\SourcePage{296}
最后考虑 $V$、$V^\circ$ 的维数大于 $3$ 的情形。尝试提出环面镜像猜想时，会立即遇到一些问题。主要问题是，$\dim V>3$ 时，极小 Calabi--Yau 簇 $V$ 可能有奇点，而我们只对光滑射影簇定义了 Gromov--Witten 不变量。这大概只是技术困难；最终希望 Gromov--Witten 理论能适用于更广泛的簇。此外，若丛同构 $\mathcal H^V_{\mathrm{simp}}\simeq\mathcal H^{V^\circ}_{\mathrm{toric}}$ 保持 Hodge 滤过，则会有
\begin{equation}\label{eq:toric-mirror-hodge-equality}
h^{p,q}(V)=h^{q,q}(V^\circ),\qquad p+q=\dim V.
\end{equation}
\hyperref[chap:mirror-constructions]{第四章}已证明 $q=1$ 的情形；对于其他 $q$，已知的最好结果是联系 $V$、$V^\circ$ 的弦理论 Hodge 数的公式~\eqref{eq:complete-intersection-string-hodge}
\[
h^{p,q}_{\mathrm{st}}(V)=h^{q,q}_{\mathrm{st}}(V^\circ),\qquad p+q=\dim V,
\]
见 \SourceCite{BD}。$V$、$V^\circ$ 光滑时，此式化为式~\eqref{eq:toric-mirror-hodge-equality}；但 $V$ 有奇点时，$h^{p,q}_{\mathrm{st}}(V)$ 可能不同于 $h^{p,q}(V)$。

当然，可假设 $V$、$V^\circ$ 光滑来避开这一困难，但这会限制猜想的适用范围。即便如此，仍需知道如何定义 $\mathcal H^V_{\mathrm{simp}}\simeq\mathcal H^{V^\circ}_{\mathrm{toric}}$。仅有 Hodge 数等式~\eqref{eq:toric-mirror-hodge-equality} 还不够，必须有实际同构 $H^{p,q}(V)\simeq H^{q,q}(V^\circ)$，其中 $p+q=\dim V$。这与\hyperref[subsec:higher-dimensional-mirror]{第~\ref*{subsec:higher-dimensional-mirror} 小节}遇到的问题相同。

因此，不在一般情形表述基于 Hodge 理论的环面镜像猜想。猜想~\ref{conj:toric-hodge-mirror-general} 的某个版本可能在高维成立，但要给出精确猜想，还需更多工作。按 \SourceCite[第 8.2 讲]{Morrison7} 中基于 Hodge 理论的镜像对称猜想的表述，最终版本应涉及某些 Hodge 结构子变分，例如 $\mathcal T\mathcal H^{V^\circ}_{\mathrm{toric}}$。不过，找到这个猜想的精确版本，仍是代数几何中的开放问题。

\subsection{结语}\label{subsec:mirror-conjecture-conclusion}
本节严格定义了数学镜像对（定义~\ref{def:mathematical-mirror-pair}、\ref{def:higher-mathematical-mirror-pair}），并考虑了基于 Hodge 理论的环面镜像猜想的不同版本（猜想~\ref{conj:toric-hodge-mirror-surjective}、\ref{conj:toric-hodge-mirror-general}）。这些表明，对于三维簇，已经较好地理解镜像对称的含义，虽然猜想~\ref{conj:toric-hodge-mirror-general} 仍需进一步研究。但进入高维以后，猜想中还有许多部分未予明确。即便在环面情形，距离基于 Hodge 理论的镜像猜想的确定版本，也还有很长的路。

不过，近期文献中已出现若干“镜像定理”，包括五次三维簇镜像对称的完整证明，以及 Calabi--Yau 环面超曲面和更一般的 Calabi--Yau 环面完全交的结果。后者使用\hyperref[sec:toric-complete-intersections]{第~\ref*{sec:toric-complete-intersections} 节}的 Batyrev--Borisov 镜像构造。\footnote{环面簇中 Calabi--Yau 完全交的讨论和例子，见 \SourceCite{BvS,LTe,HKTY2,Givental2,Givental4}，以及\hyperref[subsec:other-gw-examples]{第~\ref*{subsec:other-gw-examples} 小节}列出的参考文献。}此外，Givental 提出了适用于某些非 Calabi--Yau 环面完全交的“镜像定理”；\SourceCite{LLY} 的“镜像原理”则适用于许多超出 Calabi--Yau 三维簇及其镜像的情形。
\SourcePage{297}
\hyperref[chap:mirror-theorem]{第十一章}将看到，尽管这些“镜像定理”没有直接提及这里讨论的基于 Hodge 理论的镜像对称版本，它们仍有共同的基本思想：把计算 Gromov--Witten 不变量的困难问题转化为更简单的问题。维数大于 $3$ 时，尚不清楚这些结果与 Hodge 理论的关系。因此，正如在维数大于 $3$ 时还没有确定的“基于 Hodge 理论的镜像猜想”，也还没有该情形下确定的“镜像定理”。希望代数几何学家继续探索镜像对称的数学理论，证明更一般的“镜像定理”，提出更精确的“镜像猜想”，最终汇合为对镜像对称现象的广泛理解。

本章使我们更清楚地认识到，镜像对称的数学版本应具有怎样的形式。$\dim V>3$ 时没有精确的“镜像猜想”表述，固然令人遗憾，但这不应妨碍考察目前已经证明的“镜像定理”。带着这一目标，接下来\hyperref[chap:localization]{第九章}、\hyperref[chap:quantum-differential-equations]{第十章}将发展一些必要工具，然后在\hyperref[chap:mirror-theorem]{第十一章}研究若干“镜像定理”。
